\documentclass[11pt,a4paper]{article}

\usepackage[margin=2.1cm,top=2.3cm,bottom=2.3cm]{geometry}
\usepackage{fontspec}
\setmainfont{Cambria}
\setsansfont{Segoe UI}
\setmonofont{Consolas}[Scale=0.88]
\usepackage{microtype}
\usepackage{booktabs}
\usepackage{longtable}
\usepackage{array}
\usepackage{tabularx}
\usepackage{enumitem}
\usepackage{titlesec}
\usepackage{fancyhdr}
\usepackage{amsmath}
\usepackage{amssymb}
\usepackage[dvipsnames,table]{xcolor}
\usepackage[most]{tcolorbox}
\usepackage[hidelinks]{hyperref}
\usepackage{needspace}

\definecolor{ink}{HTML}{1A1A1A}
\definecolor{rulegrey}{HTML}{C8C8C8}
\definecolor{accent}{HTML}{7A2E2E}
\definecolor{alive}{HTML}{1F6F3D}
\definecolor{decayed}{HTML}{9A3412}
\definecolor{wastebg}{HTML}{FDF4F4}
\definecolor{wasteline}{HTML}{B45050}
\definecolor{factbg}{HTML}{F2F5F8}
\definecolor{factline}{HTML}{4A6D8C}
\definecolor{soft}{HTML}{5A5A5A}

\color{ink}
\hypersetup{colorlinks=true,linkcolor=accent,urlcolor=accent,citecolor=accent}

\setlist[itemize]{leftmargin=1.1em,itemsep=1pt,topsep=3pt,parsep=0pt}
\setlist[enumerate]{leftmargin=1.4em,itemsep=1pt,topsep=3pt,parsep=0pt}

\titleformat{\section}{\sffamily\Large\bfseries\color{accent}}{\thesection.}{0.6em}{}
\titlespacing{\section}{0pt}{1.6em}{0.7em}
\titleformat{\subsection}{\sffamily\large\bfseries}{\thesubsection}{0.6em}{}
\titlespacing{\subsection}{0pt}{1.1em}{0.4em}
\titleformat{\subsubsection}{\sffamily\normalsize\bfseries}{}{0em}{}
\titlespacing{\subsubsection}{0pt}{0.8em}{0.3em}

\pagestyle{fancy}
\fancyhf{}
\fancyhead[L]{\sffamily\footnotesize\color{soft}Trading and Quant Research Map}
\fancyhead[R]{\sffamily\footnotesize\color{soft}\thepage}
\renewcommand{\headrulewidth}{0.4pt}
\renewcommand{\headrule}{\hbox to\headwidth{\color{rulegrey}\leaders\hrule height \headrulewidth\hfill}}

\newtcolorbox{waste}{colback=wastebg,colframe=wasteline,boxrule=0.5pt,left=6pt,right=6pt,top=4pt,bottom=4pt,
  fonttitle=\sffamily\bfseries\footnotesize,title={Why this might waste your time},coltitle=white,
  colbacktitle=wasteline,arc=1pt,before skip=5pt,after skip=8pt}

\newtcolorbox{keyfact}[1]{colback=factbg,colframe=factline,boxrule=0.5pt,left=6pt,right=6pt,top=4pt,bottom=4pt,
  fonttitle=\sffamily\bfseries\footnotesize,title={#1},coltitle=white,
  colbacktitle=factline,arc=1pt,before skip=6pt,after skip=8pt}

\newcommand{\alv}[1]{\textcolor{alive}{\textbf{#1}}}
\newcommand{\dec}[1]{\textcolor{decayed}{\textbf{#1}}}
\newcommand{\fitem}[1]{\par\smallskip\noindent\textbf{#1} }
\newcommand{\rs}{₹}
\newcommand{\srcnote}[1]{{\footnotesize\color{soft}#1\par}}

\begin{document}

\begin{titlepage}
\centering
\vspace*{2.5cm}
{\sffamily\Huge\bfseries\color{accent} The Trading and Quant\par}
\vspace{2mm}
{\sffamily\Huge\bfseries\color{accent} Research Map\par}
\vspace{7mm}
{\sffamily\large A triaged survey of what is worth learning,\par}
{\sffamily\large what is tradeable, and what will waste your time\par}
\vspace{14mm}
{\large Prepared for Lalit Karthik\par}
\vspace{2mm}
{\color{soft} Convex Hedge \textbullet\ 29 August 2026\par}
\vfill
\begin{minipage}{0.84\textwidth}
\footnotesize\color{soft}
Built from fifteen independent research agents across two rounds, plus one
adversarial completeness critic. Every source was filtered mechanically: it
qualifies only if at least two of five practitioner tests hold, and any backtest
presented without a cost model was rejected outright regardless of what else it
passed. Claims a second round could not verify are marked as such rather than
smoothed over. Where the two rounds disagreed, the disagreement was adjudicated
against primary data and the adjudication is shown.

\medskip
This document reports what is documented. It contains no investment advice, no
recommended allocations and no trade suggestions. Tax and regulatory material
states rules and their arithmetic consequences; verify anything you act on with a
professional.
\end{minipage}
\vspace{1.5cm}
\end{titlepage}

\tableofcontents
\newpage

\section{How to read this}

\subsection*{The filter every source had to pass}

A source qualifies only if at least \textbf{two} of these five hold:

\begin{enumerate}
\item runnable code exists;
\item it states transaction costs, slippage or fills explicitly;
\item the author has actually traded or built it, not merely written about it;
\item \textbf{it names its own failure modes and where the technique stops working} --- this counts double;
\item it is reproducible on data you can actually obtain.
\end{enumerate}

Auto-rejected regardless of anything else: any backtest presented with no cost
model; anything whose primary funnel is a paid course, membership or signal
service; ``I made \$X'' and win-rate marketing; anonymous authors; and
broker-affiliated ``education'' where order flow is the business model.

Test~(4) carries double weight deliberately. A source that tells you where a
technique breaks is worth three that tell you it works. It is also the test that
most of the field fails.

\subsection*{What the words mean}

\begin{itemize}
\item \alv{ALIVE} --- still works in practice as of 2026, for someone in your position.
\item \dec{DECAYED} --- the effect was real and has been arbitraged, crowded or regulated out.
\item \textbf{Firm-only} --- the technique works, but the barrier is capital, exchange membership or infrastructure you cannot assemble alone.
\item \textbf{Unverified} --- reported from search metadata, or from a source that could not be opened. A lead, never evidence.
\end{itemize}

\subsection*{The one honest caveat about this document}

Two rounds produced material of genuinely different quality. The strongest
findings here come from agents that \emph{measured something} --- downloaded an
NSE bhavcopy, queried a live exchange API, cloned a repository and read its cost
code. The weakest are literature summaries. Where a finding rests on measurement,
this document says what was measured and on what date, so you can re-run it.
Where it rests on reading, it says so.

Prices, rates and market-structure facts were current on 27--28 August 2026, and
several are on a known schedule to change. Section~\ref{sec:breaks} lists the
structural breaks explicitly, because at least three of them fall inside the data
you are about to receive.

\begin{keyfact}{The single most useful habit this research suggests}
Before believing any options result --- yours or anyone's --- find the line where
the bid-ask half-spread is charged, and read the units. A published,
peer-reviewed, code-shipping 0DTE paper had that line wrong by a factor of one
hundred. When it was corrected, essentially every strategy in it went from
positive to negative net Sharpe. Section~\ref{sec:vilkov} gives the anatomy.
\end{keyfact}

\newpage
\section{The landscape map}
\label{sec:landscape}

Twenty-five fields, each with the same eight items, so you can compare them side
by side and decide what to ignore. The eighth item is mandatory and is the reason
this is a triage document rather than a menu: every survey fails by making
everything sound worthwhile. A ranked shortlist closes the section.

\subsection{Options and volatility}

\subsubsection{1. Arbitrage-free surface modelling (SVI / SSVI / local vol)}
\fitem{Is}Parametric fitting of the implied vol surface subject to static
no-arbitrage constraints, plus the Dupire local-vol bridge. Output is a clean,
interpolable, differentiable surface.
\fitem{Uniquely solves}Every other options technique needs a surface that does
not produce negative densities or fake calendar arbitrage at 09:20 when the chain
is garbage. Nothing else gives you that.
\fitem{Status}\alv{ALIVE as infrastructure}, \dec{DECAYED as alpha}. Fitting SVI
has not been a trade since roughly 2012. What survives is intraday refitting on
minute data and quote-quality filtering.
\fitem{Prereqs}You have them. Infrastructure: a tick-to-surface pipeline and a
fast constrained optimiser. Capital: zero.
\fitem{First useful result}2--3 weeks, to convert your existing research artefact
into a minute-wise NIFTY surface service with residual diagnostics.
\fitem{Solo?}Entirely.
\fitem{Sources}Gatheral \& Jacquier, \emph{Arbitrage-free SVI volatility
surfaces} (arXiv:1204.0646); the Zeliade quasi-explicit calibration white paper,
for the failure geometry rather than the formula; Bennett, \emph{Trading
Volatility}, free at \url{https://trading-volatility.com/Trading-Volatility.pdf}.

\begin{waste}
You already did it, and a better surface fit has never by itself made anyone
money --- you are polishing a lens, not finding something to look at. Worse,
Section~\ref{sec:residual} shows that on an unfiltered NSE chain your
``residual'' is substantially your model disagreeing with NSE's model.
\end{waste}

\subsubsection{2. Volatility risk premium harvesting (short vol)}
\fitem{Is}Selling implied against subsequently realised volatility,
systematically, hedged or unhedged. The most documented positive-expectancy
structure in options.
\fitem{Uniquely solves}The only options return stream that is a compensated risk
premium rather than a forecasting bet. It pays you for supplying insurance.
\fitem{Status}\dec{DECAYED in the US}, \alv{ALIVE and structurally distorted in
India}. SPX VRP is the most crowded retail-adjacent trade alive and 0DTE has
compressed it. Indian VRP is fatter, but the fat part is exactly the part that
kills you.
\fitem{Prereqs}Minimal maths. What you actually need is a margin model and a
survival-through-drawdown design. This is the one field where being small hurts.
\fitem{First useful result}3--4 weeks to a cost-honest backtest. Years to know
whether it survives.
\fitem{Solo?}Yes --- and the risk of ruin is the highest of anything on this map.
\fitem{Sources}CBOE PUT / BXM / CNDR methodologies and full history, the public
benchmark that makes an alpha claim falsifiable; SEBI's derivatives studies;
Moontower (Kris Abdelmessih, 21 years market-making at SIG and Parallax).

\begin{waste}
The edge is real, small and levered. Your backtest will show a beautiful Sharpe
that is entirely a description of what has not happened yet. Per SEBI's August
2026 behaviour study, the option sellers who lose lose roughly \rs 51.7 lakh
each --- about eleven times the buyers --- on a 44\% loss rate that looks
reassuring right up until it is not.
\end{waste}

\subsubsection{3. Dispersion and correlation trading}
\fitem{Is}Short index vol against long single-name vol, monetising the spread
between index-implied and constituent-implied correlation.
\fitem{Uniquely solves}It isolates \emph{correlation} as a tradeable asset. No
directional or single-name view survives the construction.
\fitem{Status}\alv{ALIVE but institutional} in the US; effectively
\dec{CLOSED in India}, and for structural reasons rather than liquidity ones.
Section~\ref{sec:dispersion} gives the measured verdict.
\fitem{Prereqs}Correlation estimation, vega-weighting, a margin engine, and in
India a physical-delivery operation across fifty names.
\fitem{First useful result}6--10 weeks to a paper study that will not be tradeable.
\fitem{Solo?}No.
\fitem{Sources}Bennett's dispersion chapters, the best free treatment in
existence; Bossu, \emph{A Practical Guide to Dispersion}.

\begin{waste}
You will build it, find that it works on paper, and then discover that Indian
single-stock options are physically settled, that the expiry calendars no longer
line up, and that SPAN gives you no correlation offset at all across
underlyings --- so you post margin as though you held fifty-one unrelated books.
\end{waste}

\subsubsection{4. Options market making}
\fitem{Is}Continuous two-sided quoting with inventory-aware skew, where P\&L is
edge per trade minus adverse selection minus hedging cost.
\fitem{Uniquely solves}The only options activity where you are paid the spread
rather than paying it. The sign of the cost term flips.
\fitem{Status}\alv{ALIVE in crypto}, \dec{DEAD to outsiders in Indian index
options}, where quoting is a colocation race against sub-100-microsecond stacks.
\fitem{Prereqs}Avellaneda-Stoikov and successors, queue models, real exchange
connectivity, inventory capital. Colocation for anything Indian.
\fitem{First useful result}2--3 months on crypto at small size. Never, for NSE.
\fitem{Solo?}Crypto yes, at small size, on mid-cap perps. Indian index options no.
\fitem{Sources}\texttt{hummingbot}; \texttt{hftbacktest}
(\url{https://github.com/nkaz001/hftbacktest}), which models queue position and
feed \emph{and} order latency separately; Avellaneda \& Stoikov (2008) for the
reservation-price logic, not the parameters.

\begin{waste}
Market making is 10\% model and 90\% infrastructure, operations and exchange
relationships. The part you enjoy is the part that matters least.
\end{waste}

\subsubsection{5. Options flow and dealer positioning (GEX, vanna, charm)}
\fitem{Is}Inferring the dealer hedging book from open interest and trade-side
classification, then predicting the mechanical hedging flow it forces.
\fitem{Uniquely solves}The only framework that predicts \emph{flow} --- a
forced, non-discretionary buyer or seller --- rather than value.
\fitem{Status}\dec{Half-decayed in the US} and, as ordinarily constructed,
\dec{sign-inverted on NIFTY}. Round two refuted the standard construction against
NSE's own published data; see Section~\ref{sec:gex}. What remains is a harder and
considerably more defensible project.
\fitem{Prereqs}Trade-side classification --- the hard part --- open-interest
reconciliation, and a surface for the greeks. Capital: low.
\fitem{First useful result}4--6 weeks, and only if you treat the dealer sign as a
parameter to be estimated rather than an assumption to be imported.
\fitem{Solo?}Yes, and the Indian version is genuinely unpublished.
\fitem{Sources}NSE participant-wise open interest, daily, signed by category, at
\texttt{nsearchives.nseindia.com/content/nsccl/fao\_participant\_oi\_DDMMYYYY.csv};
Barbon \& Buraschi, \emph{Gamma Fragility} --- the free, named, testable version
of the claim Karsan makes informally.

\begin{waste}
The whole edifice rests on guessing which side of each trade the dealer took.
Get that classification wrong and you have built a very sophisticated random
number generator --- and on NIFTY the standard guess is demonstrably backwards.
\end{waste}

\subsubsection{6. Event-driven volatility}
\fitem{Is}Trading the term-structure kink around a scheduled event. Implied vol
embeds an event premium that decays deterministically once the event passes.
\fitem{Uniquely solves}A dated, non-stationary edge with a known catalyst --- the
only options trade where you know in advance when you find out you were wrong.
\fitem{Status}\alv{ALIVE}. The Indian equivalents --- RBI policy, budget day,
election day --- are considerably less picked over than US mega-cap earnings.
\fitem{Prereqs}A clean event calendar, an event-vol extraction model, and brutal
cost discipline. These are wide-spread trades.
\fitem{First useful result}3--5 weeks.
\fitem{Solo?}Yes.
\fitem{Sources}Bennett's event chapters; Dubinsky \& Johannes for the
decomposition (academic, no cost model, so structure only); Klassen's
\emph{Event Modeling} talk for the two-jump Merton treatment.

\begin{waste}
The premium is real, and the bid-ask is wider than the premium in exactly the
names where the premium is largest.
\end{waste}

\subsubsection{7. Delta hedging, discrete-hedging error, and P\&L attribution}
\fitem{Is}The mechanics of turning an options position into a volatility
position: hedge frequency, band choice, and the decomposition of realised P\&L
into theta, gamma, vega, vanna and residual.
\fitem{Uniquely solves}The only thing that tells you \emph{why} you made money.
Without attribution, every options result is unattributable luck.
\fitem{Status}\alv{ALIVE, and permanently so.} Not an alpha field --- it is the
accounting layer that makes every other field legible, and the single most common
thing solo traders skip.
\fitem{Prereqs}Minute data, which you are about to have; a greeks engine; a cost
model.
\fitem{First useful result}\textbf{One to two weeks. The fastest payoff in this
entire document.}
\fitem{Solo?}Trivially.
\fitem{Sources}Section~\ref{sec:measure} in full.

\begin{waste}
It generates no signal whatsoever. It only tells you the truth about signals you
already have, which is unglamorous and therefore usually skipped.
\end{waste}

\subsection{Microstructure, execution, speed}

\subsubsection{8. Market microstructure (order-book dynamics, adverse selection, queue)}
\fitem{Is}How price forms from the order book: queue priority, order-flow
imbalance, trade classification, and the information content of quotes.
\fitem{Uniquely solves}The only field that explains \emph{why your fill differed
from your backtest}. The gap between paper and live P\&L lives entirely here.
\fitem{Status}\alv{ALIVE and undervalued by solo traders}. Order-flow imbalance
as a one-to-five-second predictor still works; it is not tradeable without low
latency, but it is usable to improve entry timing on slower strategies.
\fitem{Prereqs}Level-2 or level-3 data. Your minute data is too coarse for this
--- it needs tick. Plus serious data engineering.
\fitem{First useful result}6--8 weeks, gated entirely on data acquisition.
\fitem{Solo?}For research and execution improvement, yes. As standalone alpha, no.
\fitem{Sources}\texttt{hftbacktest}; Cont, Kukanov \& Stoikov, \emph{The Price
Impact of Order Book Events}; Bouchaud et al., \emph{Trades, Quotes and Prices}
--- and its free predecessor, Bouchaud, Farmer \& Lillo, \emph{How Markets Slowly
Digest Changes in Supply and Demand} (arXiv:0809.0822), roughly 100 pages
covering a large fraction of the book.

\begin{waste}
All the real signal lives at horizons shorter than your latency. You will learn
a great deal and be able to act on almost none of it.
\end{waste}

\subsubsection{9. Execution algorithms and transaction cost analysis}
\fitem{Is}Scheduling a parent order into children to trade off market impact
against price risk, and measuring the result against benchmarks.
\fitem{Uniquely solves}It converts a theoretical edge into a realised one. For a
strategy with 5\,bp of edge and 4\,bp of cost, execution \emph{is} the strategy.
\fitem{Status}\alv{ALIVE, permanently.} Not crowded --- it is a cost-reduction
discipline, and cost reduction cannot be arbitraged away.
\fitem{Prereqs}Almgren-Chriss, an impact model calibrated to your own
instruments, and honest fill logging.
\fitem{First useful result}2--3 weeks for TCA on your own fills; months for a
calibrated impact model.
\fitem{Solo?}Yes, and disproportionately valuable at small size in Indian
options, where spreads dominate.
\fitem{Sources}Almgren \& Chriss, \emph{Optimal Execution of Portfolio
Transactions}; Almgren et al., \emph{Direct Estimation of Equity Market Impact}.

\begin{waste}
At your size your orders have no market impact, so most of the literature's
machinery is irrelevant. Only the spread-crossing decision matters.
\end{waste}

\subsubsection{10. High-frequency trading and latency arbitrage}
\fitem{Is}Competing on speed for stale-quote and cross-venue discrepancies,
measured in microseconds.
\fitem{Uniquely solves}Captures the mechanically riskless component of price
dislocation. No forecast required.
\fitem{Status}\dec{DEAD to you.} Not decayed --- structurally inaccessible. It is
a capital-expenditure race: colocation, FPGA, microwave. There is no
evening-hours version.
\fitem{Prereqs}Seven figures of infrastructure and an exchange membership.
\fitem{First useful result}Never, as a solo.
\fitem{Solo?}Absolutely not.
\fitem{Sources}Budish, Cramton \& Shim, \emph{The High-Frequency Trading Arms
Race}, which is the best argument for not entering; Aquilina, Budish \& O'Neill
for the FCA's message-level measurement.

\begin{waste}
Reading about it is entertaining. There is a zero percent chance you can do it,
and the reading crowds out fields where you can.
\end{waste}

\subsection{Systematic equity and cross-asset}

\subsubsection{11. Statistical arbitrage and mean reversion}
\fitem{Is}Trading the residual of an asset against a hedged basket or partner,
betting on reversion of a spread that has no reason to trend.
\fitem{Uniquely solves}Market-neutral return with a mechanically defined hedge
--- the only way to make the market factor genuinely disappear.
\fitem{Status}\dec{Heavily decayed at daily frequency}, \alv{alive intraday}.
Classic daily pairs on liquid US equities has been flat for a decade.
\fitem{Prereqs}Cointegration, Ornstein-Uhlenbeck, Kalman filtering. In India,
cash-market shorting is intraday-only and the stock-lending market is thin, so
you cannot hold the short leg.
\fitem{First useful result}3--4 weeks to a backtest, and the backtest will be
wrong by roughly the size of the edge.
\fitem{Solo?}At small scale, partially.
\fitem{Sources}Avellaneda \& Lee, \emph{Statistical Arbitrage in the U.S.
Equities Market}, which includes an explicit decay analysis; Robot Wealth, who
publish negative results; Robert Carver, \url{https://qoppac.blogspot.com} and
\texttt{pysystemtrade}, where costs are a first-class object in the codebase.

\begin{waste}
Every published stat-arb result is a lower bound on how crowded it now is, and
the surviving edge is smaller than the Indian STT-plus-impact cost stack.
\end{waste}

\subsubsection{12. Factor investing and cross-sectional equity}
\fitem{Is}Sorting a cross-section on characteristics and holding the long-short
spread.
\fitem{Uniquely solves}Long-horizon compounding with capacity --- the only field
here that scales to serious assets under management.
\fitem{Status}\dec{DECAYED as alpha}, \alv{alive as a benchmark}. The
replication crisis gutted most of the zoo. Momentum and quality survive weakly.
\fitem{Prereqs}Panel-data hygiene and a survivorship-bias-free universe, which is
the real barrier in India (see Section~\ref{sec:eligibility}).
\fitem{First useful result}4--6 weeks, mostly spent cleaning data.
\fitem{Solo?}Yes, but the payoff horizon is years and the Sharpe is around 0.4.
\fitem{Sources}Hou, Xue \& Zhang, \emph{Replicating Anomalies} --- emphatically a
list of things that do not work; Harvey, Liu \& Zhu; AQR's published library.

\begin{waste}
It is slow, low-Sharpe, capacity-oriented investing dressed as quant research ---
the opposite of a minute-wise options trader's comparative advantage.
\end{waste}

\subsubsection{13. Trend following and time-series momentum}
\fitem{Is}Long-horizon directional betting on the sign of past returns across a
diversified futures basket, sized by inverse volatility.
\fitem{Uniquely solves}Positive skew and crisis convexity. The only systematic
strategy that reliably makes money when everything else is dying --- directly
relevant to a firm called Convex Hedge.
\fitem{Status}\alv{ALIVE but degraded}. Sharpe has fallen from roughly 1.0
pre-2008 to 0.4--0.5. The convexity is the point, not the Sharpe.
\fitem{Prereqs}Low maths, high diversification requirement --- 30-plus
instruments, which means global futures. Capital is the binding constraint.
\fitem{First useful result}2 weeks with \texttt{pysystemtrade}.
\fitem{Solo?}If you can access global futures. Borderline India-only.
\fitem{Sources}Carver (ex-AHL, runs it live, publishes his own failures);
Moskowitz, Ooi \& Pedersen; Hurst, Ooi \& Pedersen for the century-long evidence.
See also Section~\ref{sec:trendopts} for the trend-options identity, which is the
part that connects to your actual work.

\begin{waste}
It needs breadth and patience, both of which you currently lack. A trend book on
three Indian instruments is not trend following; it is gambling with extra steps.
\end{waste}

\subsubsection{14. Portfolio construction and risk (covariance, sizing, Kelly)}
\fitem{Is}Turning a set of signals into position sizes: covariance estimation,
shrinkage, risk parity, hierarchical risk parity, leverage choice.
\fitem{Uniquely solves}It determines the \emph{geometric} return, which no signal
research touches. Two traders with identical signals and different sizing have
unrelated outcomes.
\fitem{Status}\alv{ALIVE, permanently, and systematically underrated.} Not
crowded, because it is not a trade.
\fitem{Prereqs}Linear algebra, Ledoit-Wolf shrinkage, and an understanding of why
mean-variance optimisers explode.
\fitem{First useful result}2--3 weeks, and it immediately improves everything
else you build.
\fitem{Solo?}Yes.
\fitem{Sources}Ledoit \& Wolf; L\'opez de Prado on hierarchical risk parity, free
on SSRN with code in the paper; and Section~\ref{sec:sizing} for the
options-specific treatment, which is the version you actually need.

\begin{waste}
Optimal sizing of a nonexistent edge is exactly zero. It is very easy to spend
three months here instead of finding the edge.
\end{waste}

\subsubsection{15. Research methodology and backtest hygiene}
\fitem{Is}The statistics of strategy selection: multiple-testing correction,
deflated Sharpe, walk-forward design, and cost modelling as a first-class citizen.
\fitem{Uniquely solves}The only defence against the dominant failure mode in solo
quant work --- believing a backtest.
\fitem{Status}\alv{ALIVE, and the highest return per hour spent on this map.}
\fitem{Prereqs}Statistics you already have. Discipline you may not.
\fitem{First useful result}One weekend to implement probability of backtest
overfitting and a trial counter.
\fitem{Solo?}Yes.
\fitem{Sources}Bailey \& L\'opez de Prado on the deflated Sharpe ratio and on
combinatorially symmetric cross-validation. Note the assessment in
Section~\ref{sec:backtest}: implement \textbf{PBO via CSCV}, not DSR --- DSR needs
the number of trials, and nobody honestly reports their own.

\begin{waste}
It only ever tells you that things you liked do not work. It produces no
strategies, only funerals.
\end{waste}

\subsection{Machine learning and data}

\subsubsection{16. Machine learning applied to markets}
\fitem{Is}Supervised and reinforcement models on financial features, with the
special problems of non-stationarity, low signal-to-noise, overlapping labels and
near-zero effective sample size.
\fitem{Uniquely solves}Nonlinear interaction discovery across many weak features,
genuinely hard to do by hand once you have more than twenty.
\fitem{Status}\dec{Mostly decayed at the solo level}, \alv{alive narrowly}. It
wins where the label is short-horizon and microstructural --- next-100ms mid
move, fill probability, queue depletion --- and loses almost everywhere else.
Deep learning on daily OHLCV is a well-documented dead end.
\fitem{Prereqs}You have the maths. The real prerequisite is purged and embargoed
cross-validation, and enough independent samples --- which minute options data
gives you and daily data does not.
\fitem{First useful result}6--10 weeks, most of it spent discovering leakage.
\fitem{Solo?}For narrow, well-posed problems, yes. End-to-end, no.
\fitem{Sources}L\'opez de Prado's free SSRN corpus --- triple-barrier labelling,
meta-labelling, purged CV, fractional differentiation. The book is paid; the
constituent papers are free and carry nearly all the substance.

\begin{waste}
The most seductive and least productive field here for a solo trader. You will
build impressive pipelines whose out-of-sample Sharpe is indistinguishable from a
moving-average crossover, and the sunk cost will make you keep going.
\end{waste}

\subsubsection{17. Alternative data}
\fitem{Is}Non-price inputs --- satellite, card spend, web scrape, filings text
--- mapped to fundamentals ahead of consensus.
\fitem{Uniquely solves}True informational rather than statistical edge: actually
knowing something before the market does.
\fitem{Status}\dec{Decayed at the top}, \alv{alive in obscure corners}. Every
mainstream dataset is licensed by two hundred funds and priced accordingly. What
is alive for a solo is \emph{self-collected} data nobody sells.
\fitem{Prereqs}Scraping, ETL, storage, legal care.
\fitem{First useful result}Three months or more, and mostly engineering.
\fitem{Solo?}Only for self-collected data.
\fitem{Sources}Round one could not verify a single source in this field that
cleanly passes the filter and is free --- most of the literature is vendor
marketing. It is reported as \textbf{unsourced} rather than padded. The one
exception worth naming: NSE's own free participant-wise and open-interest
archives are alternative data that almost nobody parses properly.

\begin{waste}
For a solo trader this is 90\% data engineering and 10\% finance, and the 10\% is
only valuable if you have a stock-picking process to feed --- which you do not.
\end{waste}

\subsubsection{18. Research and data infrastructure}
\fitem{Is}The unglamorous layer: tick capture, a time-series store, point-in-time
correctness, corporate-action and contract-roll handling, deterministic
backtest-live parity.
\fitem{Uniquely solves}The only thing that makes results reproducible and
prevents look-ahead. Minute-wise Indian options data is large and messy; without
this layer nothing else is trustworthy.
\fitem{Status}\alv{ALIVE and mandatory.} Not a crowding question.
\fitem{Prereqs}Engineering. Parquet/Arrow, columnar storage, and a symbology
scheme for expiring option contracts.
\fitem{First useful result}3--4 weeks, and it gates every other field.
\fitem{Solo?}Yes. Section~\ref{sec:storage} gives a specific recommendation so
you can stop thinking about it.
\fitem{Sources}NautilusTrader, for nanosecond event-driven design with identical
strategy code in backtest and live; \texttt{hftbacktest}; ArcticDB.

\begin{waste}
Infrastructure is infinitely extensible and has no natural stopping point. You
can spend a year building a platform and zero days trading.
\end{waste}

\subsection{Crypto}

\subsubsection{19. Perpetual funding, basis, and cash-and-carry}
\fitem{Is}Harvesting the funding rate on perpetual swaps by holding the
offsetting spot or dated future, and trading the basis term structure.
\fitem{Uniquely solves}A mechanically defined, observable, high-frequency carry
stream that exists nowhere else in finance. The funding rate is \emph{printed}
every eight hours.
\fitem{Status}\alv{ALIVE but compressed.} BTC and ETH basis is heavily farmed by
ETF-arbitrage desks; annualised carry has fallen from 20\%-plus to single digits
in calm regimes. What remains alive: altcoin funding, funding dislocations during
liquidation cascades, and cross-exchange funding spreads.
\fitem{Prereqs}Low maths, high operational competence --- exchange API risk,
margin and liquidation modelling, custody, counterparty risk.
\fitem{First useful result}2--3 weeks. The fastest path to a live, cost-honest,
small-capital P\&L on this map.
\fitem{Solo?}The most solo-viable strategy here --- but see the tax verdict in
Section~\ref{sec:vda}, which reverses the conclusion for an Indian resident.
\fitem{Sources}Hummingbot's \texttt{v2\_funding\_rate\_arb.py} --- read issue
\#7295 first, which documents a double fill on Hyperliquid perps after the
strategy is stopped; \emph{Anatomy of the Oct 10--11 2025 Crypto Liquidation
Cascade} (SSRN 5611392) as the stress scenario your backtest must survive.

\begin{waste}
The carry is small and the tail is exchange insolvency, oracle failure, or a
liquidation cascade that takes your hedge leg out. You are short a catastrophe
option to earn about 6\%. On 10--11 October 2025, auto-deleveraging forcibly
closed market makers' short perp legs against their long spot --- the exchange
took away the profitable side of the delta-neutral book. That has no analogue on
NSE, and any risk framework you carry over will not model it.
\end{waste}

\subsubsection{20. Crypto options and volatility}
\fitem{Is}Volatility trading on Deribit and the Indian-registered venues, plus
the DeFi analogue of option vaults and AMM liquidity provision as a short-gamma
position.
\fitem{Uniquely solves}The place where your existing SVI/SSVI machinery transfers
immediately to a market with fat spreads, no colocation requirement, 24/7 data,
and few sophisticated participants relative to SPX.
\fitem{Status}\alv{ALIVE --- the least crowded serious options market a solo can
reach}, with two large caveats measured in Section~\ref{sec:crypto}: the fee cap
makes cheap options untradeable, and the Indian tax regime makes deployment
uneconomic.
\fitem{Prereqs}Exactly what you already have, plus API plumbing and the
inverse-contract correction.
\fitem{First useful result}4--6 weeks. Your SVI work is most of the way there.
\fitem{Solo?}Yes. Strongest fit-to-skill ratio on the entire map --- as
\emph{research}. Not as a P\&L venue.
\fitem{Sources}Echenim, Gobet \& Maurice (Kaiko), which proves that calibrating
SVI on crypto mid-prices systematically underestimates wing vol and gives an
anchor price to replace the mid; Magadini, Sinclair \& Chepal on real-time SVI
calibration; Lucic \& Sepp on inverse options with code at
\url{https://github.com/ArturSepp/StochVolModels}.

\begin{waste}
Deribit's liquidity is concentrated in a handful of BTC and ETH expiries. Outside
them, the ``opportunity'' you find in the data is a quote nobody will trade
against. And on sub-\$200 options you pay 25\% round-trip in fees plus a 24\%
spread, so your edge must exceed roughly half the premium.
\end{waste}

\subsubsection{21. On-chain analytics and MEV}
\fitem{Is}Reading blockchain state directly --- DEX pool composition,
liquidations, mempool --- and, for MEV, competing to reorder transactions.
\fitem{Uniquely solves}Complete, free, real-time visibility into a market's
entire order flow and positioning, a condition that exists in no traditional market.
\fitem{Status}MEV is \dec{decayed into a relay oligopoly}; on-chain
\emph{analytics} is \alv{alive and cheap}. The measurement is decisive: over
2023--2025, three firms captured about 90\% of extracted value by Q1 2025, and
they transfer roughly 90\% of arbitrage revenue to their integrated builders to
buy auction priority --- a subsidy you cannot match. One tracked searcher ran at
$-272\%$ margin.
\fitem{Prereqs}EVM literacy and an indexer. MEV additionally needs colocated
infrastructure and a relay relationship.
\fitem{First useful result}Analytics, 3--4 weeks. MEV, do not.
\fitem{Solo?}Analytics yes; MEV effectively no.
\fitem{Sources}\emph{Measuring CEX-DEX Extracted Value and Searcher
Profitability} (arXiv:2507.13023) --- a measurement paper that will not teach you
to build a searcher but will save you six months of trying.

\begin{waste}
It is a different discipline wearing a finance costume. You would be learning
distributed systems engineering, not trading. And the barrier is not skill, it is
a commercial builder relationship.
\end{waste}

\subsubsection{22. Volatility harvesting and the rebalancing premium}
\fitem{Is}Rebalancing between a volatile asset and cash to capture the difference
between arithmetic and geometric mean return.
\fitem{Uniquely solves}Nothing, for you.
\fitem{Status}\dec{Retail folklore, and the most misunderstood item here.}
Shannon's Demon is real arithmetic, but the ``rebalancing bonus'' is payment for
\emph{selling variance and providing liquidity} --- the payoff is short-gamma and
mean-reversion-dependent, and it inverts in a trending market.
\fitem{Prereqs}None, which is the problem.
\fitem{First useful result}Immediately, and it will be wrong.
\fitem{Solo?}Yes, and that is the problem.
\fitem{Sources}Witte, \emph{Volatility Harvesting} (arXiv:1508.05241) --- read it
to get the maths right, then stop.

\begin{waste}
In crypto it fails twice over: taker fees on every rebalance, and BTC's realised
drift has historically dominated the harvest term, making it an expensive way to
underperform holding. Note also the structural conflict with trend following ---
they are the same bet in opposite directions, and running both without an
explicit regime allocator means paying two spreads to hold no position.
\end{waste}

\subsection{India-specific}

\subsubsection{23. Indian market structure and the regulatory regime}
\fitem{Is}The rules that determine whether an Indian strategy is viable at all:
STT, stamp duty, position limits, SPAN and ELM margining, lot sizes, and the
expiry calendar.
\fitem{Uniquely solves}The only thing that converts a US-derived technique into
something that survives contact with NSE. Costs in India are not a rounding
adjustment; they are a structural asymmetry that changes which structures are viable.
\fitem{Status}\alv{ALIVE and actively changing.} This is time-sensitive
knowledge. Section~\ref{sec:breaks} lists five structural breaks in the last
twenty-four months, at least three of which fall inside your data.
\fitem{Prereqs}None mathematical. Pure diligence.
\fitem{First useful result}One week to build a correct cost-and-calendar model.
It will invalidate several strategies immediately, which is the point.
\fitem{Solo?}Mandatory, actually.
\fitem{Sources}SEBI circulars and NSE circulars, all free and machine-indexed ---
see Appendix~\ref{app:primary}. \textbf{Explicitly rejected:} Zerodha Varsity,
Sensibull education, and every broker ``learn'' portal. They are often factually
accurate on mechanics; use them as a mechanics reference, never as a source of a
conclusion.

\begin{waste}
It is compliance reading, not research, and it goes stale --- SEBI changes the
rules faster than you can encode them. The mitigation is architectural: put every
rate in a dated configuration file, never in source. See the audit in
Section~\ref{sec:indiacode} for what happens when you do not.
\end{waste}

\subsubsection{24. Indian index-options intraday structure}
\fitem{Is}The specific behaviour of NIFTY and Sensex options intraday: the open
volatility spike, midday decay, expiry-day gamma compression, and the retail flow
signature.
\fitem{Uniquely solves}The only field where your data advantage is a genuine edge
over the global research population, because almost nobody outside India studies
this and almost nobody inside India studies it rigorously.
\fitem{Status}\alv{ALIVE and genuinely under-researched.} Retail participation is
enormous, documented, and loses money --- so there is identifiable flow to be on
the other side of. The crowded part is the naive expiry-day short straddle, which
every algo desk in Mumbai is farming.
\fitem{Prereqs}Your existing stack plus the cost and calendar model from field 23.
\fitem{First useful result}3--4 weeks. The fastest path to something original.
\fitem{Solo?}Yes, and this is where solo is \emph{advantaged}.
\fitem{Sources}SEBI's FY25--FY26 studies for the shape of the losing flow; NSE
bhavcopy and participant-wise archives, free and primary.
\srcnote{Round two searched hard and could not find a single free, code-backed,
cost-explicit research source on Indian intraday options microstructure. That
vacuum is the opportunity, but it means nobody will hand you a template.}

\begin{waste}
``Under-researched'' and ``unprofitable'' look identical from the outside, and
there is a real chance the reason nobody published is that the costs eat
everything. Section~\ref{sec:coststack} gives you the number to check that
against before you start.
\end{waste}

\subsection{Adjacent}

\subsubsection{25. Tail hedging and convexity}
\fitem{Is}Deliberately paying negative carry for a payoff that is convex in a
crash --- deep OTM puts, VIX call spreads, long-vol overlays --- sized so the
cost is survivable.
\fitem{Uniquely solves}The only construction that pays off when correlations go
to one and every diversification assumption breaks at once.
\fitem{Status}\alv{ALIVE as a portfolio function}, \dec{decayed as a standalone}.
The famous 2020 results are survivorship-selected; most tail funds bleed. It works
only as an overlay whose size is chosen by portfolio maths, not conviction.
\fitem{Prereqs}Convexity decomposition and honest accounting of the bleed.
\fitem{First useful result}3--4 weeks --- and it is directly on-mission for
Convex Hedge.
\fitem{Solo?}As research yes; it needs a portfolio to overlay onto.
\fitem{Sources}Bennett on the cheapest way to buy a crash; Bhansali's free SSRN
papers on tail-hedging cost, candid about the drag; trend following as a cheaper
convexity source.

\begin{waste}
Tail hedges lose money 95\% of the time by design, and there is no backtest
length that will make you confident you sized it right.
\end{waste}

\subsubsection{Also considered, and set aside}
\textbf{Fixed income and rates volatility} --- the swaption vol cube is the only
vol market with a genuine third dimension and the modelling problem is strictly
richer than equity vol. It is also OTC, dealer-intermediated, with no retail
access and no obtainable data, so it fails test~(5) outright for you. Hagan et
al., \emph{Managing Smile Risk}, remains worth reading --- but for its diagnosis
of why model deltas are unstable, which transfers, rather than for SABR, which
is inferior to what you already have. \textit{Waste warning: you cannot get the
data or the access, so it is pure intellectual tourism.}

\subsection{The shortlist}

Ranked for your position specifically --- someone with the surface maths already
done, minute-wise NIFTY data arriving, a full-time job, and no realistic access
to US venues.

\begin{enumerate}
\item \textbf{Delta hedging, hedging bands, and P\&L attribution} (field 7).
Cheapest, fastest, and it is the measurement instrument for everything else. One
to two weeks. Without it every subsequent result you produce is unattributable.
Do this first, and in parallel with nothing.

\item \textbf{Realised volatility from minute data, and the deseasonalised-RV
versus remaining-implied-variance study} (fields 24 and 2). This is the project
where your data advantage is actually spent, it is invisible in end-of-day data,
and Section~\ref{sec:rvproject} specifies it step by step. It is also genuinely
publishable.

\item \textbf{Indian cost, calendar and margin modelling} (field 23). Not a
strategy --- a filter. It will kill several ideas you currently believe in, which
is worth more than any new idea. One week, and Section~\ref{sec:margin} has most
of it already computed.

\item \textbf{The SVI residual falsification test} (fields 1 and 24). Converts
your existing codebase from infrastructure into a falsifiable hypothesis. The
pre-registered specification is in Section~\ref{sec:residual}; the honest
expectation is that it fails, and finding that out in weeks rather than months is
the whole value.

\item \textbf{Dealer positioning on NIFTY, rebuilt properly} (field 5). Higher
risk than the others and a harder project than round one implied, but it is the
one intersection where your data, a free daily primary file with no US
equivalent, and a genuine research vacuum all coincide.
\end{enumerate}

\textbf{Explicitly deprioritised despite being interesting:} machine learning
applied to markets (highest seduction-to-yield ratio here), HFT and MEV
(structurally inaccessible), dispersion (structurally closed in India), factor
investing (wrong horizon for your data), rates vol (no data access), volatility
harvesting (folklore).

\newpage
\section{Four findings that reshape the plan}
\label{sec:constraints}

These came out of round two's verification pass. Each was asserted or assumed in
round one; each was checked against a primary document or against live data; and
each changes what is worth building rather than merely adding detail.

\subsection{The US is a source of technique, not a venue}
\label{sec:lrs}

\begin{keyfact}{Verdict: REFUTED --- and the received wisdom is wrong in a way that matters}
An Indian resident individual has \textbf{no LRS route to a US-listed option},
long or short. The chain closes without a gap, and it does not run through the
sentence everyone quotes.
\end{keyfact}

The derivation, from primary text:

\begin{enumerate}
\item \textbf{RBI Master Direction on LRS} (FED Master Direction 7/2015-16),
para 6(iii): the exhaustive list of permissible \emph{capital account}
transactions runs only through Overseas Direct Investment and Overseas Portfolio
Investment.
\item \textbf{FEM (Overseas Investment) Directions, 2022}, definition (ix)(a):
``OPI shall not be made in \ldots\ \textbf{any derivatives unless otherwise
permitted by Reserve Bank}.''
\item A US-listed option is a derivative. Therefore no LRS capital-account route
exists. Separately, para 5 independently bars ``remittance for margins or margin
calls to overseas exchanges'', which kills short options by a second route.
\end{enumerate}

\textbf{The correction to the received wisdom.} Many advisory sites quote the RBI
LRS FAQ as saying ``remittance under the Scheme is not available for margin
trading or derivative transactions.'' Round two downloaded the full FAQ and
grepped it: the strings ``derivative'' and ``option'' appear \textbf{zero times}.
The conclusion is right; the citation everybody uses is fabricated. If you ever
argue this with a bank's authorised-dealer desk, cite definition (ix)(a)(iii) of
the OI Directions --- quoting a sentence that does not exist will lose the argument.

\textbf{GIFT City does not cure it.} The July 2024 circular widened IFSC
remittance to ``all permissible purposes under LRS'' --- which imports the same
OPI carve-out. IFSC relocates the account; it does not create a derivatives
permission. NSE International Exchange's own Global Access platform says so
directly, excluding derivatives and crypto ``due to regulatory restrictions under
LRS norms.''

\textbf{What this changes.} Everything you build must clear on NSE or on a
domestically registered crypto venue. This belongs at the top of your plan, not
in a footnote. The extensive US options literature in this document is technique
transfer only, and you should read it that way --- including the parameters,
which will not transfer even where the methods do.

\srcnote{Could not be verified: no RBI document says in terms ``a resident
individual may not trade overseas exchange-traded options.'' No such sentence
appears to exist, which is precisely why secondary sources invent one. Separately
unresolved and worth a professional opinion: funds already held abroad and never
remitted under LRS sit outside the remittance question but inside FEMA's
residency net.}

\subsection{Crypto tax applies transaction by transaction, and that is a strategy filter}
\label{sec:vda}

\begin{keyfact}{Verdict: CONFIRMED transaction-level. No netting. And the section number in circulation is obsolete.}
Section 115BBH no longer exists. The Income-tax Act, 2025 replaced the 1961 Act
with effect from 1 April 2026; the successor is \textbf{s.194(1), Table serial
number 4}, with withholding at \textbf{s.393(1), Table Sl.\ 8(vi)}.
\end{keyfact}

The operative text carries two independent prohibitions. Condition (a) forbids
set-off ``in computing the income referred to in column C'' --- that is what kills
same-asset netting. Condition (b) then forbids set-off against anything else and
bars carry-forward. Note that (b) says ``any provision'', not ``any \emph{other}
provision''; the narrowing word is absent, which closes the inter-asset route.

Three converging confirmations were found, and the third is the strongest because
it is not an interpretation at all:

\begin{itemize}
\item The statutory text itself.
\item A ministerial answer in the Lok Sabha (21 March 2022) stating directly that
loss from the transfer of one VDA will not be allowed to be set off against income
from transfer of another. \srcnote{Flagged: sansad.in returned HTTP 500 on every
URL pattern; this quotation is from search-result text of that answer, not from
the PDF.}
\item \textbf{The filing utility enforces it.} CBDT's ITR-2 Validation Rules for
AY 2025-26, Rules 201--202, require the Schedule VDA total field --- literally
labelled ``sum of all \emph{positive} incomes'' --- to equal the plain sum of
column 7. That can only hold if every row is non-negative. \textbf{The software
structurally cannot accept a negative row.} Reporting is transaction-wise.
\end{itemize}

\subsubsection*{What that does to the arithmetic}

Round two derived the general result rather than asserting it. Let a round trip
have P\&L $X$ with mean $\mu$, standard deviation $\sigma$, and define the
per-trade information ratio $r = \mu/\sigma$. Tax is levied on the positive part
only, with no cost deduction:

\[
\text{after-tax per trade} \;=\; \mu - 0.312\big(\mu\,\Phi(r) + \sigma\,\varphi(r)\big)
\]

Setting this to zero gives $r^{*} = 0.1491$. A trade must carry a per-trade
information ratio of about \textbf{0.149} before the regime lets you keep
anything. Equivalently, the annualised Sharpe needed merely to break even is
$0.149\sqrt{N}$:

\begin{center}
\begin{tabular}{@{}rr@{}}
\toprule
Trades per year & Annual Sharpe needed to break even \\
\midrule
250    & 2.4 \\
1{,}000  & 4.7 \\
5{,}000  & 10.5 \\
20{,}000 & 21.1 \\
50{,}000 & 33.3 \\
\bottomrule
\end{tabular}
\end{center}

The corresponding number under the F\&O regime is \textbf{zero}. Expressed as
effective tax on gross edge, the regime charges \textbf{79\%} at $r=0.2$,
\textbf{141\%} at $r=0.1$, and \textbf{265\%} at $r=0.05$. Only above roughly
$r \approx 1$ does the 31.2\% headline rate approximate reality.

\textbf{Every delta-neutral options strategy lives between $r = 0.02$ and
$r = 0.3$.} In that band this is not a 30\% tax; it is a 79\%-to-640\% tax.

The 1\% withholding compounds it on a separate axis. With gross edge $e$ as a
fraction of transfer consideration, withholding as a share of gross edge is
$0.01/e$ --- so the book breaks the moment gross edge falls below \textbf{1\% of
turnover}. Market making, basis capture, funding harvesting and gamma scalping
earn 1--30\,bp of turnover, so they break by a factor of three to a hundred on
the withholding alone, before any tax is assessed.

\begin{keyfact}{The conclusion, stated as arithmetic rather than opinion}
The F\&O regime is \emph{tax-transparent to strategy design}: break-even gross
edge equals the cost stack, and tax is a scalar applied afterwards. The VDA
regime is \emph{not} --- because the base is $\mathbb{E}[\max(X,0)]$ rather than
$\mathbb{E}[X]$, it changes the \emph{ranking} of strategies, and it is close to
maximally hostile to exactly the delta-neutral options work you are training to
do. Build the research stack, execution stack, margin simulator and cost model
against NSE. Treat crypto as a source of data and technique, and as a place to
validate models --- not as a P\&L venue.
\end{keyfact}

\srcnote{Genuinely unresolved: whether the 1\% withholding attaches to VDA
\emph{derivatives}. A cash-settled perpetual or option on a VDA is not, on the
face of the words, a transfer \emph{of} a VDA, and no CBDT circular resolving it
was found. The no-loss-offset half applies regardless. This is a gap in the law,
not a research failure.}

\subsection{There is no NIFTY term structure to trade}
\label{sec:noterm}

SEBI's August 2026 study --- published one week before this document, and
superseding the FY25 figures round one used --- reports the turnover
concentration by time to expiry:

\begin{center}
\begin{tabular}{@{}lrr@{}}
\toprule
Share of index-option turnover & FY25 & FY26 \\
\midrule
On the expiry day itself (0DTE) & 70\% & 59\% \\
Within one day of expiry        & 80\% & 75\% \\
Within seven days of expiry     & 98\% & 97\% \\
Beyond ten days                 & \multicolumn{2}{c}{about 1\%} \\
\bottomrule
\end{tabular}
\end{center}

\textbf{Your surface has, in practice, one to two usable maturity points.}
Calendars, diagonals, term-structure relative value and any SSVI calendar-arbitrage
work have almost nothing to trade on this venue. This kills a whole class of ideas
cheaply, which is the most valuable thing a survey can do, and it should reshape
what you build before you write a line of code.

It compounds with a second structural fact: BANKNIFTY, FINNIFTY and MIDCPNIFTY
weeklies were withdrawn in November 2024, so BANKNIFTY is monthly-only and has
roughly three usable maturities. Any term-structure research on it is
structurally unavailable, not merely thin.

\subsection{Half the NSE option chain is model output, not market data}
\label{sec:imputed}

This is the finding that decides whether your SVI work contains a trade, and
round two settled it by downloading the NSE F\&O bhavcopy for 27 August 2026
(NIFTY spot 24{,}090.85; 1{,}590 NIFTY option lines across 18 expiries) rather
than arguing about it.

\begin{center}
\begin{tabular}{@{}lrrr@{}}
\toprule
& $n$ & Settlement on the \rs 0.05 tick grid & Close $\neq$ settlement \\
\midrule
Traded (volume $>$ 0)  & 769 & 90.5\% & 12\% \\
Untraded (volume $=$ 0) & 821 & \textbf{20.6\%} & \textbf{100\%} \\
\bottomrule
\end{tabular}
\end{center}

\textbf{821 of 1{,}590 lines --- 51.6\% --- did not trade at all}, and
\textbf{79.4\% of untraded settlement prices sit off the five-paise tick grid
entirely}, so they cannot be observed prices. NSE's contract specifications
confirm the mechanism: an option not traded during the day is closed at its
\textbf{Black-Scholes theoretical price}.

The imputation signature is visible directly. On the 29 September 2026 expiry,
four consecutive zero-volume call strikes returned implied vols of
\textbf{17.30, 17.21, 17.13, 17.00} --- a near-constant, monotone, near-flat
sequence sitting five to six vol points above their traded neighbours at 11.22\%
and 14.11\%. That is not a smile. That is a flat-volatility model writing its own
answer into the data.

\begin{waste}
Fit SVI to an unfiltered NSE chain and you will fit NSE's constant-volatility
assumption in the wings, then discover a ``residual'' that is your model
disagreeing with theirs. Any surface work that uses end-of-day closes without
filtering on traded volume is partly fitting NSE's own Black-Scholes output back
to itself --- which is why your wing fits will look far better than reality warrants.
\end{waste}

\textbf{And the tick-quantization explanation is wrong.} Round one blamed tick
size. Round two computed the implied-vol uncertainty from half a tick
($\Delta\text{IV} = 0.025/\text{vega}$) across the whole liquid wing of the
five-day expiry: it runs \textbf{0.00 to 0.17 vol points}, against an IV range of
6.4\% to 33\% across that same wing. A half-tick is under 0.5\% of the structure
being measured. Zero contracts settled at or below \rs 0.05 on the front expiry,
so the ``zero bid'' concern does not bind at end of day either. Right conclusion,
wrong mechanism --- and the distinction matters, because it tells you the fix is
a \emph{volume filter}, not a wider error bar.

\subsection{The structural breaks in your data}
\label{sec:breaks}

At least five, and three of them fall inside the period you are about to receive.
Any backtest spanning one of these without explicit regime-splitting is pooling
two different markets.

{\small
\setlength{\LTpre}{8pt}\setlength{\LTpost}{8pt}
\begin{longtable}{@{}p{1.9cm}p{4.7cm}p{9.0cm}@{}}
\toprule
\textbf{Date} & \textbf{Change} & \textbf{What breaks if you ignore it} \\
\midrule
\endfirsthead
\multicolumn{3}{@{}l}{\footnotesize\itshape Structural breaks, continued}\\[2pt]
\toprule
\textbf{Date} & \textbf{Change} & \textbf{What breaks if you ignore it} \\
\midrule
\endhead
\bottomrule
\endlastfoot
20 Nov 2024 & Contract size raised to \rs 15--20 lakh; weeklies cut to one
benchmark index per exchange; +2\% ELM on short options on expiry day &
BANKNIFTY / FINNIFTY / MIDCPNIFTY weeklies \emph{terminate} on 13, 19 and 18 Nov
2024. The series does not thin out --- it ends. Margin on expiry day rises by
rule. \\[2pt]
1 Feb 2025 & Upfront premium collection; calendar-spread margin benefit removed
on expiry day & A position that is margin-efficient for six days becomes
margin-expensive on the seventh, predictably and by rule. \\[2pt]
1 Sep 2025 & NSE expiry moves Thursday $\to$ Tuesday; BSE moves Tuesday $\to$
Thursday & Day-of-week seasonality, DTE mapping and expiry-day gamma all shift.
Cross-exchange work spanning this date crosses a \emph{swap}, not a one-sided
move. \textbf{And the expiry dates of already-live contracts were rewritten in
the contract master on 28 Aug 2025} --- so if you rebuild your universe from
today's masters, contracts that traded for weeks under Thursday expiry will
present as Tuesday. Snapshot the contract master per date; never reconstruct it. \\[2pt]
28 Oct 2025 & Lot sizes revised: NIFTY 75$\to$65, BANKNIFTY 35$\to$30, FINNIFTY
65$\to$60, MIDCPNIFTY 140$\to$120 & Every notional, every margin, every per-lot
cost. NIFTY's full history is 25 $\to$ 75 (Nov 2024) $\to$ 65. A P\&L-per-lot
series crossing these is wrong by 3$\times$ and then 13\%. \\[2pt]
1 Apr 2026 & STT raised again: options 0.10\%$\to$\textbf{0.15\%}, futures
0.02\%$\to$\textbf{0.05\%}, exercise 0.125\%$\to$\textbf{0.15\%} & The options
round-trip cost floor rose 27\%; futures rose 46\%. Any cost model carrying the
Oct-2024 rates is one revision stale. \\[2pt]
3 Aug 2026 & \textbf{Closing Auction Session} live, applied first to exactly the
stocks that have derivatives & The NIFTY close --- and therefore the final
settlement value of every index option --- is now derived from single-print
auction equilibrium prices of the constituents, not from any VWAP. CAS runs
15:15--15:35; equity derivatives trade to 15:40. Expiry-day strategies calibrated
on pre-August-2026 data are calibrated on a mechanism that no longer exists. \\
\end{longtable}
}

\begin{keyfact}{The settlement rule, resolved --- and it is neither answer round one had}
NSE Clearing states it plainly and the two rules differ: \textbf{index}
derivatives settle on the \textbf{close price of the underlying index};
\textbf{single-stock} derivatives settle on the \textbf{last-30-minute VWAP of
the security}. The widely repeated claim that NIFTY settles on a 15:00--15:30
VWAP of the index is \textbf{wrong}. The confusion arises because the index close
was historically computed from constituent closes, each of which \emph{was} a
last-30-minute VWAP --- but a weighted sum of per-stock VWAPs is not the VWAP of
the index series. If you have been reconstructing settlement by VWAP-ing the
index tick series, you have computed the wrong number for every expiry. And since
3 August 2026 even that is superseded by the auction.
\end{keyfact}

\newpage
\section{The volatility surface, pointed forward}
\label{sec:surface}

You can already fit SVI and SSVI. This section is about the two things you cannot
do yet: turn a fitted surface into a position, and know whether the fit means
anything on the chain you actually have.

\subsection{From a fitted surface to an actual trade}

\subsubsection{The definition of ``rich'' that survives contact with a market}

The single most useful reframing came from Kris Abdelmessih (SIG, then a
relative-value commodities vol book at Parallax): fair value is \emph{where you
can actually transact}, not where your model says. A one-vol residual inside a
three-vol-wide market is not an edge; it is a quote.

That converts the residual question from ``how big is the mispricing'' into ``how
big is the mispricing relative to the width of the market at that strike'', which
is a different and much harsher test.

\subsubsection{The error bar is the whole game}

Timothy Klassen's Vola Dynamics material is the best free source on this, and its
central design decision is the one to copy: \textbf{attach an error bar to every
fitted vol}, widening automatically as quote quality degrades. A tradeable
residual is one that exceeds the fitter's own uncertainty at that strike --- not
one that exceeds a fixed two-vol-point rule.

Round two added a second term that round one omitted entirely, and it matters:

\begin{keyfact}{The error bar has two components, not one}
An SVI slice fitted to $n$ quotes has \emph{parameter} uncertainty as well as
\emph{quote} uncertainty. The residual at strike $k$ must be compared against the
combined quote uncertainty and fit uncertainty at $k$. Comparing against quote
uncertainty alone is exactly how a fitted-surface artefact gets mistaken for
signal.
\end{keyfact}

Klassen is also explicit that \textbf{the mid is the wrong input in the wings}.
His KOSPI example shows puts quoting 0.02 by 0.03 next to 0.01 by 0.02 --- tick
quantization that makes wing mid-IVs pure artefact. Zero-bid options drag mid-IV
down in both wings; a naive fitter chases it, a robust one does not. Low-latency
desks feed size-weighted or book-depth microprices with temporal smoothing instead.

And the observation that should make you cautious about your own data source:
\textbf{settlement prices for illiquid series are themselves produced by a fitted
surface}, so a residual measured against exchange settlement data is partly
measuring someone else's fitter. Section~\ref{sec:imputed} shows this is not
hypothetical on NSE --- it is 51.6\% of the chain.

\subsubsection{What the structures actually express}

Bennett's \emph{Trading Volatility} (free PDF; the author ran derivatives
strategy at Santander and Barclays) is the structural half: when a calendar
expresses a term-structure view and when it is merely a vega/theta mismatch; why
risk reversals and put-spread financing behave differently as skew moves; how a
skew position bleeds through spot moves rather than through skew itself.

Its limits are worth stating: European institutional index vol, roughly 2012
vintage, no code, no cost model worth the name, and nothing on weeklies or 0DTE.
It tells you what trade to put on, never whether the edge survives your spread.

\begin{waste}
The academic option-return-predictability literature looks like the missing link
here and is not. The best of it (Bali et al., RFS 2023, 12 million delta-hedged
observations, 265 characteristics) reports results gross \emph{and} net of costs
--- but the effective-versus-quoted-spread assumption does most of the work in
the ``profits survive costs'' claim, and at full quoted spread much of it does
not. Treat it as a cost-sensitivity template, not a strategy. It is monthly
rebalancing on US single-name options, which you cannot trade anyway.
\end{waste}

\subsection{Does SVI even work on NIFTY weeklies?}

\begin{keyfact}{The finding first: nobody has published this, and that is real open ground}
An extensive search for practitioner work, Indian quant blogs and academic
literature applying SVI or SSVI to NIFTY, BankNifty or Sensex turned up: Indian
academic IV literature that is descriptive rather than parametric; one hackathon
repository; one first-hand LinkedIn field report; and Vola Dynamics inviting you
to ``contact us'' about fitting NIFTY, i.e.\ the people who actually solve this
treat it as a commercial product and publish nothing.
\end{keyfact}

The one genuine field report is worth quoting for its texture: a developer who
built a live NIFTY SVI surface reports that NSE chain data is ``far messier than
I initially assumed'', strikes sparse or stale, raw fits jagged and unstable near
spot, so he warm-starts each slice from the previous solution to stop
tick-to-tick jitter and weights ATM strikes more heavily. Verifiable person, no
institutional backing, no cost model. A field report, not a result --- but it
matches everything else here.

\subsubsection{Four structural facts that decide the answer}

\begin{enumerate}
\item \textbf{The forward breaks before SVI ever runs.} NIFTY weekly options have
no matching weekly future. Take the forward from the nearest monthly future --- or
worse from spot with a dividend-yield guess --- and the error enters as an
apparent put-call IV asymmetry that looks exactly like a skew dislocation.
\textbf{Imply the forward per expiry from put-call parity on the liquid near-ATM
strikes, and fit OTM-only.} Otherwise your best-looking residuals are your own
forward error.

\item \textbf{Wing quotes are tick-quantized to \rs 0.05.} A NIFTY weekly put
quoting 0.05 by 0.10 is a 67\% relative spread; in vol terms the bid-ask straddles
tens of vol points. The wings of a NIFTY weekly do not contain a fittable vol
level at all, and fitting SVI's linear-in-$|k|$ wings to them is fitting the tick
grid. \emph{Note the nuance from Section~\ref{sec:imputed}: at end of day, in the
liquid band, the half-tick contributes under 0.17 vol points. The tick problem is
real in the far wings and on live quotes; it is not what contaminates the liquid
range. The dominant contaminant there is imputation.}

\item \textbf{The front slice is 0--2 DTE most of the time.} Total variance
$w = \sigma^2 T$ is of order $10^{-4}$ there. The SVI loss surface becomes nearly
flat in $\rho$ versus $m$, and calibration pins $\rho$ at its bound --- an
identifiability failure that appears even on clean synthetic chains. SVI's
asymptotically linear wings are a large-$T$ argument; at one day to expiry the
true smile is jump-dominated and V-shaped near the money.

\item \textbf{There is no term structure} (Section~\ref{sec:noterm}). BANKNIFTY
lost its weeklies in November 2024 and has roughly three usable maturities;
NIFTY has 97\% of turnover inside seven days. SSVI calendar work has nothing to
calibrate against.
\end{enumerate}

\subsubsection{What to do instead}

Stated as a recommendation, not a finding: fit SVI \textbf{only on the band where
the market genuinely quotes two-sided size} --- roughly 10-delta to 10-delta on
NIFTY weeklies, tighter on expiry day --- in total variance, on OTM-only quotes,
with a width-based weight, and \textbf{extrapolate the wings rather than fit
them}. Publish an error bar per strike from the quote width, and refuse to call
anything rich or cheap unless the residual exceeds it.

Anyone claiming a wing dislocation on a NIFTY weekly put quoting 0.05 by 0.10 is
reporting a tick, not a trade.

\subsection{The residual: is your codebase a trade or infrastructure?}
\label{sec:residual}

Round one's two agents contradicted each other directly on this, and it is the
question that decides what your existing work is worth. One said the surviving
use of SVI is that ``the fit residual itself is the signal.'' The other said the
wings are measurement artefact. Round two adjudicated against real data.

\begin{keyfact}{The adjudication}
\textbf{Both were partly right, and both stated the wrong mechanism.} The
artefact is confirmed --- but it comes from \emph{imputed settlement prices}
(79.4\% of untraded settles are off the tick grid), not from tick quantization
(under 0.17 vol points, less than 0.5\% of the signal). And the residual claim is
refuted as stated, for a different reason: a residual measured against a surface
fitted to a chain that is half imputed is not a market signal. At end of day the
signal and the contaminant are not separable.

\textbf{Your SVI codebase is infrastructure, not a trade --- as currently fed.}
It becomes a candidate trade only under filters neither round-one agent specified.
\end{keyfact}

\subsubsection{The quote-quality filters, all mandatory, all falsifiable}

\begin{enumerate}
\item \textbf{Drop every strike with zero traded volume in the bar.} This alone
removes about 52\% of the end-of-day chain and all of the imputation signature.
\item \textbf{Never fit to settlement or close price.} Both are contaminated for
untraded strikes and close is additionally stale. Fit to intraday two-sided
quotes --- which is exactly what your incoming minute data provides and what the
end-of-day bhavcopy structurally cannot.
\item \textbf{Require a two-sided quote with bid $>$ 0} and a relative spread
below a threshold set from the data, not from taste.
\item \textbf{Screen on the tick grid as a permanent integrity check.} Any price
off the five-paise grid in an intraday feed is imputed or erroneous. Free,
decisive, and it should run always.
\item \textbf{Restrict to $|\log\text{-moneyness}| \le 0.10$ on the front weekly},
widened only where measured volume supports it.
\item \textbf{Front one or two expiries only.} Per SEBI, 97\% of turnover is
within seven days and 1\% beyond ten. Any calendar or term-structure residual
strategy is untradeable on this venue --- a hard finding that kills a class of
ideas cheaply.
\end{enumerate}

\subsubsection{The pre-registered falsification test}

Write this down before you look at the data.

\begin{keyfact}{The test}
The residual is signal \emph{if and only if} the sign of today's residual at
strike $k$ predicts the sign of tomorrow's change in that strike's mid-IV, out of
sample, \textbf{after subtracting the round-trip half-spread at that strike}, on
strikes that passed filters 1--6.
\end{keyfact}

Three specific ways it dies. Check each; do not wait to discover them.

\begin{itemize}
\item \textbf{It reverses.} The residual predicts mean reversion \emph{of the
fit}, not of the market --- it forecasts its own next recalibration. Test by
refitting with a different SVI parameterisation (raw versus SVI-JW) and checking
the residual survives. If it is parameterisation-dependent, it is an artefact of
your own model.
\item \textbf{It is real but uncapturable.} The residual is smaller than the
half-spread at every strike where it is statistically distinguishable from zero.
\textbf{This is the modal outcome}, and it is precisely what happened to the
Vilkov 0DTE results (Section~\ref{sec:vilkov}) --- an entire result set inverted
on a cost term two orders of magnitude smaller than a NIFTY wing spread. Design
to detect this early.
\item \textbf{It is a liquidity proxy.} The residual is largest exactly where
volume is lowest, meaning you have rediscovered the illiquidity premium. Test by
regressing residual magnitude on contemporaneous volume and spread; if those two
absorb it, there is nothing left.
\end{itemize}

\textbf{What changes on the answer.} If it survives all three, your SVI work is
the core of a trade and everything else is plumbing. If it does not, the codebase
is still worth having --- as the surface that prices and risk-manages positions
you take for other reasons --- but it is not the alpha, and you should stop
treating it as such. Either way the question is \textbf{intraday-only}: it cannot
be resolved on end-of-day data, and your incoming minute-wise NIFTY feed is the
correct and possibly the only instrument. That is the good news. The decisive
experiment is one you can actually run, in weeks.

\subsection{Dealer positioning: refuted as stated, and the better door it opens}
\label{sec:gex}

Round one asserted that dealer-positioning proxies are computable from open
interest plus a Black-Scholes gamma map, and presented this as the one
implementable framework in that report. Round two made the identifying assumption
explicit and then tested it.

\subsubsection{The assumption, stated}

The standard GEX construction computes $\text{Spot}^2 \times \Gamma \times
\text{OI}$ per strike, with \textbf{call OI positive and put OI negative}. That
sign convention is not a fact about the market. It is an assumption that
\emph{dealers are long every call and short every put} --- equivalently, that
customers buy puts for protection and sell calls in overwriting programmes.

Open interest is \emph{unsigned}. Every contract has a long and a short, and the
exchange publishes only the total. Nothing in an OI file distinguishes a dealer's
long from a customer's long. In US index options the assumption is a defensible
stylised fact grounded in the protection-buying bid. It is empirical, not
structural, and it carries all the weight.

Studies that actually identify dealer sign --- Ni, Pearson, Poteshman and White;
Baltussen et al. --- use \textbf{signed order-flow data} (CBOE/ISE Open-Close,
which flags buy-to-open versus sell-to-open), not open interest. That is the
honest state of the art: sign comes from flow data.

\subsubsection{On NIFTY the assumption is contradicted, and the primary file is free}

NSE publishes signed participant-category option positions \emph{daily}, at
\texttt{nsearchives.nseindia.com/content/nsccl/fao\_participant\_oi\_DDMMYYYY.csv},
with columns for Option Index Call Long, Put Long, Call Short and Put Short by
client type. \textbf{There is no US equivalent at this granularity, freely
published, daily.} Neither round-one agent mentioned it.

Three consecutive days, net positions in contracts:

\begin{center}
\begin{tabular}{@{}llrr@{}}
\toprule
Date & Participant & Net calls & Net puts \\
\midrule
25 Aug 2026 & Client & $+94{,}693$  & $\mathbf{-604{,}022}$ \\
            & FII    & $-198{,}430$ & $+549{,}436$ \\
            & Pro    & $+99{,}801$  & $+24{,}124$ \\
26 Aug 2026 & Client & $+191{,}071$ & $\mathbf{-663{,}004}$ \\
            & FII    & $-235{,}885$ & $+555{,}751$ \\
27 Aug 2026 & Client & $+289{,}153$ & $\mathbf{-760{,}983}$ \\
            & FII    & $-265{,}856$ & $+626{,}360$ \\
\bottomrule
\end{tabular}
\end{center}

\textbf{Read the put column.} The standard convention requires customers to be
net \emph{long} puts. On NIFTY, clients are net \textbf{short} puts by 600{,}000
to 760{,}000 contracts, consistently and growing, while FIIs are net long puts and
net short calls. The retail put leg has \emph{the opposite sign} to the US
stylised fact. Applying the standard GEX sign convention to NIFTY does not merely
lack evidence --- it inverts the sign of the dominant leg.

\subsubsection{What the file does and does not give you}

It is aggregated across all strikes and expiries, across all index options rather
than NIFTY alone, daily rather than intraday, and ``Pro'' is proprietary trading
members --- a proxy for market makers, not identical to them. So it does not hand
you dealer gamma.

What it does is \textbf{falsify the sign assumption and give you a daily aggregate
constraint}. Fit whatever strike-level sign model you like, then check that its
aggregate reproduces this file. That test is free, daily, and nobody in either
round proposed it.

\srcnote{Has anyone shown out-of-sample value for GEX on NIFTY? No --- a genuine
``nobody has published on this''. The search returned only commercial dashboards,
all of which auto-reject on the paid-funnel and no-cost-model criteria, and none
of which disclose their sign assumption. The academic literature is entirely S\&P
500, and the honest null there is a pre-registered SPY backtest over 1{,}972 days
finding that GEX and its relatives add little after controlling for VIX and ATM
implied vol. Benchmark anything you build against those two controls before
believing it.}

\newpage
\section{The measurement layer}
\label{sec:measure}

Three techniques, all closed-form or near it, all implementable in weeks, all
absent from round one. This is the fastest-payoff section in the document and the
one that makes every other result you produce legible.

\subsection{How often should you hedge?}

\subsubsection{The literature, in one notation}

Zakamouline's free 2005 working paper is the single most useful document here
because it contains all five strategies in one consistent notation \emph{plus}
the horse-race between them:

\begin{itemize}
\item \textbf{Leland}: inflate volatility, $\sigma_m^2 = \sigma^2(1 + \sqrt{8/\pi}\,\lambda/(\sigma\sqrt{\Delta t}))$, then hedge at the resulting delta. Note its own failure: as $\Delta t \to 0$ the modified volatility is unbounded and the hedge degenerates to holding one share.
\item \textbf{Whalley--Wilmott}: no-trade band of half-width
$H = \left(\tfrac{3}{2} e^{-r(T-t)} \lambda S \Gamma^2 / \gamma\right)^{1/3}$
around the Black-Scholes delta. \textbf{The cube root is the signature of the
asymptotic solution.} No volatility adjustment.
\item \textbf{Barles--Soner}: band plus a gamma-dependent vol adjustment.
\item \textbf{Hodges--Neuberger} and \textbf{Davis--Panas--Zariphopoulou}: the
exponential-utility formulation, solved numerically. Davis et al.'s useful result
is that the option price is approximately invariant to the \emph{form} of the
utility function --- mainly only the level of risk aversion matters.
\item \textbf{Zakamouline's approximation}: adds a term $H_0 = \lambda/(\gamma S \sigma^2 (T-t))$ that the two asymptotic strategies lack entirely.
\end{itemize}

His Monte Carlo (100{,}000 paths, 250 rebalance opportunities) ranks them:
approximation $>$ Whalley--Wilmott $>$ Barles--Soner $\approx$ delta-tolerance
$>$ Leland $>$ plain Black-Scholes.

\begin{keyfact}{The line that decides what you should build}
Zakamouline: ``When $\lambda < 0.001$, our approximation strategy and both the
asymptotic strategies produce practically identical results.'' NIFTY-futures
$\lambda$ is about \textbf{1.3\,bp $= 0.00013$} --- an order of magnitude below
that threshold. \textbf{For hedging NIFTY delta with NIFTY futures, Zakamouline
buys you nothing over plain Whalley--Wilmott.} Build WW and spend the saved weeks
on realised volatility. The calculus flips only if you hedge with options.
\end{keyfact}

\subsubsection{The other half: the cost of \emph{not} hedging}

Derman \& Kamal (Goldman Sachs, 1999) give the risk leg in closed form. For $N$
rehedges, the standard deviation of replication P\&L is approximately
$0.886\,\sigma\,\text{Vega}/\sqrt{N}$; expressed in vol points, the one-standard-deviation
replication error is about $\sigma/\sqrt{N}$. Their measured numbers: daily
hedging of a one-month ATM option gives a P\&L standard deviation of
\textbf{16.3\% of premium}; quadrupling the frequency halves it to 8.1\%.

\textbf{Mean P\&L is zero.} Discrete hedging adds noise, not bias. That is the
correct null and most people get it wrong.

\subsubsection{The answer, computed for your instrument}

Sepp's contribution is the objective function the other two are inputs to: expected
P\&L is independent of hedging frequency, P\&L standard deviation falls as
$1/\sqrt{N}$, but expected transaction costs grow as $\sqrt{N}$ --- not $N$,
because each rebalance is smaller. So \textbf{Sharpe is humped in $N$} and there
is an interior optimum.

Worked for a weekly NIFTY straddle: five sessions, 1{,}875 minutes,
$\sigma = 13\%$, $\lambda = 1.3\,$bp.

\begin{center}
\begin{tabular}{@{}lrrr@{}}
\toprule
Rebalance & $N$ & Derman noise (vol pts) & Leland cost (vol pts) \\
\midrule
Daily   & 5     & 5.80 & 0.16 \\
Hourly  & 31    & 2.30 & 0.40 \\
15-min  & 125   & 1.16 & 0.80 \\
5-min   & 375   & 0.67 & 1.35 \\
1-min   & 1{,}875 & 0.30 & 2.90 \\
\bottomrule
\end{tabular}
\end{center}

\textbf{The trade-off bottoms out around 5 to 15 minutes and is flat between
them.} Daily hedging is indefensible --- it throws away 5.8 vol points of noise to
save 1 vol point of cost. Minute hedging is indefensible in the other direction.
Redo this with your own $\lambda$ and $\sigma$; note the two columns are not
strictly additive (one is a standard deviation, one is a mean drag), so the table
approximates Sepp's Sharpe rather than replacing it.

\subsubsection{Three things about the Indian cost structure nobody states}

\begin{enumerate}
\item \textbf{The spread is irrelevant; the tax is the cost.} STT is roughly 77\%
of $\lambda$. The half-spread on NIFTY futures is 0.01\,bp --- a hundred times
smaller. Every US paper on hedging costs is about spread and impact. In NIFTY it
is about a statutory levy.
\item \textbf{STT is sell-side only, which violates the symmetric-cost assumption
of all four band models.} Leland, Whalley--Wilmott, Barles--Soner and Zakamouline
all derive a \emph{symmetric} band around delta. With a one-sided cost the true
no-trade region is asymmetric --- the sell boundary should sit further from delta
than the buy boundary. \textbf{Nobody has published this for the Indian market.}
It is a clean, small, publishable extension and it takes days.
\item \textbf{The cube root means the cost input barely matters --- and that is
the point.} Band $\propto \lambda^{1/3}$: to halve the band you must cut costs
eightfold; doubling $\lambda$ widens it only 26\%. Total cost over the option's
life scales as $\lambda^{2/3}$, so a doubling of costs costs you 59\% more, not
100\% --- \emph{provided you re-widen the band}. Keep the band fixed and you eat
the full doubling. Given the 1 April 2026 STT rise, this is live: $\lambda$
roughly doubled and your band should widen about 29\%.
\end{enumerate}

\subsubsection{Your Black-Scholes delta is systematically wrong on NIFTY}

Hull \& White's minimum-variance delta is the correction:
$\delta_{MV} = \delta_{BS} + \nu_{BS}\,\partial\mathbb{E}[\sigma_{\text{imp}}]/\partial S$,
estimated empirically as
$\delta_{MV} = \delta_{BS} + (\nu_{BS}/(S\sqrt{T}))(a + b\,\delta_{BS} + c\,\delta_{BS}^2)$,
fitted by OLS separately for puts and calls on a rolling window.

On the S\&P 500 over 1.3 million quotes, out of sample, this delivered variance
reduction of \textbf{26\% for calls and 23\% for puts}, rising to \textbf{42\% for
the highest-strike OTM calls}. The correction was negative in every
moneyness/maturity bucket --- \emph{under-hedge calls, over-hedge puts} --- and it
beat a daily-recalibrated SABR using 312 parameters against SABR's 100{,}000.

NIFTY has a steep, persistent negative skew and a strongly negative NIFTY/India
VIX correlation, so the correction is at least as large as SPX's. Two things
follow:

\begin{itemize}
\item \textbf{Nobody has published $a$, $b$, $c$ for NIFTY.} Novel, cheap, high value.
\item Hull \& White explicitly confined their test to maturities beyond 13 days,
noting that including very short maturities worsens results because of large
gammas near the money. \textbf{That exclusion removes almost the entire NIFTY
weekly complex}, which is where all the volume is. So you cannot import their
coefficients; you must estimate your own, conditioned on days to expiry --- and
because NSE now runs a weekly Tuesday cycle, the coefficient will almost certainly
be a function of position in that cycle, a structure the SPX literature does not have.
\end{itemize}

\subsubsection{Expiry day: the formulas fail, and they fail loudly}

As $\tau \to 0$, ATM gamma diverges, so the Whalley--Wilmott band widens as
$\tau^{-1/3}$ --- past the entire feasible delta range. \textbf{The model is not
asking for infinite trading. It is telling you the position is unhedgeable and
you should stop.} Away from the money, gamma goes to zero but Zakamouline's $H_0$
diverges. Leland degenerates regardless.

The practitioner override is not a better band. It is to stop treating expiry day
as a hedging problem and treat it as a \emph{sizing} problem:

\begin{enumerate}
\item \textbf{Cap dollar gamma ex ante} --- decide the maximum you will carry into
the final session at position-entry time, when you can still choose. No
rebalancing rule rescues a position that is too big on Tuesday afternoon.
\item \textbf{Your futures hedge has structural basis risk on expiry day by
construction.} Final settlement is the index close (and since August 2026, an
auction print), not the futures. A delta hedge in futures cannot replicate it.
This is precisely what SEBI's +2\% expiry-day ELM is compensating for --- the
circular's stated rationale is the basis risk of the expiring contract.
\item Roll or close before the final session if gamma is material, or hedge gamma
with a same-expiry option rather than delta with futures.
\item After the close you hold an unhedgeable position through the settlement
window. Size for that, not around it.
\end{enumerate}

\subsubsection{What minute bars can and cannot measure}

NSE's continuous session is 375 minutes. Time-based rules --- rebalance every
$\Delta t \ge 1$ minute --- are \textbf{exactly} measurable on minute bars.

\textbf{Band rules are not, and the bias is one-directional.} A delta band is
touched \emph{inside} a minute. Check only bar closes and you under-count
crossings, hedge late, and overstate the strategy. Use bar high/low to detect the
touch and assume a fill at the band level and you have assumed a resting limit
order filled at the touch --- zero adverse selection. That is the standard silent
optimism in band backtests.

\begin{keyfact}{The rule of thumb}
NIFTY per-minute return standard deviation is about $13\%/\sqrt{250 \times 375}
\approx 4.2\,$bp. If the band's half-width in spot terms, $h/(\Gamma S)$, is less
than roughly twice that (about 8\,bp), a typical minute contains multiple
crossings and minute bars cannot resolve the strategy. Narrow bands are exactly
what low $\lambda$ and high $\Gamma$ produce --- i.e.\ exactly NIFTY weeklies near
expiry. \textbf{Calibrate time-based rules on minute bars honestly; treat
band-rule minute-bar results as an upper bound and validate the winner on ticks.}
\end{keyfact}

\subsection{Realised volatility from minute data}
\label{sec:rv}

Your single largest new asset, and no round-one agent said how it should be measured.

\subsubsection{The decision, already made on evidence}

Liu, Patton \& Sheppard tested 420 estimators across 31 assets, five asset
classes and eleven years of tick data, with proper multiple-testing correction.
\textbf{Nothing significantly beats 5-minute realised variance} except for one
asset. Four operational findings matter more than the headline:

\begin{itemize}
\item Subsampling helps at 5 and 15 minutes; it \emph{hurts} at 1 minute and below.
\item Transaction prices beat mid-quotes at low frequencies; quotes win only at
one tick or one second.
\item Tick-time sampling beats calendar-time at 5 and 15 minutes.
\item \textbf{The line that matters most for you:} ``more sophisticated realized
measures generally perform significantly worse in non-US markets than in US
markets.'' The realized-kernel and two-scale machinery is tuned to NYSE/CME
microstructure. It has no demonstrated advantage on NSE and some demonstrated
disadvantage outside the US.
\end{itemize}

\subsubsection{The diagnostic you must run first --- 30 minutes of work}

Andersen, Bollerslev, Diebold \& Labys's volatility signature plot: average
realised volatility against sampling interval. Two shapes with opposite meanings.

\begin{itemize}
\item \textbf{Liquid asset:} RV highest at the smallest interval, falling and
stabilising. Cause is \emph{negative} serial correlation from bid-ask bounce.
Naive 1-minute RV is biased \textbf{up}.
\item \textbf{Illiquid asset:} RV \emph{lowest} at small intervals, rising.
Cause is \emph{positive} serial correlation from stale trading. Naive 1-minute RV
is biased \textbf{down}.
\end{itemize}

\begin{keyfact}{This decides your instrument choice, and the answer is not obvious}
The \textbf{NIFTY 50 index} is a weighted average of fifty stocks that do not
print synchronously. At high frequency it shows the \emph{illiquid} signature ---
positive autocorrelation from staleness --- so 1-minute RV on the index is biased
\textbf{downward}, opposite to the standard bid-ask-bounce story. The
\textbf{front-month NIFTY future} is a single extremely liquid instrument and will
show the liquid signature. \textbf{Compute RV on the future, not the index.} Draw
both signature plots and look at them before writing another line of code. Note
also that the originators of the tool recommended 20 minutes for their own data,
not 5.
\end{keyfact}

\subsubsection{The overnight gap, which is decisive for weeklies}

Hansen \& Lunde give the MSE-optimal linear combination
$RV_{\text{day}} = \omega_1 r^2_{\text{overnight}} + \omega_2 RV_{\text{open-close}}$,
and rule out the two things everyone does: \textbf{do not} simply add the squared
overnight return, and \textbf{do not} scale open-to-close RV by a constant. Both
are dominated. Their DJIA finding puts about 80\% of daily volatility in the
active session, with weights tilted heavily toward RV because a single squared
overnight return is a far noisier estimator of its period than RV is of the session.

\textbf{The DJIA number does not transfer.} NIFTY's overnight window is 15:30 to
09:15 --- \textbf{17.75 hours}, spanning the entire US session and GIFT Nifty's
overnight trading, plus a pre-open auction. \textbf{Nobody has published the
active/inactive variance ratio for NIFTY.} You must estimate it, and it is
genuinely important: a weekly option prices seven calendar days of variance while
you can only measure five sessions of 375 minutes. Get this wrong and the entire
implied-versus-realised comparison is biased.

\subsubsection{NSE intraday seasonality is not the textbook U}

The only free study measuring NIFTY intraday volatility on real minute data
(1-minute futures, 2011--2018) reports hourly-scaled volatility:

\begin{center}
\small
\begin{tabular}{@{}lrrrrrrr@{}}
\toprule
& 09:15--10 & 10--11 & 11--12 & 12--13 & 13--14 & 14--15 & 15--15:30 \\
\midrule
Vol (\%) & \textbf{0.350} & 0.238 & \textbf{0.225} & 0.246 & 0.285 & \textbf{0.327} & 0.280 \\
\bottomrule
\end{tabular}
\end{center}

Open-to-trough ratio is 1.56, a variance ratio of 2.4. But \textbf{the last half
hour is lower than 14:00--15:00}. NSE's shape is a \textbf{J with a clipped right
arm}, not a U --- the afternoon peaks at 14:00--15:00 and volatility \emph{falls}
into the close. Import a US-calibrated U-shaped diurnal factor and you will
over-deseasonalise exactly the window that sets the daily settlement price and the
expiry settlement value.

\srcnote{Stated plainly: this source does \emph{not} pass the filter. No cost
model, not a practitioner, a student working paper in a non-refereed archive, and
its methodology text describes converting per-minute to hourly volatility by
multiplying by 60 --- which is wrong (it must be $\sqrt{60}$), and the reported
levels are only consistent with correct scaling, so the description and the
arithmetic disagree. Use the \emph{shape} as a prior to check your own estimate
against, and re-estimate everything on your own data. An honest weak source
flagged as weak beats a confident US import.}

Three NSE-specific facts verified from primary exchange documents that no paper covers:

\begin{enumerate}
\item A \textbf{pre-open auction runs 09:00--09:08}, so the 09:15 open already
absorbs part of the overnight gap. The first minute bar is a post-auction bar, not
a gap bar. Do not treat 09:15 as the gap.
\item The \textbf{15:00--15:30 window is the daily settlement window}. Volume and
volatility there are mechanically driven by settlement, not information.
\item On expiry day that same window determined final settlement for every
expiring index option --- and \textbf{moved from Thursday to Tuesday} for
contracts expiring on or after 1 September 2025, then was \textbf{superseded by
the Closing Auction Session} on 3 August 2026. Estimate the periodicity factor
separately per regime. That is at least four regimes, arguably six.
\end{enumerate}

\subsubsection{The recommended estimator, and the first project specified}
\label{sec:rvproject}

\begin{keyfact}{Recommendation}
\textbf{Subsampled 5-minute realised variance on the front-month NIFTY future,
computed in tick time, on transaction prices, open-to-close, plus a
Hansen--Lunde-weighted overnight term --- with the sampling interval confirmed,
not assumed, by your own volatility signature plot.}
\end{keyfact}

The project, step by step:

\begin{enumerate}
\item \textbf{RV pipeline.} Front-month future, minute bars, aggregated to 5-minute
in tick time, subsampled across the offset grids. Compute bipower variation and
MedRV on the same grid for the jump decomposition.
\item \textbf{Deseasonalise} with a jump-robust periodicity estimator, fitted
\emph{separately per regime}. Use the deseasonalised series for spot-vol and jump
detection; use raw RV for the daily aggregate.
\item \textbf{Overnight.} Estimate NIFTY's active/inactive split yourself, then
combine with Hansen--Lunde MSE-optimal weights. Report the estimated split --- it
is a publishable number in its own right.
\item \textbf{Match the clocks.} Remaining implied variance at minute $t$ is
$\sigma_{\text{imp}}(t)^2 \tau(t)$ with $\tau$ in \emph{calendar} years to
settlement, because the option prices calendar time including weekends. Remaining
realised variance is the sum over surviving minutes \emph{plus every overnight to
expiry}. Mismatch these two clocks and the whole study is noise.
\item \textbf{Weight by dollar gamma, not equally.} This is the step that makes it
a trade rather than a statistic. Delta-hedged short-option P\&L is
$\tfrac{1}{2}\int \Gamma_s S_s^2 (\sigma_{\text{imp}}^2 - \sigma_s^2)\,ds$ --- a
\emph{dollar-gamma-weighted} average of the variance gap. A variance risk premium
that lives in the low-gamma part of the option's life is not tradeable. Compute
both the plain and the gamma-weighted VRP and report the difference; if they
disagree, the plain one is a mirage.
\item \textbf{Decompose and stress.} Split RV into continuous and jump components,
locate the jump minutes, and re-run excluding jump days. Compare implied against
RV \emph{including} jumps --- that is what the hedged P\&L actually eats --- but
report the jump-free version alongside. The gap between the two tells you whether
you are harvesting a variance premium or writing crash insurance.
\end{enumerate}

\srcnote{Toolchain note, and an honest negative. The complete, peer-reviewed, free
implementation is the R package \texttt{highfrequency} --- realised kernels,
two-scale, pre-averaging, bipower, semivariance, Barndorff-Nielsen--Shephard and
Lee--Mykland jump tests, and the intraday periodicity estimators you specifically
need. \textbf{There is no Python equivalent of comparable completeness or peer
review.} The closest is Kevin Sheppard's MATLAB MFE Toolbox, which is the code the
Liu--Patton--Sheppard study actually used. If you are a Python shop you will be
writing RV estimators yourself or calling R. Better to know that now.}

\subsection{Sizing and risk of ruin}
\label{sec:sizing}

\subsubsection{Why full Kelly is wrong for options}

\begin{enumerate}
\item \textbf{Kelly is optimal for the distribution you feed it, and your backtest
does not contain your loss distribution.} A short-gamma book's P\&L is a series of
small positive numbers punctuated by a large negative one. With zero or one tail
events in sample, $\mu$ is biased up and $\sigma^2$ biased down, so
$f^{*} = \mu/\sigma^2$ is biased up \emph{multiplicatively in both terms}. This is
the dominant error, not a small one.
\item \textbf{The moment correction points the same way.} Sinclair \& Brooks derive
closed-form skewness and kurtosis of European option \emph{returns} and show the
higher moments shift the Kelly fraction materially --- and for short options they
shift it \emph{down}. \srcnote{Flagged: SSRN blocks programmatic access, so the
abstract and affiliation were read but not the 13-page body, and no formula from
it is reproduced here. It is free and short; download it yourself.}
\item \textbf{Kelly's own properties are bad.} Full Kelly's expected maximum
drawdown is unbounded in the limit, and Thorp's own ``when will Kelly dominate''
analysis shows the horizon is long enough that a solo trader will be stopped out
--- by his own nerve or his broker --- long before the asymptotics arrive.
\end{enumerate}

\subsubsection{The fix is not ``use half-Kelly''}

Busseti, Ryu \& Boyd's risk-constrained Kelly makes the drawdown tolerance the
\emph{primitive} rather than a fudge factor. Maximise $\mathbb{E}\log(r^\top b)$
subject to the usual constraints plus $\mathbb{E}(r^\top b)^{-\lambda} \le 1$, with
$\lambda = \log\beta / \log\alpha$ where $\alpha$ is the drawdown level and $\beta$
the tolerated probability.

It is a \textbf{convex} problem, its entire specification is about six lines of
CVXPY, and the code is published. Its properties:

\begin{itemize}
\item The constraint bounds the \emph{entire} CDF of minimum wealth:
$P(W_{\min} < \alpha) \le \alpha^{\lambda}$ for all $\alpha$, not just the one you
specified. Their example: $\alpha = 0.7$, $\beta = 0.1$ gives $\lambda = 6.46$, so
the probability of ever losing half your capital is bounded by $0.5^{6.46}
\approx 1.1\%$.
\item $\lambda = 0$ recovers unconstrained Kelly exactly; $\lambda$ is a
risk-aversion dial.
\item It \emph{dominates} fractional Kelly in simulation --- higher growth at the
same drawdown risk, or lower risk at the same growth.
\end{itemize}

\textbf{The caveat that applies to all of this:} the method takes the return
distribution as given. Feed it the empirical distribution of your backtested
short-gamma trades and it will hand you a confidently oversized position, because
the empirical distribution does not contain the loss that has not happened yet.

\subsubsection{Margin expansion is the constraint that binds first}

This has nothing to do with the maths and will kill you before the maths does.
SPAN risk parameters are revised \textbf{six times a day} --- 11:00, 12:30, 14:00,
15:30, 17:00, end of day and beginning of day. \textbf{Your margin requirement can
be raised four times during the session you are losing money in.} That is the
documented, scheduled, public mechanism of margin expansion under stress. Add the
+2\% expiry-day ELM, which raises margin on the day your gamma is largest, and the
removal of calendar-spread benefit on expiry day, and the \rs 15 lakh minimum
contract size, which quantises any Kelly fraction below a certain account size.

\begin{keyfact}{The one-line rule}
Size every short-gamma position so that a single overnight gap of the worst
magnitude you can \emph{imagine} --- not the worst in your sample --- combined
with a SPAN scan-range widening and the expiry-day ELM add-on, still leaves you
funded; \emph{then} take the risk-constrained-Kelly fraction for your stated
drawdown tolerance and halve it, because your edge is estimated with far more
error than your loss.
\end{keyfact}

The halving is not superstition. It is a crude stand-in for the estimation
variance of $f^{*}$, and you should replace it with the actual confidence interval
once you have read Sinclair's paper on confidence intervals for the Kelly criterion.

\subsection{P\&L attribution: how you find out whether any of this works}

\subsubsection{The free corpus, and a name correction}

\begin{keyfact}{Timothy Klassen, not Kevin}
PhD particle physics, then equity derivatives quant under Emanuel Derman at
Goldman, then \textbf{designed the 2003 ``new VIX'' for Cboe}, then built Getco's
options analytics infrastructure and quant team, then founded Vola Dynamics in
2016 with Jim Gatheral as an advisor. He is the single richest free source in this
entire corpus for spot-vol dynamics and P\&L attribution, and round one missed him
by searching for ``PnL attribution'' as an academic topic rather than as a desk
engineering problem.
\end{keyfact}

His \emph{Spot-vol Dynamics and Deltas for SPX Options} (13 pages, free) is
directly executable on your minute data. It gives the estimation recipe verbatim:
``the (realized) vol sensitivities can be determined reasonably well using, say,
one day's worth of 1-minute or 5-minute snapshot data using a linear regression,
with some outlier treatment for robustness.'' That is a one-day project for you.

Concrete magnitudes of the error you are currently making: ATM delta adjustment
about $-0.06$; put-wing peak about $-0.08$; and in the call wing the correct delta
is essentially \emph{zero} while the naive Black-Scholes delta is still 0.03 --- a
100\% relative error on tradable OTM calls.

And the explicit falsification, which is the part to internalise:

\begin{quote}\small
``Sticky-strike scenarios are still commonly used by risk departments. In
equities, at least, neither [sticky-strike nor sticky-delta] has happened for
20y+. Sticky-delta implies SSR$=$0 for all terms. Never happens. Sticky-strike
implies SSR$=$1. Even when vols are sticky-by-strike around ATM, they never are in
the wings.''
\end{quote}

Typical skew-stickiness ratio is about 1.3, rising toward 2.0 on large down days,
with \textbf{KOSPI often above 2} --- the one named non-US index, and plausibly a
better prior for NIFTY than SPX is. Note also that SSR now has a term structure
and is horizon-dependent: estimate it from 1-minute NIFTY data and you will get a
different number than from close-to-close.

The free video series is the curriculum, in this order: \emph{PnL Attribution \#1:
The Greeks} $\to$ \emph{\#2: Scenarios} (full revaluation versus greek attribution
--- i.e.\ where the unexplained residual comes from) $\to$ \emph{Thetas and Time
Conventions} $\to$ \emph{Sticky-by-Strike or Sticky-by-Delta? Neither!} $\to$
\emph{If sticky strike is not the answer, what is?} Then read the Spot-vol
Dynamics PDF with a notebook open.

\srcnote{Caveat, stated because it is load-bearing: this is marketing collateral
for a paid C++ library, and the self-serving claim is explicit --- ``good vol
curves that can actually match the market implied vols for liquid names (none in
the public domain do)''. Since you built SVI, treat that as a hypothesis to test.
His slides also argue SVI has insufficient shape flexibility and that SPX needs
16--18 parameters, with an empirical scaling law of roughly
$n_{\text{params}} \approx (n_{\text{options}}/5)^{1/3}$. Note what that law
implies: \textbf{NIFTY chains are far smaller than SPX's, so NIFTY needs
materially fewer parameters --- possibly making SVI adequate where it is not for
SPX.} Test it; do not assume either way.}

\subsubsection{The regulator already answered ``how much unexplained P\&L is too much''}

BCBS d457's P\&L Attribution test is the only quantitative, published,
universally applied standard for whether a risk model explains its P\&L. Compare
risk-model P\&L against hypothetical P\&L (yesterday's positions revalued with
today's market data, no intraday trading, no new deals) on a rolling 250-day
window, using two statistics: \textbf{Spearman rank correlation} and the
\textbf{Kolmogorov--Smirnov} statistic.

\begin{center}
\begin{tabular}{@{}lll@{}}
\toprule
Zone & Correlation & KS statistic \\
\midrule
Green & $> 0.80$ & and $< 0.09$ \\
Amber & \multicolumn{2}{l}{between --- carries a capital surcharge} \\
Red   & $< 0.70$ & or $> 0.12$ \\
\bottomrule
\end{tabular}
\end{center}

The footnote that answers your theta question directly: P\&L due to the passage of
time must be included and \emph{treated consistently} across both measures. The
regulator's answer to calendar-versus-business time is \textbf{pick one and be
consistent} --- inconsistency there is itself a model failure.

The instructive negative: Basel's own vega bucketing for equities is defined along
\emph{one} dimension, option maturity, mapped to 0.5, 1, 3, 5 and 10 years. It
does \textbf{not} bucket equity vega by strike or moneyness at all; skew risk is
pushed into a curvature charge. So the regulator's scheme is deliberately cruder
than your book needs. Use it as a \emph{floor} --- a strike-bucketed vega grid is
a strict improvement --- and take the published maturity-decay correlation kernel
as a defensible free starting point.

\srcnote{The honest limit: the PLA test compares model P\&L to model P\&L, not
model P\&L to realised cash. It answers ``is my risk system consistent with my
pricing system'', which is weaker than ``is my model right about the market''. It
is also distributional over 250 days, so it gives you \emph{no per-day threshold}
for ``today's residual is too big''. For that you must set your own bar, and
nobody has published one.}

\subsubsection{The Indian business-time problem, which is now a moving target}

With a \textbf{Tuesday} expiry, the weekend sits \emph{in the middle} of a weekly
option's life rather than at the start. Under a calendar-day convention your
Monday and Friday theta are the same number; under trading-day they are not, and
the difference is the entire P\&L of a short weekly in a holiday week. India runs
14--16 exchange holidays, unevenly clustered around October and November, so a
Tuesday-expiry weekly spans seven calendar days but between \textbf{three and five
trading days} depending on the week.

\textbf{Nobody has published on business-time theta calibration for NIFTY
weeklies.} With minute data you are positioned to close it: estimate intraday
versus overnight versus weekend variance shares directly from 1-minute returns,
then compare the implied vol you back out under each time convention.

\srcnote{Honest negative worth knowing before you go looking: \textbf{there is no
credible open-source P\&L-explain engine.} The best-known hit is a single Jupyter
notebook. Everything else is a zero-star toy repository, and every production
implementation is commercial and closed --- Vola Dynamics, Murex, Bloomberg MARS.
You will write this yourself, and that is a genuine, defensible piece of work.}

\newpage
\section{Backtests that do not lie, and spending the minute-data advantage}
\label{sec:backtest}

\subsection{The most instructive artefact in the entire corpus}
\label{sec:vilkov}

\begin{keyfact}{A peer-reviewed 0DTE result that died when its cost model was fixed}
Grigory Vilkov (Frankfurt School) is one of the few options academics who ships
his replication code. In August 2026 he posted a correction notice on the
repository, verbatim:

\smallskip
\emph{``A transaction-cost unit-scale error caused the bid-ask half-spread to be
charged at 1/100 of its true size (about 0.022\,bp instead of 2.2\,bp). The code
is fixed in this commit, and the correction reverses the sign of the
net-of-cost conclusions: no strategy or basket retains a positive net Sharpe
ratio. Gross results are also affected.''}
\end{keyfact}

\textbf{One correction to how round one reported this.} The repository's
\texttt{KNOWN-ISSUES} file states the result more carefully --- ``no structure
retains a \emph{materially} positive net Sharpe ratio'' --- and its own table does
not support the stronger sentence in the banner:

\begin{center}
\begin{tabular}{@{}lrrr@{}}
\toprule
Structure & SR at mid & SR net (before) & SR net (after) \\
\midrule
Put ratio spread     & $+1.06$ & $+0.84$ & $\mathbf{-0.61}$ \\
Risk reversal        & $+0.65$ & $+0.44$ & $\mathbf{+0.10}$ \\
Bear put spread      & $+0.48$ & $+0.30$ & $\mathbf{-0.73}$ \\
Strangle / straddle  & $-0.27$ & $-0.51$ & $-0.97$ \\
Iron butterfly / condor & $-0.56$ & $-0.96$ & $-2.67$ \\
\bottomrule
\end{tabular}
\end{center}

Risk reversal survives at $+0.10$. Use the careful wording, not the banner.

\subsubsection{The mechanism, exactly}

The data panel mixed scales: bid-ask spread, mid and time value were \emph{fractions}
of spot, while return columns were \emph{percent} of spot. The cost line subtracted
a fraction from a percent, while the flat fee term was correctly scaled. Entry
turnover for a two-leg strangle went from an apparent 0.24\,bp to a true 24\,bp.

\subsubsection{The part round one missed, and the most transferable lesson here}

From the notice's own scope note:

\begin{quote}\small
``Note that gross columns also move, because the binary target
$y = \mathbb{1}\{\text{pnl\_net} > 0\}$ depends on the cost. Correcting the cost
relabels the training data, so the classifier learns a different --- and weaker
--- signal. Anyone rescaling the cost only at the evaluation step, without also
rebuilding the target, will get a more favorable number than the one reported here.''
\end{quote}

\textbf{If your cost model feeds your labels, fixing the cost model is not an
evaluation-time patch.} You must retrain. This is the subtle version of the error
and it will survive a casual audit.

\begin{keyfact}{What actually changes for you}
A cost bug of exactly two orders of magnitude survived peer review, an SSRN
posting and a public replication package, and was caught by an outside replicator
reproducing a table to four decimals. \textbf{Treat every published 0DTE net
Sharpe as unproven until you have read the cost line and checked its units.} Note
the scale: 2.2\,bp of half-spread was enough to invert every result on SPX. A
NIFTY weekly wing spread is orders of magnitude larger than that.
\end{keyfact}

\subsection{How options backtests lie}

In descending order of how much they catch, in options specifically:

\begin{enumerate}
\item \textbf{Re-price the whole backtest at three cost levels}: mid, mid $\pm$
25\% of half-spread, and full spread crossing. \textbf{If the sign of P\&L
changes, you have a cost-model result, not a strategy.} This is the Vilkov failure
and it is the modal failure.
\item \textbf{Timestamp-shift test.} Move every entry and exit by $\pm$5, $\pm$15,
$\pm$30 minutes. A real edge degrades smoothly; a microstructure artefact or a
look-ahead bug falls off a cliff. \textbf{Daily-bar backtests cannot run this test
at all} --- this is one thing your minute data buys immediately.
\item \textbf{Quote-staleness audit before any P\&L is computed.} Per strike per
minute, compute time since the quote last changed and spread as a percentage of
mid. Recompute excluding any leg stale beyond $k$ minutes at decision time.
\textbf{The delta between the two runs is the size of your lie.}
\item \textbf{Live-versus-backtest reconciliation on paper fills.} Run forward,
log your intended price and the market's next print, and track the distribution of
the gap. After about 200 fills you have an empirical slippage model that replaces
every guess. Nothing else substitutes for this.
\item \textbf{PBO via combinatorially symmetric cross-validation} over your
configuration grid, plus walk-forward with a refit cadence matching production ---
including the fit cost in the out-of-sample window.
\item \textbf{Assignment and settlement replay} as a separate unit test, not part
of the strategy loop. Force expiry-day paths through the settlement rule and
verify cash by hand for a few dates.
\end{enumerate}

\begin{keyfact}{On deflated Sharpe versus PBO --- a blunt assessment}
The deflated Sharpe ratio is academically popular and practically half-useless,
because it requires $N$, the number of trials, and \textbf{nobody honestly reports
their own $N$ --- including the people citing DSR}. \textbf{PBO via CSCV is the
one to implement}: it splits the backtest into combinatorial train/test partitions
and reports the probability that your in-sample-best configuration ranks below
median out of sample. It needs no $N$ and is computable from artefacts you already
have. Its limit: it catches \emph{selection} lies, not \emph{execution} lies. Both
are live in options.
\end{keyfact}

\subsection{Open-source options backtesters: specific verdicts}

\textbf{Trustworthy with work --- QuantConnect LEAN.} A real event-driven engine
with a reality-modelling layer and source-readable option pricing. But read their
own docs first: the default fill model \emph{does not consider size or volume},
and their greeks come from a QuantLib Black-Scholes that admittedly diverges from
vendor implied vols. Out of the box it will fill you at prices no one would give
you. \textbf{It has no meaningful Indian options support}, so for NIFTY you are
writing your own engine regardless --- use LEAN as a design reference for the
reality-model separation, not as your runtime.

\textbf{Honest about itself, wrong resolution --- optopsy.} Better than its
reputation: it models commissions and several slippage schemes, and its README
disclaims exactly the right things. But no margin model, no assignment model,
chain-snapshot oriented rather than intraday, and the project is forked across
three organisations --- a reliable signal of abandonment. Screen structures on
end-of-day chains with it; do not use it for anything at minute resolution.

\textbf{Actively misleading if used for options.} \texttt{vectorbt},
\texttt{backtesting.py} and \texttt{backtrader} with an option chain bolted on.
None are options engines: vectorised frames make look-ahead free and silent, none
carry a contract-lifecycle or assignment concept, and \texttt{backtrader} has no
options support at all despite endless blog posts pretending otherwise.

\textbf{Closed-box Indian simulators} --- AlgoTest, Sensibull, Opstra and the rest
--- have undisclosed fill assumptions and are product funnels. As evidence they
are worth zero. As a \emph{cross-check on your own chain reconstruction} (does my
strike list on a given date match theirs?) they are mildly useful.

\textbf{Right philosophy, wrong asset class --- \texttt{hftbacktest}.} Rust and
Python, reconstructs the book from level-2 and level-3 feeds, models \emph{feed
latency and order latency separately}, and simulates fills using queue position.
No options support --- its value to you is as the reference implementation of what
an honest fill simulator looks like, and as the tool for the crypto leg. It is
also the concrete demonstration of why 1-minute options bars cannot answer
passive-fill questions.

\subsection{Verify what you actually received --- this decides everything below}

\begin{keyfact}{Ask your data provider one question}
Is each 1-minute bar OHLC of the \emph{last traded price}, or does it carry a
\emph{best bid / best ask snapshot}? Most Indian ``1-minute options data'' is
last-traded-price only. If so, then for any strike outside the few liquid ones the
``price'' in a bar can be a trade from minutes ago, arbitrarily far from the live
quote, and \textbf{every fill assumption you build on it is fiction}.

LTP-only minute data is still excellent for realised-variance and timing work and
nearly useless for fill-sensitive work. Do not proceed until you know which you have.
\end{keyfact}

\subsubsection{What is tractable at one minute, and what is not}

\textbf{Tractable:} time-of-day realised-variance seasonality; intraday variance
risk premium; intraday surface dynamics and per-minute refits; discrete
delta-hedging error and gamma-scalping P\&L; entry and exit timing sensitivity;
expiry-day pin behaviour; synthetic-forward versus futures basis; event studies
around the open and the settlement window.

\textbf{Not tractable --- needs level-1 with timestamps at minimum, level-2
properly:} passive fill probability, queue position, adverse selection on resting
orders, effective versus quoted spread, order-book imbalance signals, market
making, and anything holding less than roughly 10--15 minutes.

\begin{keyfact}{The working rule}
At 1-minute resolution, only backtest strategies that \textbf{cross the spread on
entry and exit and hold at least 15 minutes}. Everything else is a guess wearing a
Sharpe ratio.
\end{keyfact}

\subsection{Storage: pick this and stop thinking about it}
\label{sec:storage}

Full NIFTY and Sensex option chains at 1-minute for several years is \textbf{tens
of gigabytes compressed}, not terabytes. That is small, and it settles the question.

\begin{itemize}
\item \textbf{Hive-partitioned Parquet on local disk, with DuckDB as the query
engine.} Partition by date and underlying, sort within file by expiry, strike,
right and timestamp, compress with ZSTD. DuckDB queries Parquet directly and keeps
peak memory bounded. \textbf{Do not use DuckDB's native file as your system of
record} --- its on-disk format is not stable across versions. Parquet is truth.
\item \textbf{ArcticDB} is real, Man Group-built, and its bitemporal versioning is
a genuine argument given that vendor restatements of options data are common. But
it is last-writer-wins with no transaction isolation, and versions cost storage.
\textbf{Verdict: not worth it at your scale.} Adopt it only if you start receiving
daily corrections you must audit.
\item \textbf{kdb, ClickHouse, QuestDB, TimescaleDB: no.} They are right at
tick-scale on a shared team feed. At single-researcher minute scale they are a
second job, which is the thing you said you did not want.
\item \textbf{The piece everyone skips and everyone gets burned by:} a separate
contract-metadata table with \texttt{valid\_from} and \texttt{valid\_to} --- lot
size, tick size, expiry date, strike interval. \textbf{Never join on lot size as a
constant.} See Section~\ref{sec:breaks} for why.
\end{itemize}

\subsection{Edges that only exist at minute resolution}

This is where the advantage is actually spent. Six, roughly in order of value.

\begin{enumerate}
\item \textbf{Time-of-day-adjusted realised variance against remaining implied
variance.} Daily data cannot tell you that 1-minute RV of 12\% annualised is
\emph{high} at 12:30 and \emph{low} at 09:20. Deseasonalise, compare against the
option's implied variance over its \emph{remaining life only}, and you have an
intraday VRP signal. Specified fully in Section~\ref{sec:rvproject}. Highest value.

\item \textbf{The intraday implied-vol diurnal pattern is an execution edge before
it is a strategy.} Both calls and puts show lowest implied vol at the open with
call implied vol drifting to a peak at the close. Even if you never trade it
directionally, it means ``sell premium'' backtests that fill at the close are
systematically flattered and ``buy premium at open'' backtests systematically hurt
--- a bias daily-bar work cannot even detect, let alone remove.

\item \textbf{Timing decomposition of any existing daily strategy.} Take any
published Indian short-premium or calendar result, re-run it at every minute of
the day, and plot P\&L against entry time. In almost every case the majority of
the ``edge'' concentrates in two 15-minute windows --- which tells you it is a
microstructure artefact rather than a risk premium. \textbf{This is a falsification
tool daily data physically cannot provide, and it will kill more of your ideas
than anything else you build.}

\item \textbf{Honest gamma-scalping and hedge-error estimation.} The variance of
short-gamma P\&L depends almost entirely on hedge frequency, and daily-bar
backtests silently assume one hedge per day, thereby overstating short-straddle
Sharpe. With minute data you can compute the actual hedge-error distribution as a
function of rebalance interval and price the trade-off. \textbf{The single most
decision-relevant thing minute data gives an options book.}

\item \textbf{Per-minute surface refits} --- the direct extension of work you have
already done. Track the time constant of smile relaxation after a jump, intraday
ATM-skew mean reversion, and \emph{persistent} no-arbitrage violations. Violations
lasting a minute or two are stale quotes, which gives you a free model-based
staleness detector. Violations persisting longer are either a real
calendar/butterfly trade or a data defect --- and distinguishing them is exactly
what is tradeable and what nobody with daily data can even look at. India's
expiry-day chains, with the deepest short-dated liquidity in the world, are
arguably the best dataset on earth for this.

\item \textbf{Synthetic forward versus futures, per minute.} The put-call-parity
forward and the futures price should track. When they diverge, one of your three
inputs is stale. A cheap continuous data-quality monitor, and occasionally an
actual box-spread signal. Impossible at daily resolution because everything is
synchronised at the close by construction.
\end{enumerate}

\subsection{Indian option-chain traps, specifically}

\begin{itemize}
\item \textbf{Lot sizes are point-in-time facts, not constants} --- see
Section~\ref{sec:breaks}.
\item \textbf{Expiry day changed twice} in recent history, and once was a
cross-exchange \emph{swap}. Any day-of-week feature, any expiry flag and any DTE
calculation must be derived from a dated contract master, never from a weekday rule.
\item \textbf{Illiquid wings quote garbage.} Enforce a hard filter --- drop legs
with spread above a threshold percentage of mid, or with zero volume in the bar
--- \emph{before} the strategy sees the chain, and \textbf{report how many
candidate trades the filter killed}. That number is a result in itself.
\item \textbf{Forward, not spot.} Compute implied vol against the synthetic
forward from put-call parity per minute, never against the index spot. India's
index carry is non-trivial and dividends are lumpy; using spot puts a fake skew in
your surface that your fit will faithfully reproduce and you will then trade.
\item \textbf{Expiry-day settlement mechanics.} STT on exercised in-the-money
options is charged on intrinsic value --- a large, asymmetric, frequently omitted
cost that specifically penalises held-to-expiry positions. Section~\ref{sec:stt}
shows the direction of that incentive has now \emph{reversed} relative to the
folklore.
\end{itemize}

\newpage
\section{India: the rules, the arithmetic, and the margin}
\label{sec:india}

Round one computed the crypto tax arithmetic in detail and never stated the NSE
F\&O treatment at all, then implicitly compared venues across the two. Round two
computed both halves from the statute and the exchange's own files. Everything in
this section is derived from primary documents; the unverified residue is listed
at the end.

\subsection{What it costs to trade, per round trip}
\label{sec:coststack}

\begin{keyfact}{One correction that changes the premise}
The Finance Act 2026 raised derivative STT again with effect from 1 April 2026:
options \textbf{0.10\% $\to$ 0.15\%}, futures \textbf{0.02\% $\to$ 0.05\%},
exercise \textbf{0.125\% $\to$ 0.15\%}. Any cost model carrying the October 2024
rates --- which is most of what is published, and was round one's basis --- is one
revision stale.
\end{keyfact}

\subsubsection{Index options, per rupee of one-way premium}

\begin{center}
\begin{tabular}{@{}llr@{}}
\toprule
Component & Rate & bp of one-way premium \\
\midrule
STT, sell leg only & 0.15\% of premium & \textbf{15.00} \\
Exchange transaction charge + IPFT, both sides & \rs 3{,}553/crore $\times$ 2 & \textbf{7.11} \\
GST at 18\% on exchange charges & --- & \textbf{1.28} \\
Stamp duty, buy leg \emph{(unverified)} & 0.003\% & 0.30 \\
SEBI turnover fee \emph{(unverified)} & about \rs 10/crore $\times$ 2 & 0.02 \\
\midrule
\textbf{Regulatory total, round trip} & & \textbf{about 23.7\,bp} \\
Brokerage & broker-specific; \rs 20/order is the discount standard & $+$ \\
\bottomrule
\end{tabular}
\end{center}

Pre-April-2026 the same stack was about 18.7\,bp. \textbf{The options round-trip
floor rose 27\%.}

\subsubsection{Index futures, per rupee of notional}

STT 5.00\,bp (sell leg only) $+$ exchange 3.66 $+$ GST 0.66 $+$ stamp about 0.20
$=$ \textbf{about 9.5\,bp}, against about 6.5\,bp before. \textbf{Futures
round-trip cost rose 46\%.}

\begin{keyfact}{The consequence you actually need}
A short NIFTY weekly straddle collects \rs 11{,}495 per lot on \rs 15{,}71{,}417
of notional (real 28 August 2026 prices). \textbf{One full round trip of that
notional in futures costs 9.5\,bp $\times$ \rs 15.71 lakh $=$ \rs 1{,}493 ---
13.0\% of the premium collected.} Delta-hedging that straddle with total hedge
turnover of twice position notional over its life burns \textbf{26\% of the
premium}. Three round trips of notional and the trade is a coin flip before you
have priced a single greek.

\textbf{The Finance Act 2026 futures STT increase is, functionally, a tax on
delta-hedging.} Combine this with Section~\ref{sec:measure}: hedging cadence is
now the single largest controllable cost line in a short-premium book, and the
band should widen by roughly $(2.8/1.3)^{1/3} \approx 1.29\times$.
\end{keyfact}

\subsection{Tax: the F\&O regime, from the statute}

Section numbers are \textbf{Income-tax Act, 2025}, in force 1 April 2026. The 1961
equivalents are given once in brackets, because most published material still
cites those.

\subsubsection{Non-speculative status}

s.66(31) [old 43(5)] excludes ``a specified derivative transaction'' from the
definition of a speculative transaction. s.66(33) defines that as trading in SCRA
derivatives, carried out through a SEBI-registered broker electronically on a
recognised exchange, supported by a \textbf{time-stamped contract note stating the
unique client code and PAN}.

Note what is \emph{not} required: unlike the commodity limb of the same clause,
the equity-derivative limb does not condition non-speculative status on the
transaction being chargeable to STT. Screen-based, through a broker, with a
compliant contract note --- that is the whole test.

\textbf{Consequence:} NIFTY futures and options are \textbf{business income},
neither speculation business nor capital gains. Intraday cash equity, by contrast,
\emph{is} speculative and lives in a worse ring-fenced regime.

\subsubsection{Deductibility and set-off --- the structural advantage}

\begin{itemize}
\item \textbf{s.34(1)} [37(1)]: brokerage, exchange charges, GST on those, data
feeds, VPS, terminal charges, internet, depreciation on hardware, interest on
borrowed trading capital --- all deductible.
\item \textbf{s.32(1)(k)} [36(1)(xv)]: \textbf{STT itself is deductible}, provided
the income is under business income. By contrast s.72(3)(b) expressly bars STT
deduction when computing capital gains --- so the moment you are classified as an
investor rather than a business, your largest cost stops being deductible.
\item \textbf{s.108}: loss from one source nets against any other source under the
same head --- a losing NIFTY book nets against a winning BANKNIFTY book automatically.
\item \textbf{s.109(1)}: business loss sets off against \emph{any other head
including capital gains} --- \textbf{except} s.109(1)(a), which bars set-off
against \textbf{salary}.
\item \textbf{s.112}: unabsorbed business loss carries forward \textbf{eight
years}, against business profits only.
\item \textbf{s.121}: a loss not determined in a return filed by the due date
\textbf{cannot be carried forward at all}. File on time or the loss is gone permanently.
\end{itemize}

\textbf{The salary carve-out is your live constraint.} A salaried quant with an
F\&O book cannot use trading losses against salary in the loss year. They can be
used against interest, rent, capital gains or other business income; otherwise
they wait up to eight years.

\subsubsection{Turnover, audit, and the presumptive-taxation trap}

Derivative \emph{turnover} for audit purposes is the sum of favourable and
unfavourable differences on squared-off trades \textbf{plus the premium received on
sale of options}. That matters enormously to a seller: a book selling eight NIFTY
straddles a week writes about \rs 47.8 lakh of premium a year before a single
rupee of P\&L is counted. Scale it threefold and you are at \rs 1.5 crore of
``turnover'' on a \rs 25 lakh account.

The audit threshold is \rs 1 crore, \textbf{substituted to \rs 10 crore} where
cash receipts and cash payments are each at most 5\% of the total. A trader whose
entire flow is bank-to-broker satisfies both, so \textbf{the effective threshold
is \rs 10 crore}.

\begin{waste}
\textbf{Do not touch presumptive taxation (s.58, old 44AD).} Turnover includes
written premium, which bears no relation to profit --- so declaring the deemed 6\%
of a \rs 2 crore turnover means declaring \rs 12 lakh of income on a book that may
have lost money. Decline to declare it in a later year, as any honest loss year
forces, and s.58(7) locks you out of the section for \textbf{five years} while
s.58(8) makes books and audit compulsory. It is a one-way door.
\end{waste}

\subsubsection{The STT folklore is now false, and the incentive has reversed}
\label{sec:stt}

Three regimes, from the primary record:

\begin{enumerate}
\item \textbf{Before 1 September 2019}: exercise STT was levied on the
\emph{settlement price} --- the full notional. A NIFTY call \rs 50 in the money
attracted STT of roughly \rs 1{,}970 on a contract worth \rs 3{,}250. That is the
massacre everyone remembers, and it was real.
\item \textbf{1 September 2019 to 31 March 2026}: the base changed to
\textbf{intrinsic value}. The base collapsed by roughly 500-fold. A rate gap
remained --- exercise 0.125\% versus sale 0.10\% --- so letting an ITM long expire
still cost more STT than selling it.
\item \textbf{From 1 April 2026}: \textbf{both entries are 0.15\%}.
\end{enumerate}

\begin{keyfact}{The arithmetic that inverts the standard advice}
At expiry an ITM option's market price is at least its intrinsic value. So for the
holder of a long ITM option: selling it costs $0.15\% \times$ \emph{premium}
(intrinsic plus residual time value); letting it be exercised costs $0.15\%
\times$ \emph{intrinsic}. \textbf{Exercise is now weakly cheaper than sale, and
strictly cheaper whenever any time value remains.}

The ``always square off before expiry to dodge exercise STT'' rule is a 2019-era
artefact. What still argues for squaring off is settlement-price risk and, for
stock derivatives, the delivery-margin ladder --- not STT.
\end{keyfact}

\textbf{The asymmetry that survives:} exercise STT falls on the \emph{purchaser},
who paid no STT on entry. A long ITM position therefore carries a 15\,bp exit tax
on intrinsic that a naive P\&L model omits entirely. On a synthetic long built
from a deep ITM call --- long the 22{,}000 strike with NIFTY at 24{,}176,
intrinsic 2{,}176 --- that is \textbf{\rs 212 per lot, every expiry, forever}, on
a position whose premium is a small fraction of notional.

\subsubsection{Advance tax turns a volatile P\&L into a working-capital problem}

Advance tax is payable if the liability is \rs 10{,}000 or more, and F\&O business
income has no withholding against it, so effectively the whole liability must be
prepaid: at least 15\% by 15 June, 45\% by 15 September, 75\% by 15 December, 100\%
by 15 March. Shortfalls carry interest.

\textbf{The structural problem nobody states:} you must estimate the whole year's
P\&L by 15 June to get the first instalment right. A book that makes its money in
Q4 overpays; a book that gives back Q1 profits in Q4 has prepaid tax on income that
no longer exists and gets it back as a refund a year later. The salary carve-out
compounds it --- you cannot even net a Q3 trading loss against salary already taxed
at source.

\subsection{The two regimes side by side}

Same strategy, same gross edge, same turnover. This is the comparison round one
never made.

\subsubsection{Case 1 --- a low-turnover short-vol book (the favourable case for crypto)}

Real inputs: NIFTY 24{,}175.65 on 28 August 2026, lot 65, notional per lot
\rs 15{,}71{,}417, one 24{,}200 weekly straddle collecting \rs 11{,}495. Book:
\rs 25 lakh equity, eight lots, one cycle per week, 52 cycles. Outcome shape 72\%
win at $+$\rs 40{,}000 and 28\% lose at $-$\rs 63{,}000, giving mean \rs 11{,}160,
standard deviation \rs 46{,}247, and $r = 0.241$ --- a genuinely good short-vol
book with annual Sharpe 1.74.

\begin{center}
\begin{tabular}{@{}lrr@{}}
\toprule
& \textbf{NSE F\&O} & \textbf{VDA regime} \\
\midrule
Gross annual edge & \rs 5{,}80{,}320 & \rs 5{,}80{,}320 \\
Tax base & net of year, after costs & \rs 14{,}97{,}600 (winners only) \\
Tax at 31.2\% & \rs 1{,}81{,}060 & \rs 4{,}67{,}251 \\
\textbf{Net} & \textbf{\rs 3{,}99{,}260 (16.0\%)} & \textbf{\rs 1{,}13{,}069 (4.5\%)} \\
Withholding in-year & --- & \rs 47{,}819 $=$ 42\% of the year's net profit \\
\bottomrule
\end{tabular}
\end{center}

Same strategy, same edge, same turnover: \textbf{the crypto regime keeps 28\% of
what the F\&O regime keeps} --- and this is the \emph{flattering} case, because
turnover is only 1.9 times equity.

\subsubsection{Case 2 --- a high-turnover delta-neutral book (where it breaks)}

Same \rs 25 lakh equity, 5{,}000 round trips a year, \rs 5 lakh notional each ---
\rs 250 crore of annual turnover, a thousand times equity. Per-trade edge 2\,bp of
notional against per-trade volatility of 40\,bp, so $r = 0.050$. Gross edge
\rs 5{,}00{,}000.

\begin{center}
\begin{tabular}{@{}lrr@{}}
\toprule
& \textbf{NSE F\&O} & \textbf{VDA regime} \\
\midrule
Gross edge & \rs 5{,}00{,}000 & \rs 5{,}00{,}000 \\
Regulatory costs & $-$\rs 71{,}100, \textbf{deductible} & $-$\rs 71{,}100, \textbf{not deductible} \\
Tax & \rs 1{,}33{,}817 & \textbf{\rs 13{,}24{,}255} \\
\textbf{Net before withholding} & \textbf{$+$\rs 2{,}95{,}083 ($+$11.8\%)} & \textbf{$-$\rs 8{,}24{,}255 ($-$33.0\%)} \\
1\% withholding on transfers & --- & \rs 2.5 crore $=$ \textbf{ten times the entire equity} \\
\bottomrule
\end{tabular}
\end{center}

The tax alone converts a $+$20\% gross book into a $-$33\% net book. The
withholding makes it unfundable before the tax question is even reached.

\subsection{Margin as a simulable object}
\label{sec:margin}

Named as a requirement in round one and never sourced. Round two implemented the
model from published parameters and public daily files.

\subsubsection{The specification, exactly}

\textbf{PRISM} is NSE Clearing's real-time position monitoring and disablement
layer; \textbf{SPAN} is the margin computation inside it.

\begin{itemize}
\item \textbf{Price Scan Range}, index derivatives: six standard deviations scaled
by $\sqrt{2}$ (the two-day margin-period-of-risk scaling), \textbf{subject to a
floor of 9.3\%} of the underlying.
\item \textbf{Volatility Scan Range}, index: 25\% of annualised EWMA volatility,
\textbf{floor 4\%}.
\item \textbf{Sixteen scenarios}: price moves of $0, \pm\tfrac13, \pm\tfrac23,
\pm1$ times PSR crossed with volatility up and down, all at 100\% weight, plus
$\pm 2 \times$ PSR at \textbf{35\% weight}. Worst-scenario loss is the maximum
across all sixteen.
\item \textbf{Extreme Loss Margin} is a second, \emph{non-portfolio} layer: 2\% of
notional on short index options, 3\% deep OTM, 5\% beyond nine months, plus
\textbf{an additional 2\% on options expiry day}. Notional for a short option is
the value of the underlying index, \textbf{not the premium}.
\item \textbf{Calendar spread charge}: 1.75\% of the far-month contract for index
derivatives, charged \emph{in addition} to worst-scenario loss --- and the benefit
is \textbf{withdrawn entirely on expiry day}.
\end{itemize}

\begin{keyfact}{The structural fact for a spread trader}
SPAN is portfolio-based. \textbf{ELM is not.} It is levied on short option
contracts and \emph{long options do not reduce it}. This is why buying wings
collapses your SPAN and barely touches your total margin.
\end{keyfact}

\subsubsection{The numbers, computed from public files}

Inputs, all from NSE's own files for 28 August 2026: underlying 24{,}175.65, lot
65, notional per lot \rs 15{,}71{,}417; daily volatility 0.847876\%, annualised
16.1986\%; $6\sigma\sqrt{2} = 7.19\%$, so \textbf{the 9.3\% floor binds}, giving
PSR $=$ 2{,}248.3 points and VSR $=$ 4.05 vol points.

\begin{center}
\small
\begin{tabular}{@{}lrrrrr@{}}
\toprule
Position (1 lot) & Premium & SPAN WSL & ELM & \textbf{Initial margin} & \% notional \\
\midrule
Weekly ATM short straddle, 4 DTE & 11{,}495 & 1{,}36{,}769 & 62{,}857 & \textbf{1{,}99{,}625} & 12.70 \\
Weekly $\sim1\sigma$ short strangle, 4 DTE & 718 & 1{,}21{,}546 & 62{,}857 & \textbf{1{,}84{,}402} & 11.73 \\
Monthly ATM short straddle, 32 DTE & 36{,}189 & 1{,}19{,}923 & 62{,}857 & \textbf{1{,}82{,}780} & 11.63 \\
Monthly short strangle, 32 DTE & 3{,}916 & 89{,}198 & 62{,}857 & \textbf{1{,}52{,}055} & 9.68 \\
\textbf{0DTE ATM short straddle} & 11{,}495 & 1{,}36{,}769 & \textbf{1{,}25{,}713} & \textbf{2{,}62{,}482} & \textbf{16.70} \\
0DTE short strangle & 718 & 1{,}21{,}546 & 1{,}25{,}713 & \textbf{2{,}47{,}259} & 15.73 \\
Iron condor (straddle vs wings) & 11{,}008 & \textbf{21{,}492} & 62{,}857 & \textbf{84{,}349} & 5.37 \\
\bottomrule
\end{tabular}
\end{center}

\textbf{Verdict on round one's figures.} The \rs 1.5--2 lakh per lot for a short
strangle is \textbf{correct} --- but for a different reason than supposed: the lot
is now \textbf{65, not 75}. The \rs 3--3.25 lakh for a 0DTE short straddle is
\textbf{about 15\% too high}; the current figure is \rs 2.62 lakh. The old number
was right at lot 75.

\subsubsection{Three non-obvious results worth keeping}

\begin{enumerate}
\item \textbf{SPAN margin for a short option is almost time-insensitive.} The 0DTE
and 4-DTE straddles show the \emph{same} worst-scenario loss. With a 9.3\% price
scan, intrinsic dominates regardless of time to expiry. \textbf{The entire 0DTE
margin uplift --- 24\% --- is the expiry-day ELM, nothing else.}
\item \textbf{Buying wings collapses SPAN by 84\% and total margin by only 58\%},
because ELM does not net across contracts. The iron condor's margin is 75\% ELM.
\textbf{Any margin-efficiency work on defined-risk structures must optimise
against ELM, not against SPAN.}
\item \textbf{The 9.3\% PSR floor was binding} ($6\sigma\sqrt{2} = 7.19\%$), and
has been for most of recent years. \textbf{In calm markets Indian index option
margin is not a VaR number at all --- it is a flat 9.3\% plus 2\% administrative
floor.} Margin therefore does \emph{not} fall when volatility falls, which is
precisely when a short-premium book wants leverage. Model it as a floor, not as a
risk measure.
\end{enumerate}

\subsubsection{The fact that makes a margin-aware backtest feasible}

\begin{keyfact}{Beginning-of-day parameters}
NSE Clearing's own risk FAQ: ``The computation of the margin requirement of
SPAN$+$ELM (both intra-day and EOD) \ldots\ shall be based on the \textbf{fixed
Beginning of Day (BOD) margin parameters}.''

\textbf{Your intraday margin path is fully reconstructible from the BOD parameters
plus the intraday price path.} You do not need the intraday SPAN files --- you
need one parameter set per day and a tick series.
\end{keyfact}

And you must simulate \textbf{peak} margin, not end-of-day margin. Clearing
corporations take a minimum of four random client-wise snapshots per day, and the
higher of the end-of-day and peak shortfalls is penalised, at 0.5\% or 1.0\% per
client per segment per day. \textbf{A strategy that ends the day flat but peaks at
three times margin mid-session has a real, penalised requirement it never shows on
the closing book.}

\textbf{Forced liquidation is an event, not a slippage haircut.} Warnings fire at
70\%, 80\% and 90\% utilisation; \textbf{at 100\% the trading facility is
withdrawn} --- not a margin call, the terminal stops --- and on failure to restore,
NSE Clearing may close out positions by placing counter orders \emph{without any
notice}. Simulate it as a hard state transition: no new orders, and existing
positions liquidated at market by a counterparty who does not care about your exit
price. For a short-gamma book this is the mechanism by which a four-sigma day
becomes a twelve-sigma P\&L.

\subsubsection{Other components a backtest must carry}

\begin{itemize}
\item \textbf{Margin on consolidated crystallised obligation} --- net payable
across premium and futures mark-to-market, blocked as margin.
\item \textbf{Delivery margin} for stock derivatives, laddered from four days
before expiry: 10\%, 25\%, 45\%, 70\%. \textbf{Index derivatives are cash-settled
and escape this} --- which is exactly what kills dispersion (Section~\ref{sec:dispersion}).
\item \textbf{Additional margin on extreme scenarios} --- a second, wider stress
test at a 20\% market fall and 17.74\% rise, sitting behind the 9.3\% PSR and
targeting precisely the naked short-option book.
\item \textbf{Cash collection of option value} on deep-OTM shorts --- requires
\emph{cash}, not collateral.
\item \textbf{At least 50\% of liquid assets must be cash or cash equivalents.} To
carry a \rs 1.84 lakh short strangle you must hold at least \rs 92{,}000 in cash
or T-bills earning a T-bill yield at best. The true funding cost of an Indian
options book is therefore roughly half a repo rate plus half the opportunity cost
of a haircut-reduced pledge, and that belongs in your cost model alongside the
23.7\,bp. \srcnote{Verified at clearing-member level in the SEBI Master Circular;
the identical requirement is applied down to the client by brokers, but the
primary circular imposing it at client level was not reached.}
\end{itemize}

\begin{keyfact}{You do not need the SPAN risk-parameter file, and you cannot get it}
The \texttt{.spn} file is distributed to members via the NSE Clearing extranet;
every anonymous download path returns 404, and the page titled ``SPAN risk
parameter files'' contains only the parameter table with no file links. There is
also \textbf{no public historical archive}, so you cannot backtest historical
margin expansion from primary files.

\textbf{The good news:} everything that matters for index derivatives is either a
published constant (9.3\% PSR floor, 25\%/4\% VSR, 1.75\% calendar spread,
2\%/3\%/5\% ELM plus the 2\% expiry-day add-on) or in a public daily file --- the
volatility file for $\sigma$, the market-lot file for lot size, the bhavcopy for
prices. \textbf{The entire table above was produced from public files and
published parameters, and it is about sixty lines of Python.} Expect single-digit
percentage deviation from the official number, not orders of magnitude. \textbf{No
public, runnable Indian SPAN implementation exists} --- this is a genuinely
publishable artefact and a better use of a week than reading another paper.
\end{keyfact}

\subsection{Indian code and craft: the honest finding}
\label{sec:indiacode}

\begin{keyfact}{Verdict: essentially nobody clears the bar for Indian options}
A systematic GitHub search across every plausible query, reading the actual cost
code in each candidate, found three populations and none of them is what you want:
broker-API wrappers and data libraries (real and useful, but plumbing with no cost
model); live order-placement scripts for short straddles (the bulk of it, with no
backtester behind them at all); and recently vibe-coded ``platforms'', mostly
pushed in the last six months, near-zero stars, README claims far ahead of the code.

Searching ``nifty implied volatility surface'' returns 375 repositories, \emph{none}
of which is a serious NIFTY vol-surface implementation. Searching ``dispersion
trading implied correlation'' returns \emph{zero} India-relevant repositories.

\textbf{For someone who has already implemented SVI and SSVI, there is no Indian
open-source prior art to stand on. You are not behind the field; there is no field.}
\end{keyfact}

What that implies practically: the Indian options-research commons is a
\emph{data} commons, not a \emph{methods} commons. Your comparative advantage
against the domestic retail population is not the model --- it is having a correct
cost and settlement model at all. Treat US and European sources as the methods
library and Indian sources as the data-and-rules library, and do not expect to find
the two joined anywhere.

\subsubsection{An audit, because ``realistic Indian cost model'' was a claim worth checking}

Round one's social sweep flagged one repository as ``the single most on-target''
Indian item, on the strength of a README claiming a realistic Indian cost model.
Round two cloned it and read the code. \textbf{The claim does not hold.} Four
defects, two disqualifying:

\begin{enumerate}
\item \textbf{The STT rate is pre-October-2024}, hardcoded at 0.0625\% with a
docstring claiming currency ``as of 2025''. The bundled sample data spans
July 2025 to June 2026 --- entirely inside the post-change regime. For a
short-premium strategy STT is the largest cost line, and this understates it by
37.5\% of its true value \emph{in the favourable direction}.
\item \textbf{The expiry/exercise path is dead code.} The exercise flag is never
passed anywhere in the repository, and the engine squares off unconditionally at
15:15 --- so the single most-traded Indian structure, an expiry-day short straddle
carried into settlement, is not backtestable at all. The flag is also on the wrong
side: exercise STT is paid by the \emph{buyer}, but the code only reaches that
branch on sells.
\item \textbf{Lot sizes are a static snapshot with no time dimension}, and two are
wrong against the very circular the file cites.
\item \textbf{Slippage is one tick, flat} --- which is 0.025\% on a \rs 200 ATM
option and 2.5\% on a \rs 2 wing. No spread model, no depth, no queue, and fills
occur at the close of the same bar the decision was made on.
\end{enumerate}

\textbf{Salvage:} the architecture is decent and the historical expired-chain
integration is genuinely valuable --- expired option chains without a paid vendor.
Fork it for the data layer, throw away the cost model, and never use the shipped
sample results.

\subsubsection{The two Indian sources that do survive}

\textbf{A first-person cost post-mortem.} An independent write-up documenting a
signal with a monthly Spearman information coefficient of 0.1758 ($n=106$,
$t \approx 12.4$) --- genuinely predictive --- that nonetheless returned
\textbf{$-$7.15\% a year with Sharpe $-$3.27} over 2015--2024, from 541 trades a
year at about 0.35\% round trip. The cost stack is modelled from first principles
including Almgren-Chriss market impact and stock-lending cost, not a flat
basis-point assumption, and the limitations section names survivorship bias as the
single most important caveat and flags that the short leg ignores the F\&O ban
list. Passes tests (2), (3) and (4). \emph{Caveat: the code is not published, so
it fails (1), and it is cash equities rather than options. Read it for the cost
discipline.}

\textbf{One correctly-built cost model in open source.} A small repository whose
cost module gets delivery equity right and --- critically --- \textbf{puts every
rate in a configuration file rather than in source}, precisely so a regulatory
rate change does not silently invalidate old runs. That is the design decision the
audited repository got wrong, and it is the one to copy. It also benchmarks
against a \emph{costed} buy-and-hold and publishes a table in which three of its
four strategies lose to it, and says so. \emph{Caveat: no options.}

\textbf{One genuinely useful instrument.} A live-trading framework that
\emph{records realised slippage} --- market price at decision versus actual fill,
per order type, to CSV. That is an empirical slippage distribution for Indian
options rather than an assumption. There is no backtest cost model in it; use it
as an instrument, not a backtester.

\textbf{And one academic thread worth pulling.} Kumar, Agarwalla, Varma \&
Virmani, \emph{Harvesting the volatility smile in a large emerging market: A
Dynamic Nelson--Siegel approach} (\emph{Journal of Futures Markets} 43(11), 2023)
models the NSE smile as a term structure \emph{in delta}, estimates it by Kalman
filter, sorts on model-implied mispricing and trades the extreme deciles. This is
the closest published thing to ``SVI/SSVI, but for NSE, evaluated as a trade.''
\textbf{It is paywalled and round two could not confirm whether it nets
transaction costs} --- which, given the strategy trades single-stock options
monthly across deciles, is the question that decides whether the result survives.
\textbf{Pull it via IIT KGP and read the cost treatment first.} This is the single
most important open question in the Indian half of this report.

\subsection{Dispersion in India: the measured verdict}
\label{sec:dispersion}

Computed from the actual NSE F\&O bhavcopy for 27 August 2026, cross-referenced
against the live NIFTY 50 constituent list.

\begin{center}
\begin{tabular}{@{}lrrrr@{}}
\toprule
& Underlyings & Series listed & Series with a trade & Notional turnover \\
\midrule
Stock options & 210 & 23{,}741 & 7{,}893 (33\%) & \rs 3.25 lakh crore \\
Index options & 6 & 4{,}700 & 1{,}083 (23\%) & \textbf{\rs 77.3 lakh crore} \\
\bottomrule
\end{tabular}
\end{center}

\textbf{Index options are twenty-four times the entire single-stock option
complex.} For the NIFTY 50 constituents specifically: all fifty have listed
options, so on a naive screen the trade looks feasible. But aggregate constituent
option notional is \rs 1.41 lakh crore against NIFTY index options at \rs 74.9
lakh crore --- a \textbf{53:1} ratio. Median per constituent is \rs 1{,}789 crore
of notional across 46 traded strikes, and \textbf{eleven of fifty are below
\rs 1{,}000 crore a day}, the weakest at \rs 214 crore across 36 traded strikes.

\begin{keyfact}{Verdict: not viable --- and liquidity is the \emph{least} of the four problems}
\begin{enumerate}
\item \textbf{Single-stock options are physically settled.} Index options are
cash-settled; stock options are not. A long constituent leg held into expiry
generates a delivery obligation of \emph{full notional} on every ITM leg, with
delivery margins ramping from four days out. Across fifty legs this is not a margin
inconvenience --- it is a different business.
\item \textbf{The expiry calendars no longer line up.} Since 1 September 2025
single-stock options are \textbf{monthly only}. NIFTY has weeklies. The natural
weekly dispersion trade has no stock leg at all, and the monthly trade puts your
entire short-index leg into the one expiry where the $+$2\% ELM and the removed
calendar-spread benefit both bite.
\item \textbf{Margin has no correlation offset across underlyings.} SPAN scans each
underlying independently. The whole economic point of dispersion --- that the
basket is nearly hedged --- earns you nothing.
\item \textbf{Eligibility churn kills the backtest even where it spares the trade.}
See below.
\end{enumerate}
\end{keyfact}

\textbf{And the ratio points the wrong way anyway.} Published work on Indian
single-stock options finds the embedded risk premium is \emph{higher} than in other
markets and higher than theory suggests --- implied volatility consistently exceeds
realised ``by too large an amount to be easily explained as a rational risk
premium.'' That is the wrong sign for classic dispersion, where you want to
\emph{buy} cheap constituent vol. If anything the Indian data argues for
\emph{reverse} dispersion --- short constituent vol, long index vol --- which is
exactly the leg physical settlement punishes hardest.

If you still want correlation exposure, BANKNIFTY against its twelve constituents
is the one place the basket is small enough and every name is top-decile liquid.
Both legs are now monthly-only, which at least makes the calendars match --- at the
cost of one observation per month.

\subsubsection{Why the single-stock universe is non-stationary}
\label{sec:eligibility}

Eligibility is determined by a \emph{rolling six-month window on cash-market
statistics}, re-reviewed continuously. In August 2024 the thresholds were tripled
in one step: median quarter-sigma order size \rs 25 lakh $\to$ \textbf{\rs 75
lakh}, market-wide position limit \rs 500 crore $\to$ \textbf{\rs 1{,}500 crore},
average daily delivery value \rs 10 crore $\to$ \textbf{\rs 35 crore}. Fail any
criterion for three consecutive months and the stock exits, with \textbf{no
re-inclusion for one year}. A separate product-success framework adds turnover and
open-interest floors.

\textbf{Any study that takes today's roughly 210-name F\&O list and runs it
backwards has survivorship bias of exactly the worst kind}: names were removed
\emph{because} their liquidity deteriorated. There is no free point-in-time
constituent file --- you must build one by differencing the daily bhavcopy
universe, which you can do.

\subsection{What could not be verified in the Indian material}

Stated explicitly rather than glossed.

\begin{itemize}
\item \textbf{SEBI turnover fee} and \textbf{stamp duty} rates on derivatives ---
the regulations page serves its text via a JavaScript call that could not be
resolved, and the Stamp Act schedule 404'd. Combined they are about 0.3\,bp and
move no conclusion.
\item \textbf{GST at 18\% on brokerage} --- stated from general knowledge, not
verified against the primary notification.
\item \textbf{Whether NSE transaction charges are still at the September 2024
level} in August 2026. Check the circular archive before using the number.
\item \textbf{The Income-tax Act 2025 ``as amended by Finance Act 2026''}
consolidated PDF is 403-blocked to non-browser clients. The gap was closed by
reading the Finance Bill's amendment clauses directly and confirming none of the
load-bearing sections were touched --- but re-verify in a browser before relying on it.
\item \textbf{Whether the November 2024 $+$2\% expiry-day ELM has since been
modified} by any later risk circular. It took effect; it was not traced forward.
\item \textbf{Whether the DNS smile paper nets transaction costs.} The single most
important open question in this section.
\end{itemize}

\newpage
\section{Crypto}
\label{sec:crypto}

Crypto was your top-weighted area, and it splits cleanly in two. As a
\textbf{research venue} it is excellent and under-exploited: free tick data, real
term structure, a risk-based margin engine at small account size, and a surface
whose behaviour genuinely differs from equity index vol in ways nobody has
published with costs. As a \textbf{P\&L venue} for an Indian resident it is, on
the arithmetic in Section~\ref{sec:vda}, close to unusable for exactly the
strategies you are training to build.

This section is what round two measured against live APIs on 28 August 2026. Every
number below is reproducible with an unauthenticated request.

\subsection{Strategy families: alive or decayed}

\subsubsection{Perpetual funding carry --- alive, but the risk is not where retail thinks}

Funding has structurally compressed: 2021 bull-phase prints of 0.15--0.20\% per
eight hours (over 200\% annualised) have given way to a 2024--2026 baseline near
\textbf{0.01\% per eight hours, about 11\% annualised} before costs, spiking only
in leverage-euphoria windows.

\begin{keyfact}{The failure mode now has a date and a body count}
On 10--11 October 2025, roughly \$19\,bn of open interest was liquidated across
1.6 million accounts, \$3.21\,bn in a single minute. The mechanism that matters:
\textbf{auto-deleveraging forcibly closed market makers' short perp legs against
their long spot} --- the exchange took away the profitable side of the
delta-neutral book, leaving participants unhedged long into a 14\% drawdown. A
delta-neutral yield product marketing about 18\% subsequently depegged more than
77\%, with post-mortems tracing the damage to ADL breaking hedges.

\textbf{ADL is an uncompensated, unhedgeable forced exit that fires precisely at
the tail}, and it ranks counterparties by profit times leverage --- so being
\emph{more} profitable and \emph{more} hedged puts you \emph{earlier} in the
queue. This has no analogue on NSE or the OCC. Whatever risk framework you carry
over from the index work will not model it.
\end{keyfact}

Capital is meaningful only above roughly \$50--100k, because you run two collateral
pools and must hold buffer against margin calls on both. Latency is low for entry
and \emph{high for the unwind} --- October punished slow unwinds. A runnable
multi-venue implementation exists in Hummingbot; read its open issue documenting a
double fill on Hyperliquid after the strategy is stopped before you trust it. Your
hedge is only neutral if the order state machine is correct.

\srcnote{The claim to disbelieve: a ``19.26\% annual, 0.80\% max drawdown, positive
all 12 months of 2025'' figure circulates widely with no attribution and no cost
model. A strategy that was positive in all twelve months of 2025 was, by
construction, not honestly accounting for October.}

\subsubsection{Market making --- alive on CEX perps, folklore on AMMs}

Solo market making is reachable on a small number of mid-cap perps, with a
matched-region VPS. It is not reachable on BTC/ETH top-of-book against
professional firms. Before assuming the economics, look at what a protocol's own
professionally-run market-making vault earns \emph{with} rebate and liquidation
privileges you will not have: historical returns in the range of 3--8\% annualised
with negative months, and profit concentrated in a few high-volatility events.
Your unlevered solo expectation should be \emph{below} that, not above it.

\begin{waste}
The AMM liquidity-provision version is mostly a losing trade dressed as yield. An
LP is profitable only if fees exceed loss-versus-rebalancing, and one analysis of
seventeen concentrated-liquidity pools found \$199.3m of fees against \$260.1m of
impermanent loss --- LPs were \$60.8m worse off than simply holding. Concentrated
liquidity is a short-gamma position sold to arbitrageurs at a price you do not
set, and headline fee APR is gross of LVR. Skip unless you are specifically
researching LVR mitigation.
\end{waste}

\subsubsection{Cross-exchange and triangular arbitrage --- decayed to folklore}

No source passing the filter shows a solo retail trader profitably running latency
arbitrage in 2025--2026. The arithmetic: cross-venue spreads sit roughly
0.1--2\%, and at 0.1\% taker fee per side you need 0.2\% to break even --- before
slippage, before the fact that capturing it requires pre-positioned inventory on
both venues, and before the funding cost of that idle split capital. The one
legitimate residual is \emph{funding-rate dispersion across venues}, which is
slow-moving and capital-rewarded rather than speed-rewarded --- and that is the
carry trade, not this.

\subsubsection{Momentum --- partly decayed, and the decay is regime-shaped}

Time-series momentum produced negative annual returns through the choppy
2022--2023 regime, while cross-sectional momentum limited losses better because
its market-neutral construction survives the absence of a trend. The effect is now
regime-conditional rather than unconditional.

\begin{waste}
Crypto time-series momentum is the most academically over-published and
practically fragile item here. \textbf{Every backtest you will find applies no
funding cost to the perp leg} --- and holding a levered long perp through a
positive-funding regime costs you precisely the 11\% a year the carry trade is
trying to earn. Momentum backtests that ignore funding are not backtests. Build
the funding-inclusive version; it will look materially worse than every published
result, and that gap is the honest answer.
\end{waste}

\subsection{Deribit market structure, measured}

\subsubsection{The listing ladder is a step function, and that is the big difference}

Deribit lists four dailies, three weeklies, three monthlies and four quarterlies,
with strike intervals that step by tenor and moneyness --- \$500 at four days or
less, \$1{,}000--\$5{,}000 at a week to a month, \$2{,}000--\$10{,}000 at three
months. Measured ladders confirm the published policy exactly.

\textbf{This is the single biggest structural difference from NIFTY.} Deribit's
strike spacing is a step function of tenor and moneyness, not a constant, and the
far expiries carry strikes to four times spot. NIFTY gives you a dense uniform
ladder over a narrow band; Deribit gives you a sparse non-uniform ladder over an
enormous one. Your surface fitter cannot assume a uniform grid.

Open interest is extremely concentrated: the two big quarterly expiries carried
\textbf{68.7\% of all BTC open interest} at the time of measurement.

\subsubsection{The fee cap is the number with no NIFTY analogue}

Options fees are \textbf{3\,bp of underlying notional, capped at 12.5\% of
premium}, maker and taker alike --- there is no maker rebate in options. Worked
through at spot near \$79{,}400:

\begin{center}
\begin{tabular}{@{}rrrrrl@{}}
\toprule
Premium & as \% spot & Fee one way & \% of premium & Round trip & Binding \\
\midrule
\$100   & 0.13\% & \$12.50 & 12.50\% & \textbf{25.0\%} & cap \\
\$190   & 0.24\% & \$23.75 & 12.50\% & \textbf{25.0\%} & cap \\
\$437   & 0.55\% & \$23.82 & 5.45\%  & 10.9\% & notional \\
\$1{,}826 & 2.30\% & \$23.82 & 1.30\%  & 2.6\%  & notional \\
\$7{,}275 & 9.16\% & \$23.82 & 0.33\%  & 0.7\%  & notional \\
\bottomrule
\end{tabular}
\end{center}

The cap binds below a premium of \textbf{0.24\% of spot}. Below that you pay a
flat 12.5\% each way. Above it you pay a fixed dollar amount per BTC \emph{regardless
of the option's value}.

\begin{keyfact}{NSE charges on premium; Deribit charges on notional --- and it inverts which part of the surface is tradeable}
Every OTM wing and every short-dated option on Deribit carries a fee that is
enormous relative to its premium, while long-dated expensive options are nearly
free. Any strategy that trades cheap options frequently --- gamma scalping the
dailies, wing buying, ratio structures --- is dead on arrival at retail fee tier.
This is the exact opposite of the NIFTY incentive.
\end{keyfact}

\subsubsection{The tick is a second hidden tax}

Deribit's inverse-option tick is a fixed \emph{BTC} amount, so its dollar value
scales with spot, and it steps up at higher prices. Measured across two-sided BTC
options:

\begin{center}
\begin{tabular}{@{}lrrrr@{}}
\toprule
Expiry & $n$ & Tick as \% of premium & Quoted spread \% & Median premium \\
\midrule
Front daily   & 41  & 2.41\% & \textbf{16.87\%} & \$1{,}648 \\
Next daily    & 48  & 2.80\% & \textbf{16.58\%} & \$1{,}241 \\
$\sim$1 week  & 61  & 2.08\% & 6.67\% & \$1{,}866 \\
$\sim$1 month & 112 & 0.66\% & 3.18\% & \$6{,}044 \\
$\sim$10 month& 100 & 0.33\% & 2.06\% & \$12{,}148 \\
\bottomrule
\end{tabular}
\end{center}

Across all two-sided BTC options with premium under \$500: \textbf{median tick is
9.09\% of premium and median quoted spread is 24.00\%}. Combine with the fee
table: on a sub-\$200 option you pay 25\% round-trip in fees \emph{plus} cross a
roughly 24\% spread. \textbf{Your edge must exceed roughly half the premium.} The
Deribit wings and dailies are untradeable at retail size, and that is a fee-and-tick
fact, not a liquidity opinion. Concentrate on the big quarterlies where open
interest is and spreads are 2--3\%.

\subsection{The crypto surface behaves differently, and it is measurable in thirty lines}

Live 25-delta skew term structure, computed from marks at spot \$79{,}207 (BTC)
and \$2{,}505 (ETH):

\begin{center}
\begin{tabular}{@{}lrrrrr@{}}
\toprule
Expiry & Days & BTC ATM & BTC RR25 & BTC BF25 & ETH RR25 \\
\midrule
Front daily & 0.7   & 33.1 & \textbf{$+$1.14} & $+$0.86 & \textbf{$+$2.82} \\
1 week      & 6.7   & 36.8 & \textbf{$+$2.44} & $+$0.76 & \textbf{$+$3.66} \\
1 month     & 27.7  & 37.5 & $+$0.26 & $+$1.05 & $+$1.11 \\
2 months    & 62.7  & 37.5 & $+$0.06 & $+$1.17 & $+$0.99 \\
3 months    & 90.7  & 39.4 & \textbf{$-$0.28} & $+$0.78 & $+$0.48 \\
4 months    & 118.7 & 40.3 & \textbf{$-$0.38} & $+$0.75 & $-$0.27 \\
10 months   & 300.7 & 42.5 & \textbf{$-$0.85} & $+$1.35 & $+$0.27 \\
\bottomrule
\end{tabular}
\end{center}

Two things an equity-index trader would not predict, both directly observable:

\begin{enumerate}
\item \textbf{The skew sign flips with tenor.} Short-dated BTC risk reversals are
\emph{positive} --- calls bid over puts by one to 2.4 vol points --- and go
negative only beyond about ninety days. ETH is call-skewed nearly the whole curve.
In SPX the risk reversal is negative at every tenor. \textbf{Any model you port
from equities with a hard-coded negative correlation parameter will fight the
front of this surface.}
\item \textbf{Term structure is in steep contango while the front is call-skewed}
--- cheap near vol, call-skewed, upward sloping. That combination is the crypto
signature.
\end{enumerate}

Butterflies are uniformly positive and remarkably stable at $+$0.7 to $+$1.4
across the whole curve for both coins, which is the one thing that does look like
equities.

\subsubsection{The best answer to ``does SVI calibrate here, and what broke''}

Echenim, Gobet \& Maurice (CNRS / \'Ecole Polytechnique with Kaiko's quant data
lead), free full PDF, using real Deribit BTC data. It contains:

\begin{itemize}
\item A \textbf{theorem} that calibrating SVI on mid-prices produces a
\emph{systematic underestimation} of implied volatility in the wings. Not a
heuristic --- a proof, with the direction of the bias.
\item The reason: ``the relative tick size (and thus the relative spread) for
options on cryptos is roughly \textbf{10 to 50 times larger} than for options on
traditional markets.'' The measurement above independently confirms it.
\item An \textbf{anchor price} built from bid and ask whose bias is third-order in
the tick size, replacing the mid.
\item A data-augmentation procedure for one-sided and missing quotes --- exactly
the failure you will hit on Deribit wings.
\item Hard liquidity counts: roughly \textbf{30 BTC option trades per hour across
all strikes} for one-week and one-month expiries, about 5 per hour for three- and
six-month, and 1--2 per hour at nine months to a year. That is the whole market.
\item Measured improvement: calibration error down 25--30\%.
\end{itemize}

\srcnote{Caveat: no code is released and the data is paid, so you implement the
anchor price and the augmentation yourself. The sample predates the current
liquidity regime --- re-measure the trade counts before believing them. Pair it
with Magadini, Sinclair \& Chepal on real-time crypto SVI calibration, which gives
you the architecture and problem framing but describes its calibration process as
``proprietary''. The Kaiko paper gives the maths the Sinclair paper withholds;
together they are a working method.}

\subsubsection{Two corrections to make before you fit anything}

\begin{enumerate}
\item \textbf{The inverse-contract numéraire.} BTC options on Deribit are
coin-margined: premium paid in BTC, settled in BTC. In the BTC numéraire this is
not a vanilla --- it is a self-referencing payoff, and the implied vol quoted is
the USD-numéraire vol of a payoff you never receive. Build a surface from Deribit
marks and price a stablecoin-settled structure off it and you are wrong by the
correlation between underlying and settlement asset, which here is one. Deribit
lists linear options too, and the two families carry different vols for the same
strike.
\item \textbf{Use the forward, not the index.} Deribit publishes its exact
convention, defining the rate as $R = \ln(F/X)/T$ and stating that ``the implied
volatility displayed uses the forward, rather than the index, as the underlying
price.'' Get this wrong and your implied vols will disagree with the screen by a
spread-sized amount, and you will spend a week chasing it.
\end{enumerate}

\begin{keyfact}{The weekend effect is real, and it is not a modelling artefact}
Deribit's own documentation settles it: the exchange uses \textbf{pure calendar
time, with no business-day or weekend weighting whatsoever}. So the weekend dip in
quoted implied vol going into Friday is not something you must undo --- it is a
genuine, tradeable consequence of the exchange charging calendar theta the market
knows will not be delivered. \textbf{Whether it survives a 25\%-of-premium
round-trip on short-dated options is precisely the question, and nobody has
published the answer with costs.}
\end{keyfact}

\subsection{DVOL: an excellent instrument and a non-existent market}

DVOL's construction differs from VIX in ways that matter. It uses a
\textbf{depth-weighted price}, not a mid --- walking five levels or two BTC of
cumulative size. It \textbf{discards ITM options and anything under five delta},
explicitly to avoid strikes where tick size swamps vega. It then applies an
exponential moving average over the last 240 one-second points.

What it therefore does \emph{not} measure: it is not a pure 30-day variance swap
rate --- the five-delta truncation systematically \textbf{understates} the true
variance swap strike, in the same direction as the mid-price bias above. The EMA
means it \emph{lags} by construction and is unsuitable for anything faster than a
few minutes. And the depth-weighted construction means DVOL moves when
\emph{liquidity} changes, not only when volatility does.

\begin{keyfact}{Is it tradeable? Measured answer: no}
Querying every listed instrument returns exactly \textbf{one} DVOL future. Live
state at measurement: bid 40.0 / ask 41.0 (a full vol point wide), \textbf{open
interest 8.1 contracts --- about \$324 notional --- and 24-hour volume of \$160},
with one trade that day. There is no ETH DVOL future at all. The listing policy
confirms this is by design: one monthly contract.

\textbf{Treat DVOL as a data series, never as a hedge.} The valuable part is that
the \emph{index history is free and multi-year} at arbitrary resolution, verified
returning data for March of 2022 through 2026. That is your variance-risk-premium
measurement input.
\end{keyfact}

\subsection{The perp-option-future triangle is closed}

Put-call parity for inverse options gives $F = K/(1 - (C-P))$. Comparing the
parity-implied forward against the listed future's mark across the whole curve:

\begin{center}
\begin{tabular}{@{}lrrrr@{}}
\toprule
Expiry & Strike & $F$ implied & $F$ future & Difference \\
\midrule
Front daily & 79{,}000 & 79{,}194.03 & 79{,}183.36 & $+$0.0135\% \\
1 week      & 79{,}000 & 79{,}277.47 & 79{,}253.65 & $+$0.0301\% \\
1 month     & 79{,}000 & 79{,}436.90 & 79{,}422.16 & $+$0.0186\% \\
3 months    & 79{,}000 & 80{,}142.02 & 80{,}111.69 & $+$0.0379\% \\
7 months    & 80{,}000 & 81{,}321.47 & 81{,}324.95 & $-$0.0043\% \\
10 months   & 80{,}000 & 82{,}346.89 & 82{,}334.72 & $+$0.0148\% \\
\bottomrule
\end{tabular}
\end{center}

\textbf{Maximum deviation across the whole curve: 0.038\%, or \$30 on \$79{,}000.}
Against frictions on the same book: one option tick is \$8 to \$40, the median
front-month ATM spread is 6.67\% of premium, and the fee is \$23.82 per BTC per
leg. \textbf{A parity arbitrage needs four legs. The mispricing is an order of
magnitude smaller than the cost of one.}

Not exploitable at retail size, and the reason is structural rather than cyclical:
three markets on one venue, one margin engine, one index, arbitraged continuously
by professionals. \textbf{What \emph{is} usable is the inverse} --- the
parity-implied forward is a free, accurate, per-expiry forward that you should use
as the underlying in your implied-vol calculations, exactly as Deribit itself
does. That materially changes your fitted skew.

The futures term structure is a clean flat carry curve at about \textbf{4.8\%
annualised} from three months to ten, which is exactly why the triangle closes.

\srcnote{An important tension worth flagging. A published \emph{Journal of
Financial Markets} paper on crypto derivative arbitrage reports strategies
remaining profitable ``even under conservative transaction cost scenarios'' ---
more optimistic than this live measurement. Two likely reconciliations: their
sample is older and efficiency has improved, and academic ``conservative
transaction costs'' rarely include the 12.5\%-of-premium fee cap or the
9\%-of-premium tick that bite hardest on the cheap options where deviations look
largest. \textbf{Read their profitability claims as an upper bound and re-test on
current data with the cost model above.} That re-test is itself worthwhile work.}

\subsection{Where Indian residents can actually trade}

\subsubsection{Deribit: not restricted, not registered --- and those are different things}

Deribit's restricted-jurisdictions list was read directly. It names about
twenty-five jurisdictions. \textbf{India does not appear.} From Deribit's side, an
Indian resident is not blocked. Deribit is also not on the FIU-IND enforcement
lists --- neither the nine offshore providers served notices in December 2023 nor
the twenty-five named in the October 2025 action.

\textbf{But Deribit is not FIU-IND registered.} That is a live regulatory gap, not
a clearance. Note the direction of travel: the October 2025 action shows FIU-IND
enumerating offshore venues serving Indian users and issuing takedowns, and states
that the obligations are activity-based and not contingent on physical presence in
India. Deribit is now Coinbase-owned, and Coinbase itself registered with FIU-IND.
\textbf{Treat access as revocable, and do not build anything whose viability
depends on Deribit remaining reachable from India.}

\srcnote{And a caveat on the registration claim for the domestic venue: \textbf{FIU-IND
does not publish a machine-readable public register.} Three official PDFs were
fetched and rendered --- they turn out to be AML guidelines and the registration
\emph{process} circular, not a list of registrants. Aggregate counts come from
press releases and parliamentary answers. So a venue's stated registration number
is a \emph{self-declaration} that could not be independently confirmed against a
government register, because no such public register exists. It is credible but
unverified, and you should know the difference.}

\subsubsection{The measured depth comparison, which reverses round one}

Round two measured both books properly --- not five levels, which is meaningless
across venues with different ticks, but \textbf{notional resting within $\pm$2\%
and $\pm$5\% of mid premium}, which is tick-independent.

\begin{center}
\begin{tabular}{@{}llrrrrr@{}}
\toprule
Contract & Venue & Mid \$ & bid $\pm$2\% & ask $\pm$2\% & bid $\pm$5\% & ask $\pm$5\% \\
\midrule
$\sim$1-day ATM call  & Deribit & 437 & \textbf{0.0} & \textbf{0.0} & \textbf{0.0} & \textbf{0.0} \\
$\sim$1-day ATM call  & Domestic & 488 & 24.1 & 14.9 & 58.6 & 36.7 \\
$\sim$1-week ATM call & Deribit & 1{,}826 & 0.0 & 0.0 & 49.4 & 51.6 \\
$\sim$1-week ATM call & Domestic & 1{,}838 & 18.8 & 12.2 & 40.3 & 23.4 \\
$\sim$1-month ATM call& Deribit & 3{,}593 & 37.8 & 33.5 & 86.3 & 111.8 \\
$\sim$1-month ATM call& Domestic & 3{,}595 & 12.7 & 8.5 & 40.1 & 21.1 \\
\bottomrule
\end{tabular}
\end{center}

\textbf{The zeros are not missing data.} On Deribit's daily expiry, one tick is
already more than 2\% of premium, so \emph{nothing can rest inside a $\pm$2\%
band}. The domestic venue's finer tick makes its short-dated book genuinely
tighter --- measured spreads of 1.75--2.34\% on the one-day ATM against Deribit's
10.9--19.1\% on the same contracts.

\begin{keyfact}{The honest reversal}
For \textbf{short-dated} BTC options the Indian-registered venue is \emph{tighter
and more accessible} than Deribit, because Deribit's fixed-BTC tick is punitive on
cheap options. \textbf{Deribit wins decisively at one month and beyond} --- three
to five times the resting notional, 130 strikes against 46 --- and wins
overwhelmingly on tenor and strike range. Round one had this backwards.
\end{keyfact}

The domestic venue's own book measured about \textbf{\$486m of open interest and
\$1.10bn of 24-hour turnover, with every single listed strike two-sided} across
nine expiries --- reflecting enforced market-maker obligations, against Deribit's
41 of 54 two-sided on its front expiry. Its costs are structurally the same
notional-based charge with a slightly tighter cap, \emph{plus 18\% GST on the
fee}. Lot size is far smaller. There is \textbf{no India price premium}: measured
simultaneously, spot differed by $-$0.007\% and the perp mark by $-$0.037\%. The
folklore applies to INR-funded spot venues, not to a USD-margined derivatives
book. There is no cross-venue basis trade here.

\subsection{Data: what is free, and what you must start recording today}

\begin{itemize}
\item \textbf{Deribit's own API is mostly a dead end for history.} Verified
negatives, so you do not waste a week: trade history by currency returned
\textbf{zero} trades for the same date in 2023, 2024, 2025 and 2026; expired
instruments returned only the most recent expiry, so you cannot enumerate
historical contracts. \textbf{The exception is DVOL} --- multi-year, free, at
arbitrary resolution.
\item \textbf{The realistic free deep history} is a first-day-of-every-month
options-chain snapshot, downloadable with no API key, with a verified schema
carrying microsecond timestamps, \emph{both exchange and local timestamps} (so you
can measure feed latency), full greeks, both-side implied vols and sizes, and open
interest. Coverage starts March 2019. \textbf{Caveat: about 4\,GB compressed per
day}, and the free days are monthly snapshots --- fine for surface cross-sections
and calibration research, \textbf{useless for event studies or path-dependent
backtests}.
\item \textbf{Consequence: if you want continuous Deribit options history, you
must either pay or start recording now.} A single call returns the entire chain
with mark implied vols and open interest. Stand up a snapshot recorder today; in
six months you own something you cannot otherwise get free.
\end{itemize}

\srcnote{Honest negative on code, after searching by both relevance and recency:
\textbf{there is no maintained, credible, open-source crypto options analytics
stack.} Almost everything currently active in this space is zero-star, days old
and AI-generated. The one solid recommendation is a maintained, MIT-licensed
stochastic-volatility pricing library that is the reference implementation behind
the published inverse-options paper --- but it is pricing analytics only, with no
data layer and no execution. \textbf{Given you have already implemented SVI and
SSVI, you are better served writing your own thin layer than adopting any of these.}}

\subsection{The crypto variance risk premium, and the gap worth filling}

The most rigorous current measurement puts the Bitcoin variance risk premium at
about \textbf{14\% annualised} at the one-month horizon against roughly 2\% for
the S\&P 500, with a Sharpe around 0.8. The regime result is the useful part:
\textbf{the premium is high in the \emph{low}-volatility regime and lower in
high-volatility states} --- the opposite of equity intuition. And the return
decomposition finds large \emph{positive} returns account for 39\% of the Bitcoin
premium, against under a third of SPX's premium coming from positive returns at
all. \textbf{The right tail carries the information in crypto}, which fits the
positive front-end call skew measured above.

\textbf{But there are no transaction costs anywhere in that paper, and the sample
ends before the 2025--26 regime.} A 14\% variance premium is comfortably inside
the round-trip costs computed above for short-dated options. \textbf{Do not read
14\% as a strategy Sharpe} --- it is a statement about the risk-neutral versus
physical gap, not about a harvestable return.

\begin{keyfact}{The highest-value unclaimed work in the crypto half}
Every 2026 claim about the crypto VRP traces to vendor blogs or exchange-affiliated
commentary with no published methodology. DVOL history is free back to at least
2022, and minute-wise BTC spot is trivially obtainable. \textbf{Compute
$\text{DVOL}^2$ against subsequent 30-day realised variance yourself, split by
volatility regime.} It is a two-day project, it is genuinely original for 2026,
and it settles the question with your own measurement rather than an assertion.
\end{keyfact}

\newpage
\section{Regime detection and your senior's stack, stress-tested}
\label{sec:regime}

Your colleague built these. This section is about which survive contact with live
trading. The default posture is scepticism, because much of this is academically
popular and practically fragile --- and where that is true it is said directly.

\subsection{Hidden Markov models: the weakest link, and the reason is operational}

\begin{keyfact}{The bug that accounts for most impressive HMM backtests}
A Gaussian-emission HMM fitted by Baum--Welch gives you \emph{smoothed} state
probabilities $P(s_t \mid y_1 \ldots y_T)$. Every published equity curve using
smoothed states \textbf{uses the future to label the present}. The only legitimate
object is the \emph{filtered} probability $P(s_t \mid y_1 \ldots y_t)$, and
filtered probabilities are far noisier and lag the turn by roughly the persistence
horizon implied by the transition matrix.

\textbf{If your colleague's HMM looks good, the first thing to check is whether it
called a Viterbi decode over the full sample rather than running an
expanding-window filter.} The demonstration takes about thirty lines: plot filtered
and smoothed side by side, plot the equity curve under each, show the gap. That is
more persuasive than any paper.
\end{keyfact}

\textbf{Second defect: state-count instability.} Information criteria on financial
returns are close to flat between two, three and four states, and the argument
minimum moves as you extend the sample. A four-state model typically resolves into
something interpretable on one window and into a degenerate near-empty state on the
next. \textbf{A state that exists in one training window and vanishes in the next
is not a regime; it is a fitting artefact.} Test it directly: refit on rolling
five-year windows and check label correspondence via the Hungarian algorithm on
state means and covariances. If labels permute or a state collapses, the model is
not deployable.

\textbf{Third: Gaussian HMMs reproduce fat tails and zero return autocorrelation
but fail to reproduce volatility clustering} --- the one property you actually
need for options. \srcnote{Flagged as unverified: stated in the jump-model
literature and not independently confirmed. It is cheap to check with a simulation
and worth doing, because if true it means an HMM regime label is a poor
conditioning variable for a vol-selling book precisely because it does not capture
the autocorrelation of absolute returns.}

\subsubsection{What to use instead}

Statistical jump models, with a scikit-learn-style API and discrete, continuous and
\emph{sparse} variants. The jump penalty is an \emph{explicit} regulariser on state
persistence --- you are tuning the thing HMMs leave implicit and unstable in the
transition matrix. The sparse variant is the one that matters: it tells you which
candidate features actually carry regime information, which kills a lot of
feature-fishing.

\srcnote{Real caveat: no evidence any of the authors ran this with client money.
``Jump model beats HMM beats buy-and-hold'' is an in-sample-of-the-research-programme
result, and the jump penalty is a free parameter that does exactly what state count
did in the HMM --- it moves the overfitting one level up. Tune it inside a purged
cross-validation, not on the full sample.}

\subsubsection{Who has run it with money}

\begin{keyfact}{Honest answer: nobody identifiable}
No named fund could be verified as attributing P\&L to an HMM regime filter. The
claim that large systematic managers ``use regime models'' is folklore that could
not be sourced, and you should not repeat it.
\end{keyfact}

\textbf{Where it stops working is the moment you want it.} Regime detection is a
change-point problem with a detection delay bounded below by the signal-to-noise of
the feature. In a fast volatility event the label flips \emph{after} the gap move.
For a short-gamma book, a regime filter that flips after the gap is worse than
useless --- it makes you flatten into the bottom.

\textbf{Use regime labels only for slow decisions} --- position sizing, hedge
budget --- \textbf{never as a gate on entering or exiting short premium.}

\subsection{Spillover and connectedness: excellent instrument, close to worthless as a signal}

What it actually is: a rolling-window generalised forecast-error variance
decomposition of a reduced-form VAR, renormalised so rows sum to 100. It is a
\emph{measurement} of realised co-movement transmission. It is not a forecast, and
the originators frame it as measurement.

The implementations are genuinely good and ship replication data for a battery of
published papers, so the field passes tests (1) and (5). It passes neither (2) nor
(3): \textbf{there is no cost model anywhere in this literature because there is no
trade in it.}

\subsubsection{Four failure modes, stated bluntly}

\begin{enumerate}
\item \textbf{The rolling-window artefact.} The index jumps when a crisis
\emph{enters} the window and drops when it \emph{exits}, $W$ days later, with no
market event at all. Practitioners repeatedly mistake the exit-drop for
``spillover normalising''. Every such time series has this sawtooth baked in.
Time-varying-parameter VAR exists specifically to kill it and is the version to use
if you use any.
\item \textbf{Generalised FEVD does not identify anything.} It sidesteps ordering
dependence, but the price is that directional spillovers are not structural ---
they are variance shares under a non-orthogonal shock, and rows only sum to 100
after an ad-hoc renormalisation. \textbf{You cannot read ``A causes B'' off the
matrix.} Papers do this constantly. Do not.
\item \textbf{It is contemporaneous with volatility, so it adds nothing to a vol
signal.} Total connectedness rises in crises --- so does realised volatility,
correlation and VIX.
\item \textbf{Out-of-sample instability.} Horizon, lag order and window length are
three free parameters and the index level is sensitive to all three. The
\emph{shape} over time is far more robust than the \emph{level}, so never threshold
on the level.
\end{enumerate}

\begin{keyfact}{The one test that decides whether this survives, and it costs an afternoon}
Before building anything on a connectedness matrix, \textbf{regress your
connectedness index on cross-sectional realised correlation and average realised
volatility.} If the residual is small --- and it usually is --- the matrix is a
repackaged volatility signal with twenty times the maintenance cost. Run this
first, and tell your colleague the result either way.
\end{keyfact}

\textbf{The one use worth defending} is as a cross-asset diversification
\emph{diagnostic}, not a signal. Net directional connectedness tells you which of
your positions are net \emph{receivers} of shocks, and therefore where your
apparent diversification is fake. For a book short NIFTY vol and long crypto vol,
``how much of BTC's forecast error variance is explained by equity and vol shocks
in a stress window'' is a legitimate, decision-relevant number. Run it monthly,
look at the matrix with your eyes, and do not automate a trade off it.

\textbf{And note a specific trap for your case:} NIFTY, BTC and SPX do not share a
session. A VAR estimated on misaligned bars creates spurious directional spillover
\emph{from the market that closes later}. That is a mechanical artefact, not a finding.

\subsection{Clustering for parameter selection: it does not do what it claims}

\begin{keyfact}{The honest framing}
Picking a ``robust plateau'' in parameter space via clustering is a \emph{smoothing
operator} on the backtest-Sharpe surface. Smoothing reduces the variance of the
selected parameter's out-of-sample performance and slightly reduces selection bias
versus taking the argmax. \textbf{It does not address the fact that you ran $N$
trials.} If you tried 500 combinations and then clustered them, you still tried 500
combinations, and your best cluster centroid still needs deflating by the same
multiple-testing correction.

\textbf{Clustering that just relabels overfitting is exactly the right suspicion.}
\end{keyfact}

\textbf{The recursion is the fatal objection.} A density-based clusterer has its
own hyperparameters --- so using it to pick your strategy's hyperparameters means
you now have hyperparameters selecting hyperparameters, tuned on the same data.
People then Bayesian-optimise \emph{those}, adding a third layer. Each layer is
another set of trials that nobody counts.

A nearest-neighbour approach for this purpose is simply local averaging of the
Sharpe surface. If that is what you want, say so and use a Gaussian kernel smoother
--- one line, one bandwidth parameter, and honest about being a smoother.

\srcnote{A targeted search found \emph{no} credible finance practitioner source
applying density-based clustering to strategy-parameter selection. What comes back
is microscopy, astronomy and the documentation.}

\textbf{What to do instead:} drop the clustering. Replace it with (a) sensitivity
of Sharpe to each parameter one at a time, plotted; (b) a requirement that the
strategy be profitable across a \emph{contiguous region}, not at a cluster
centroid; and (c) \textbf{count every trial you ran, including the ones you
abandoned, and deflate}. Item (c) is the one everyone skips and the one that matters.

\srcnote{One caveat that cuts against the standard advice for \emph{your} case:
purged cross-validation matters most when labels overlap. \textbf{For a
minute-wise intraday NIFTY options strategy that opens and closes within the
session, label overlap is near zero and purging buys you almost nothing.} Your real
leakage risks are different --- expiry-day survivorship in the chain,
strike-selection using the day's realised ATM, and quoting off a stale mid.}

\subsection{Delta-neutral construction: this one survives}

It is the only item in your colleague's stack where the cadence question is
genuinely solved with closed-form answers you can implement this week. That
material is Section~\ref{sec:measure} and is not repeated here. The single most
important point to carry back:

\begin{keyfact}{The correction your colleague may be missing}
The optimal response to transaction costs is \textbf{not} ``hedge less often on a
timer''. It is a \textbf{no-trade band around delta whose width scales as the cube
root} of (cost $\times$ gamma $\times$ spot / risk aversion). Time-based
rebalancing is dominated by band-based rebalancing at essentially every cost level,
and the gap widens with gamma. With minute NIFTY data you can measure this directly.
\end{keyfact}

\textbf{Where it stops working:} the band solutions all assume a diffusion. In a gap
--- budget day, RBI policy, a crypto liquidation cascade --- the framework is silent
and the gamma cost is realised instantly regardless of cadence. And on expiry day
the asymptotic assumptions break entirely (Section~\ref{sec:measure}).

\textbf{The crypto-specific gap worth knowing:} Deribit is 24/7 with no weekend gap,
so the diffusion assumption behind the bands holds \emph{better} than in NIFTY ---
but perpetual funding introduces a carry term in the hedge leg that none of the
classical band literature models, and that term can dominate the transaction-cost
term. That is unfilled, and it is a genuinely publishable piece of work.

\subsection{Trend following and options are the same trade}
\label{sec:trendopts}

The best-founded item in this section, and the only one with a named person
managing real money who published the analysis free.

The identity is operational, not decorative. Replicated straddle P\&L equals the
payoff range minus $\tfrac{1}{2}\sum \Gamma_{s-1}(S_s - S_{s-1})^2$, and
\textbf{that quadratic-variation term is named as the slippage cost of
replication}. Theta in the option is identified as the \emph{same quantity} as
whipsaw plus transaction cost in the trend program. That gives you a direct
arithmetic test:

\begin{keyfact}{A runnable decision rule}
Compute $\tfrac{1}{2}\sum \Gamma (\Delta S)^2$ for your trend rule and compare it
to the actual traded straddle premium. \textbf{Whichever is cheaper is how you
should express the trade.} You can run this on minute NIFTY data.
\end{keyfact}

The rest of the mapping: the moving-average crossover threshold is the straddle
strike; position-sizing ramp is delta; ramp speed is gamma. \textbf{Strangles
correspond to entry thresholds} --- spreading the strikes is exactly ``require
confirmation before entering'', and the stated cost is slower convexity accrual,
``the cost of patience''. \textbf{Stop-losses correspond to knock-out options}: a
trend program with stops is \emph{short a tail put below the stop level} and
exposed to jump risk there. The prescription, if the option is cheap enough, is to
replace the stop with the option outright.

\srcnote{The caveat that matters: \textbf{there is no backtest in that paper.} No
cost-modelled equity curve, no parameter set, no data appendix --- it is an
analytic framework plus illustrative exhibits. It also has an obvious commercial
interest: a tail-risk manager arguing that trend followers should buy options is
arguing for its own product. The framework is correct and useful regardless, but
do not treat it as evidence that the overlay pays after costs. \textbf{You have to
build that test, and you can}, because the whole thing reduces to comparing
$\tfrac{1}{2}\sum\Gamma(\Delta S)^2$ against traded straddle prices. To their
credit they volunteer the negative: over the most recent two decades their 50/50
allocation did \emph{not} keep up with 100\% equity on CAGR and won only on
risk-adjusted terms.}

\subsubsection{The finding that invalidates most options-overlay backtests}

\begin{keyfact}{Rebalance timing luck}
A self-financing three-month put-spread collar was replicated in implementations
differing \emph{only in rebalance date}. The result: \textbf{annualised tracking
error in excess of 400\,bp between implementations.} For equity factor indices the
luck term exceeded 100\,bp a year.

\textbf{Any options overlay you backtest on a fixed roll schedule has a
several-hundred-basis-point random component from the roll date alone.} The fix is
overlapping tranches --- roll a fraction of the book each period. If your backtest
platform rolls on a single fixed date, its results are not reproducible in the
sense that matters.
\end{keyfact}

\subsection{The empty \texttt{regime} column in your platform}

Not an oversight to fix by picking a model. It is a schema stub, and what it should
hold depends on a decision nobody has made. In production systems this field is
normally one of four things:

\begin{enumerate}
\item \textbf{A low-cardinality categorical label per bar} from a small enumerated
set. Almost always \emph{not} from an HMM --- it is a deterministic rule over slow
observables: trailing realised volatility bucketed by its own trailing percentile,
term-structure slope sign, a drawdown-from-high flag. Deterministic rules are used
because they are auditable, have no refit, and cannot silently leak.
\item \textbf{A continuous score plus a threshold}, stored as the score, bucketed at
query time. Preferred by teams burned by state-count instability.
\item \textbf{An as-of-bar filtered probability vector} with an explicit model
version and fit-as-of column. The presence of a fit-as-of column is the tell that
whoever built it understood look-ahead.
\item \textbf{A post-hoc analyst label} used for slicing P\&L reports and
explicitly forbidden from entering signals.
\end{enumerate}

\textbf{The distinction to enforce before anyone populates it:} (1) to (3) are
point-in-time and legal as inputs; (4) is retrospective and is an attribution field
only. If it is going to be filled, it needs a sibling boolean or the two uses will
get mixed within a quarter --- and a smoothed retrospective label leaking into a
signal is exactly the bug at the top of this section.

\begin{keyfact}{The useful thing to say in the meeting}
\textbf{An empty regime column is a better state than a wrongly populated one.}
Ask what decision the field is supposed to change. If the answer is ``position
size'' or ``hedge budget'', a deterministic volatility-percentile rule is
sufficient and costs a day. If the answer is ``whether to enter short premium'',
\textbf{the field should stay empty}, because detection lag makes it actively
harmful there.
\end{keyfact}

\newpage
\section{People, books, and where to actually read}
\label{sec:people}

\subsection{The seed video, audited}

You sent one video as the exemplar. It is worth stating plainly what the audit
found, because the pattern generalises and recognising it will save you far more
time than any single recommendation here.

\begin{keyfact}{The exemplar fails the filter on five of six counts}
\begin{itemize}
\item Primary call to action is a paid coaching programme at roughly \$3{,}000 to
\$5{,}000, \textbf{including trade alerts and algorithmic signals}. Course plus
signals is the exact shape the filter excludes.
\item Brokerage-sponsored with an affiliate link, and carrying the platform's own
paid-promotion badge.
\item \textbf{The ``verified performance'' link does not survive a click.} The
linked portfolio shows a name, a city, a Verified badge --- and then ``Trades are
hidden for this user.'' No equity curve, no trades, no drawdown. The verification
is a badge with the substance switched off.
\item \textbf{The trade evidence is about 23 seconds.} By the video's own chapter
list, ``Real Trade Examples'' runs 10:32 to 10:55. The other nine minutes are
whiteboard mechanics.
\item Every title in the channel is a dollar amount plus a win rate --- 89\%,
91\%, 92\%, 96\% --- advertised without loss-size context, which is precisely the
metric that flatters short premium and hides the tail.
\item The channel is fourteen months old and already sells a multi-thousand-dollar
programme.
\end{itemize}
\end{keyfact}

\textbf{To be fair to it: double calendars are a completely legitimate structure},
and the mechanics section may well be sound. The problem is the register and the
unfalsifiable record, not the instrument.

\begin{keyfact}{The generalisable tell}
\textbf{Any claimed track record whose per-trade fills you cannot inspect is not a
track record.} Several better-credentialled US names, and most of the Indian
option-selling ecosystem, are the same object at different production budgets.
\end{keyfact}

\subsubsection{What to read instead, on the same subject matter}

{\small
\setlength{\LTpre}{8pt}\setlength{\LTpost}{8pt}
\begin{longtable}{@{}p{4.6cm}p{11.0cm}@{}}
\toprule
\textbf{Question} & \textbf{Honest source} \\
\midrule
\endfirsthead
\toprule
\textbf{Question} & \textbf{Honest source} \\
\midrule
\endhead
\bottomrule
\endlastfoot
Does premium selling have an edge at all, and how do I size it? &
Sinclair's position-sizing work --- the only treatment that starts from ``here is
the cost of being wrong about the fourth moment'' rather than from a win rate. \\[3pt]
Double calendars and diagonals &
\textbf{There is no good dedicated free source, and that is informative.} The
correct framing is \emph{forward variance and term-structure slope}, which you
already have. A double calendar is a bet on forward vol between two expiries plus
a short skew position --- price it as such and the folklore evaporates. Nobody
teaching calendars as a ``range play'' is pricing the forward vol. \emph{And on
NIFTY, per Section~\ref{sec:noterm}, there is barely a second expiry to trade.} \\[3pt]
Delta-neutral short vol: the failure modes &
Kris Sidial on liquidity gaps and why the hedge does not fill; Benn Eifert on why
the structural short is a levered carry trade; and the QVR wind-down itself. \\[3pt]
The actual net cost of doing this in India &
SEBI's studies --- roughly \rs 25{,}000 crore a year of transaction costs is the
number every Indian income-strategy backtest omits. \\
\end{longtable}
}

\subsection{Practitioners, gated on a verifiable CV}

Verification standard: \textbf{V} means confirmed against at least one source
independent of the person's own marketing; \textbf{S} means self-reported or
repeated only by promotion-adjacent outlets.

\subsubsection{The highest-value names}

\textbf{Euan Sinclair} --- \emph{the single highest-value name here for you.}
PhD theoretical physics; roughly 25 years professional options trading; risk
management and strategy design at a proprietary options firm. \textbf{V}. His free
SSRN papers are directly implementable: the risk-reversal premium, the skewness
and kurtosis of European option returns and the implications for trade sizing,
volatility estimation via first exit times. \textbf{The trade-sizing paper is the
one to read first} --- it connects the third and fourth moments of an option's own
return distribution to position sizing, which is exactly the gap between your
surface work and an actual position. Passes tests (2), (3), (4) doubly, and (5).
\emph{Caveat: US-centric, and some of the specific premia have decayed. Transfer
the method, not the parameters.}

\textbf{Timothy Klassen} --- see Section~\ref{sec:measure}. Goldman under Derman,
designed Cboe's 2003 VIX, built Getco's options analytics. \textbf{V}. The richest
free corpus in this document for spot-vol dynamics and P\&L attribution.
\emph{Caveat: vendor, selling a library, and his SVI criticisms are made by an
interested party --- benchmark them against your own NIFTY fits.}

\textbf{Kris Abdelmessih (Moontower)} --- SIG trainee, then market maker across
several option pits, then a relative-value commodities vol book. \textbf{V}. The
best free writing anywhere on thinking in payoff space, and repeatedly explicit
about which retail vol ideas are simply the bid-ask spread in disguise. \emph{He
writes for intuition, not rigour --- almost nothing there is a finished testable
specification, and converting his intuitions into hypotheses is your job. Skip
roughly 40\% of the Substack, which is personal-finance and lifestyle.}

\textbf{Kris Sidial (Ambrus)} --- co-CIO of a live volatility-arbitrage and
tail-risk fund; previously an exotic derivatives desk at a bank and proprietary
trading. \textbf{V}, corroborated across several independent outlets. Unusually
explicit about \emph{how a tail hedge actually fills} --- liquidity, cost of carry
on convexity, and the fragility of options-market plumbing. Passes test (4) hard:
he names where his own trade fails, which is that a slow grind-down bleeds him out
and he needs discontinuity. No paid product for retail.

\textbf{Cem Karsan} --- mostly verified, with one claim that does not check out.
Founded an SPX volatility-arbitrage market-making firm; his current firm is a
registered commodity trading advisor, which is a regulator record rather than a
bio. \textbf{V}. \emph{But the widely circulated ``32 years of exchange trading
experience'' is arithmetically impossible against a 1999 undergraduate degree ---
it appears to be firm-level language that migrated into his personal bio via
speaker directories. His real book is roughly 2000 to present, which is still
about 25 years. Marketing drift, not fraud, but he has never corrected it.}
\textbf{Good for} the mechanical, flow-based half of index vol.
\textbf{Ignore} the macro and political narrative layer --- it is unfalsifiable
and unrelated to the part of his work that has edge.

\begin{keyfact}{The severe caveat on Karsan, which is load-bearing}
\textbf{He has never published the computation.} In none of the free long-form
does he specify how gamma or vanna exposure is estimated from open interest, what
the dealer-sign assumption is, or how it is validated. \textbf{The framework as he
presents it is unfalsifiable as stated}, and the commonly cited underlying
methodology is anonymous and attached to a paid data product --- an auto-reject on
two counts under your own filter.

\textbf{Read Barbon \& Buraschi's \emph{Gamma Fragility} first} --- free, named,
peer-reviewed, and it finds intraday momentum under negative aggregate dealer gamma
and reversal under positive, \emph{conditional on illiquidity}. That conditioning
is the failure mode Karsan never states. Then read Karsan as colour. With NIFTY
open-interest data and minute bars this is directly replicable on an Indian index
and, as far as can be determined, unpublished there.
\end{keyfact}

\textbf{Benn Eifert (QVR)} --- PhD economics; head of quantitative research and
derivatives trader on a bank proprietary desk; co-founder of a real vol-arbitrage
manager. \textbf{V}. Years of free, technically precise threads on structured
product flows, dispersion, how bank vol supply distorts surfaces, and a sustained
public demolition of retail ``options income'' strategies.

\begin{keyfact}{The most instructive single fact in this section}
\textbf{QVR's multi-strategy fund lost roughly 30\% between January and April 2026
and is winding down}, from a peak around \$1.6bn, with the management business up
for sale. This does not invalidate his analysis. It is the empirical footnote:
\textbf{the most technically credentialled, most cost-aware, most
failure-mode-literate vol manager in the retail-visible world still blew up.} Read
the back catalogue \emph{and} read the blow-up. Do not treat anyone's process as
proof of survivability --- including your own.
\end{keyfact}

\textbf{Harel Jacobson (volquant)} --- roughly 15 years in FX volatility trading;
all free, no paywall, no funnel. \textbf{The closest thing to a free written
market-maker course.} His FX volatility smile series matters to you even though you
trade index: \textbf{Part II explains \emph{why} the smile has the shape it has as
a consequence of dealers hedging exotic books} --- risk reversals to hedge vanna
from barriers, butterflies to hedge volga from digitals. That is a supply-side
explanation of skew and convexity that no equity-options text gives you. One piece
includes runnable code, so he passes test (1). \emph{Caveat: FX conventions do not
transfer without translation, and he says so.}

\subsubsection{Verified and rejected, with reasons}

\begin{itemize}
\item \textbf{Imran Lakha} --- background is the cleanest of any candidate
(\textbf{V}: Citigroup head of European index options, then head of a bank's
European options desk, then macro PM). \textbf{Auto-rejected on outlet}: the
primary funnel is a paid membership plus a course. His free weekly vol commentary
remains usable as \emph{market colour} --- how a sell-side index-vol book thinks
about skew and implied correlation --- never as a source of a testable claim, and
the free tier is deliberately incomplete.
\item \textbf{A well-known calendar-and-butterfly educator} --- CV is
\textbf{S} only: every restatement traces to bios he supplied, and no independent
primary record was found. Auto-rejected on the paid mentoring funnel. His
pre-course-era webinars are pedagogically decent and contain no reproducible
testing. He is the closest US analogue to the video you sent --- better
credentialled, same commercial structure.
\item \textbf{A prominent proprietary firm's ``options desk''} --- \textbf{reject,
and the implied pedigree is weaker than presented.} What actually checks out
(\textbf{V}) materially contradicts the marketing impression: the lead is a career
property-and-casualty insurance executive who left insurance in 2006 and has traded
options since --- not an ex-market-maker or ex-desk trader. The firm is real; the
options \emph{desk} is substantially an education product. Skip entirely.
\item \textbf{A large brokerage's media arm} --- the founder's CV is genuine and
large (\textbf{V}: about 20 years as a floor market maker, founded two brokerages).
\textbf{Auto-rejected twice}: it is the media arm of a broker, so order flow is
literally the business model, and their ``studies'' are the marketing. Documented
failures to reproduce their backtests exist. Use it for vocabulary and the
\emph{mechanics} of managing a short-premium book --- rolling, assignment, margin
--- and then test every quantitative claim yourself. \textbf{Do not treat a single
one of their numbers as an input.}
\end{itemize}

\subsection{Indian practitioners: the honest answer}

\begin{keyfact}{Almost nobody survives, and the bar was not lowered to manufacture entries}
Checked and rejected: the major broker-integrated platforms (order flow is the
business model, and one widely repeated ``ex-Goldman'' claim does not check out);
the best free Indian options primer, which is broker-affiliated and also below your
level; the largest quant-education blog, whose parent is a genuine high-frequency
firm with verifiable people but which exists to funnel into a paid programme; and
the closed-box backtesters, several of which \emph{do} let you set slippage and
brokerage --- more than most US retail content offers --- but whose fill logic you
cannot inspect, so nothing produced there is reproducible.

\textbf{And the structural problem: the serious Indian quant firms are real, with
verifiable people, and they publish nothing.} The competent Indian options
community is entirely inside proprietary firms and does not write.
\end{keyfact}

\textbf{Two Indian sources do pass, and they are not people:}

\begin{enumerate}
\item \textbf{SEBI's studies on individual traders in equity derivatives.} The only
Indian document that states the true all-in cost drag. It tells you precisely who
is on the other side of your fills and what the toll booth costs. \emph{Caveat:
regulator aggregate data with a policy agenda; it tells you the population loses,
not why any individual strategy fails.}
\item \textbf{Academic work on Indian market microstructure} by the people who
built the institutions --- foundational quantitative work on Indian microstructure
and index construction, published free, plus first-person writing on Indian
derivatives structure. \emph{Caveat: policy and microstructure economics, not a
trading strategy, and much of the derivatives-specific work is from the late 1990s
and 2000s. Read it for how NIFTY liquidity and impact cost are actually structured.}
\end{enumerate}

\begin{keyfact}{The blunt conclusion}
There is no Indian equivalent of Sinclair or Abdelmessih. Given that you will hold
minute-wise NIFTY and Sensex data and can model costs against SEBI's published
aggregates, \textbf{the honest recommendation is to import technique from the US
and crypto names above and become the missing Indian source yourself}, rather than
keep searching for one.

\textbf{The same applies to video.} All Indian long-form options content is
broker-affiliated education, course funnels, or ``I made \rs X'' content. For
Indian-market-specific craft there is \emph{no free practitioner video layer, and
you should not go looking for one.} Your Indian inputs should be exchange and
regulator primary documents plus your own data; the technique layer should come
from elsewhere, translated.
\end{keyfact}

\subsection{The book canon, ranked for your position}

Ranking premise: you have the surface maths. What you lack is (i) how a book
behaves over time, (ii) what the costs actually are, (iii) how to tell whether a
model is broken. Ranked by marginal value \emph{given that}, not by general quality.

\subsubsection{Access reality first}

\begin{keyfact}{Do this before spending anything --- it is a two-minute email}
The national journals consortium your library leads with is a \textbf{journals}
deal and does not cover monographs. Indian institutional subscriptions typically
include a publisher's \emph{journals}, not its \emph{online books} --- and Sinclair,
Taleb, Gatheral, L\'opez de Prado and Rebonato are all with the same publisher, so
the likely outcome is that \textbf{none of them is reachable and you buy them}.
Springer/Palgrave e-books are commonly held, which makes one title on this list
most likely free to you. McGraw-Hill Professional is effectively never in an Indian
academic package.

\textbf{Ask the library for the e-books A--Z list, and specifically whether the
publisher's \emph{Online Books} and Cambridge Core \emph{Books} --- not just
journals --- are in the subscription.} That single question resolves roughly
\rs 40{,}000 of purchasing decisions.
\srcnote{Flagged: the library site is login-gated, so all access claims here are
inference from publisher and consortium patterns, not verification.}
\end{keyfact}

\subsubsection{Tier 1 --- in this order}

\textbf{1. Sinclair, \emph{Volatility Trading} (2nd ed.).} The only book that
treats an options trade as a \emph{statistical bet with a cost structure} rather
than a payoff diagram. The hedging chapter derives the path-dependence of a
delta-hedged option's P\&L explicitly --- why your realised P\&L is not the vol
difference times vega but a gamma-weighted path integral. The money-management
chapter is where most quant-trained people are weakest. \textbf{Chapters that
matter:} volatility measurement, hedging, money management, trade evaluation.
\textbf{Skippable:} the Black-Scholes derivation and the options-basics chapter.
\textbf{Dated:} pre-0DTE, US-centric, and its GARCH-family forecasting chapter is
the weakest part. \textbf{No free equivalent exists.} This is the book with no
substitute.

\textbf{2. Bergomi, \emph{Stochastic Volatility Modeling} --- and chapters 1 and 2
are free on the author's own site.} The correct answer to ``what happens to my book
when spot moves''. Chapter 1 is titled around \emph{How (in)effective is delta
hedging?} and derives the P\&L of a delta-hedged short option from break-even
accounting, then shows empirically how much hedging error survives. That is the
attribution content, free, from a bank's head of quantitative research.
\textbf{Chapters that matter:} 1 and 2 (free), then the smile-of-stochastic-vol
chapter where the skew-stickiness ratio is derived, and the static-versus-dynamic
and what-causes-equity-smiles chapters. \textbf{Free equivalent:} the two free
chapters plus his \emph{Smile Dynamics} papers cover perhaps 40\% of the practical
value at zero cost --- \textbf{read those before deciding to buy.} \textbf{Caveat:}
dense idiosyncratic notation and no code. You will be re-deriving.

\textbf{3. Taleb, \emph{Dynamic Hedging}.} The \emph{bookkeeping} of an options
book --- how a market maker organises risk across hundreds of positions in one
name. The only book-length treatment of \textbf{vega bucketing} by someone who did
it on a desk, and the only serious treatment of what happens \emph{between} the
greeks: pin risk, wing risk, early-exercise pathologies. \textbf{Dated and in
places wrong:} it predates SVI, SABR, everything you have implemented, and has no
coherent spot-vol dynamics --- precisely the gap Bergomi and Klassen fill. Its
exotics pricing is obsolete. \textbf{Read it for the taxonomy of risk, not for
models.} Badly organised and badly edited; budget for frustration. Nothing has
replaced it.

\textbf{4. Bouchaud et al., \emph{Trades, Quotes and Prices}.} The cost side,
quantitatively: square-root impact, order-flow long memory, the propagator model,
why your backtest's fill assumptions are the dominant error term. \textbf{Real free
equivalent:} Bouchaud, Farmer \& Lillo, \emph{How Markets Slowly Digest Changes in
Supply and Demand} (arXiv:0809.0822), about 100 pages covering a large fraction of
the relevant parts. \textbf{Read that before buying.} \textbf{Caveat:} options
microstructure is essentially absent, which matters because NIFTY option books
behave differently from the underlying.

\subsubsection{Tier 2 --- conditionally}

\textbf{Sinclair, \emph{Positional Option Trading}.} The value is the negative
result: most named option ``strategies'' carry no edge, and the edge where it
exists is in sizing, selection and exit. Keep the bet-sizing and process chapters;
skip the strategy catalogue. Substantial overlap with \emph{Volatility Trading} ---
buy second, not first.

\textbf{Gatheral, \emph{The Volatility Surface}.} You have absorbed the
parameterisation content already. What remains valuable is the skew-dynamics
chapter and the variance-swap chapter. \textbf{His free Baruch lecture notes
overlap the book heavily --- use the notes, skip the book.}

\textbf{Natenberg, \emph{Option Volatility and Pricing}.} For you, roughly 70\% is
below your level. The residue that is not: how a spread's greeks migrate as spot and
time move, and the market-maker chapters. \textbf{Fails test (2) --- no cost model
anywhere. Recommendation: skip the purchase.}

\textbf{A short Palgrave monograph on the greeks and hedging} whose publisher
description names exactly the gaps above --- day conventions, interest rates,
sticky deltas. \textbf{The title on this list most likely to be free to you.}
Caveat: thin; it will not replace Bergomi.

\textbf{L\'opez de Prado, \emph{Advances in Financial Machine Learning}.} Include
it \emph{only} for four things: triple-barrier labelling, meta-labelling,
purged/embargoed cross-validation, and the deflated Sharpe ratio. Everything else
is contested. \textbf{Nearly all the substance is in his free SSRN papers.}
\srcnote{Honest negative: a targeted search for a substantive published critique of
this book's reproducibility found none from a credible named practitioner --- the
criticism circulates as forum folklore, not as a citable document. Treat its methods
as \emph{unreplicated} rather than either validated or refuted.}

\subsubsection{Tier 3 --- do not buy}

\textbf{Rebonato, \emph{Volatility and Correlation}} --- the majority is
interest-rate volatility, irrelevant to NIFTY index options, and the relevant parts
are superseded by Bergomi. \textbf{Sinclair, \emph{Option Trading}} --- largely
superseded by his other two. \textbf{A 2024 retail-oriented options title} --- given
your level, skip; credible review describes it as oscillating between explaining
what a long call is and computing rolling volatility z-scores.

\textbf{Books that earn a slot and were not on the original list:} Derman \&
Miller, \emph{The Volatility Smile} (though Derman's Columbia lecture notes are free
and cover most of it); Larry Harris, \emph{Trading and Exchanges}, still the
canonical market-structure text and cheap; Lehalle \& Laruelle, \emph{Market
Microstructure in Practice}, the execution-cost companion to Bouchaud.

\newpage
\section{Income and honest base rates}
\label{sec:income}

Factual only. This section reports what is documented --- capital levels, cost
structures, published loss rates, and which skills monetise where. It contains no
recommendations, no allocations and no trade suggestions.

\subsection{Two facts that reorder the question before any maths}

\begin{enumerate}
\item \textbf{US options are closed to you as a venue} (Section~\ref{sec:lrs}).
Anything you read about US retail options economics is technique transfer, and the
constraint is legal rather than stylistic.
\item \textbf{The crypto tax regime is anti-selected against exactly the strategies
you are training to build} (Section~\ref{sec:vda}), and the dominant crypto options
venue is not FIU-registered.
\end{enumerate}

So of the candidate strategy classes, \textbf{two are venue-constrained before you
reach a line of maths}, and the realistic choice set is NSE index derivatives
versus Indian-domiciled crypto venues.

\subsection{Capital versus viability, with numbers}

\subsubsection{Indian index options: the floor is the flat fee, then the margin}

The ad-valorem charges are not the problem at small size --- a one-lot round trip
on a short option costs roughly 1.3\% of a good gross P\&L. \textbf{The flat
per-order brokerage is the problem}, because it is per \emph{order}, not per lot,
and so dilutes about tenfold going from one lot to ten.

A four-leg iron condor is eight orders. At one lot, brokerage plus GST alone is
roughly 7\% of maximum profit --- and if you exit at 30\% of maximum, which is the
standard management rule, \textbf{costs are about 23\% of the realised target}.

\begin{keyfact}{The two capital floors}
\textbf{Defined-risk multi-leg selling} becomes fee-viable at roughly ten lots,
which is about \rs 1.5--2 lakh of margin. \textbf{Undefined-risk selling} has a
much higher floor driven by margin, not fees: from the computed table in
Section~\ref{sec:margin}, \textbf{\rs 1.52--1.84 lakh per lot for a short strangle
and \rs 2.62 lakh for a 0DTE short straddle}. Add the requirement that at least
half of margin be cash or cash-equivalent, so a \rs 20 lakh ``account'' is not
\rs 20 lakh of deployable margin.
\end{keyfact}

\textbf{Scaling arithmetic --- division, not a return assumption.} At a net 1\% a
month on deployed margin you need roughly \rs 1 crore of margin to gross \rs 1 lakh
a month; at 2\% a month, about \rs 50 lakh. Then F\&O profit is business income at
slab. \textbf{There is no configuration of this arithmetic in which a sub-\rs 25
lakh account produces income competitive with a finance salary inside two years.}

\textbf{And the regulator is an active adversary of the small-account version.} The
November 2024 package --- contract size raised to \rs 15 lakh-plus, weeklies cut to
one index per exchange, STT raised, then raised again in April 2026 --- was designed
to raise the cost floor for exactly this activity.

\subsubsection{Crypto}

The lowest capital floor mechanically --- no lot size, fractional sizing --- with
the floor set by fees, transfer costs and margin buffer. \textbf{Then the tax code
eats it} (Section~\ref{sec:vda}): for an Indian resident the after-tax carry can be
negative while the gross carry is positive. That is the single most
under-appreciated fact in this section.

For crypto \emph{options} specifically, serious-participant treatment on the deep
venue starts around six figures of account equity and market-maker terms require a
negotiated agreement --- and it is not the registered venue anyway.

\subsection{The base rates}

\subsubsection{SEBI --- and note the numbers you had are superseded}

\begin{center}
\begin{tabular}{@{}lrrr@{}}
\toprule
& FY24 & FY25 & FY26 \\
\midrule
Net losses (\rs crore) & 74{,}812 & \textbf{1{,}12{,}000 (revised)} & \textbf{91{,}685} \\
Loss-maker share & 91.1\% & \textbf{90.9\%} & \textbf{87.7\%} \\
Unique individual traders (lakh) & 86.3 & 96.0 & \textbf{87.5} \\
\bottomrule
\end{tabular}
\end{center}

The August 2026 study revised FY25 \emph{upward} from \rs 1.06 to \rs 1.12 lakh
crore --- so the figure in circulation is stale. FY26 shows the first
year-over-year decline in active traders since FY16, down 18\%, with new entries
down 40\%. \textbf{But average loss per person rose 2\%}, and 23\% of traders
account for about 90\% of total losses.

Options accounted for \textbf{92\% of aggregate individual losses}, and 93\% of
individuals traded \emph{only} options. \textbf{99\% of foreign-investor and
proprietary profits came from algorithmic entities.}

\begin{keyfact}{A claim from round one that must be retracted}
Round one inferred from these studies that ``a large share is naive short vol.''
\textbf{The studies cannot support that and do not attempt to.} A full-text search
of the current study for direction-related terms returns six hits in 91 pages, none
about strategy composition. The methodology table lists the complete categorisation
--- product category, age, income, gender, city tier, turnover, portfolio size,
frequency, experience. \textbf{There is no direction field. No long/short. No
buy/write.} The sign of the position is never in the dataset, and SEBI states its
own limit: ``the analysis is descriptive and does not establish cause-and-effect.''
\end{keyfact}

\subsubsection{And SEBI has now published the document that settles it}

A companion \emph{behaviour} study classifies traders by strategy using
clearing-corporation settled-value, traded-value and margin data, on a random
sample of 5{,}050 individuals. FY26:

\begin{center}
\begin{tabular}{@{}lrrr@{}}
\toprule
Category & Share of traders & Loss-makers & Avg.\ loss per loss-maker \\
\midrule
Only options buyers      & \textbf{93\%}  & 90\% & \rs 1.3 lakh \\
Majorly options buyers   & 4.2\%          & 75\% & \rs 4.6 lakh \\
\textbf{Majorly options sellers} & \textbf{2.1\%} & \textbf{44\%} & \textbf{\rs 51.7 lakh} \\
Majorly futures traders  & 0.9\%          & 61\% & \rs 1.7 lakh \\
\bottomrule
\end{tabular}
\end{center}

\begin{keyfact}{The retail base is 97\% option \emph{buyers}, not sellers}
The sellers are the smallest group and have the \emph{lowest} loss rate --- and
lose eleven times more when they lose. The category still net-lost money in FY26.
SEBI's own framing: ``relatively high win rates may mask a tail-risk profile.''

\textbf{This is the single most important number in this document for your risk
management}, and it points the same way as everything in Section~\ref{sec:sizing}:
the danger in short premium is not the mean, it is the skew.
\end{keyfact}

\srcnote{Honest caveat, which SEBI states itself three times: the seller cell is
\textbf{105 traders in FY26}, of whom 46 were loss-making. The \rs 51.7 lakh
average is a 46-observation mean of a fat-tailed variable. \textbf{Treat the
direction as solid and the magnitude as indicative.}}

\subsubsection{The two studies that do what SEBI's cannot}

\textbf{Brazil, day traders in index futures.} Regulator records on 19{,}646
individuals who began day trading, followed for years. \textbf{97\% of those who
persisted more than 300 days lost money}; 1.1\% earned more than the minimum wage;
0.5\% more than a bank teller's starting salary, and those with high risk.
\textbf{No evidence of learning by doing.} \emph{What it does not show:} a single
instrument, a single style, retail with no stated prior training. It says nothing
about multi-day holding periods or people entering with institutional-grade
quantitative preparation. The ``no learning'' result is the part that generalises
least comfortably and is also the part most worth sitting with --- it is the direct
empirical rebuttal to ``I will be profitable once I have done my two years.''

\textbf{Taiwan, full population, 1992--2006 --- the only study here with a proper
persistence design.} \textbf{Less than 1\% of day traders predictably and reliably
earn positive abnormal returns net of fees.} Predictably profitable traders were
under 3\% and generated only about 10\% of volume. Sorting on one year's returns,
the top 500 went on to earn $+$37.9\,bp a day after fees the following year; the
bottom ranked $-$28.9\,bp.

\begin{keyfact}{The most useful number here for you}
That study says skill is \textbf{real, persistent, and detectable} --- and
concentrated in a sub-1\% tail. It also says the persistent winners are
\emph{identifiable in advance from track record}. So the honest test is empirical
and self-administered: \textbf{does your own live, cost-inclusive record show
persistence?} Nothing else answers it.
\end{keyfact}

\textbf{Proprietary-firm evaluations.} The best available dataset --- roughly
300{,}000 accounts across 100{,}000 traders at ten firms --- gives about
\textbf{14\% passing an evaluation}, about \textbf{45\% of funded traders ever
receiving a payout}, so roughly \textbf{7\% of everyone who buys a challenge ever
sees money back}, and 1--3\% still withdrawing beyond six months. Structural
caveats matter more than the percentages: most are simulated accounts, the
challenge fee is the revenue model (which aligns the firm against your passing),
and the model is forex and futures --- \textbf{it does not test the skill you have.}

\begin{keyfact}{The study you want does not exist}
There is \textbf{no published longitudinal study of the population you belong to}
--- a quantitatively trained person with a full-time job trading own capital
systematically. Brazil and Taiwan study retail day traders. Anyone who tells you
the base rate for your case is extrapolating.
\end{keyfact}

\subsection{Solo versus firm, per skill}

{\small
\setlength{\LTpre}{8pt}\setlength{\LTpost}{8pt}
\begin{longtable}{@{}p{3.1cm}p{3.0cm}p{9.5cm}@{}}
\toprule
\textbf{Skill} & \textbf{Solo?} & \textbf{The honest reason} \\
\midrule
\endfirsthead
\multicolumn{3}{@{}l}{\footnotesize\itshape Solo versus firm, continued}\\[2pt]
\toprule
\textbf{Skill} & \textbf{Solo?} & \textbf{The honest reason} \\
\midrule
\endhead
\bottomrule
\endlastfoot
Vol surface modelling & \textbf{No, not directly} &
It is an \emph{input}, not a strategy. At a firm it is table stakes --- every desk
has a calibrated surface and nobody is paid for having one. Its solo routes are
intra-surface relative value, which needs quoting infrastructure you cannot access,
or generating a rich/cheap signal, which collapses back into premium selling.
\textbf{You have already built the thing; the thing is not the edge.} \\[3pt]
Regime detection & \textbf{Yes} &
The only cleanly solo-scalable item --- no infrastructure, no latency, no venue
relationship. \textbf{Blunt counterweight:} it also has the weakest out-of-sample
evidence base in all of quant finance (Section~\ref{sec:regime}). \\[3pt]
Market making & \textbf{No} &
Requires exchange membership or a negotiated agreement, quoting obligations,
colocation, rebate-tier economics and inventory capital. \textbf{There is no retail
version.} Every ``retail market making'' claim is wide-spread scalping or marketing. \\[3pt]
Statistical arbitrage & \textbf{Partially, and worse than it looks} &
Indian cash shorting is intraday-only and stock lending is thin, so you cannot hold
the short leg. Crypto cross-venue is reachable but capital- and
infrastructure-hungry, margins have compressed, and the tax regime taxes the gross
legs of a spread independently. \\[3pt]
Options premium selling & \textbf{Yes --- fully} &
Which is exactly why it is crowded, and exactly what the SEBI numbers measure. The
one item an individual with margin can run at institutionally identical mechanics.
\textbf{The risk is not mean edge --- it is skew.} A single gap can exceed multiple
years of carry, and nothing about surface competence protects against it: the
exposure is to the jump, not to the surface. \\
\end{longtable}
}

\begin{keyfact}{The uncomfortable summary, in one line}
\textbf{The skill you have invested most in --- surface modelling --- is the least
solo-monetisable. The one that is fully solo-monetisable --- premium selling --- is
the most crowded and the most tail-exposed.}
\end{keyfact}

\subsection{Quant-adjacent income that is not a software job}

Assessed honestly; most of it does not pay on your stated horizon.

\textbf{Crowdsourced quant platforms.} Real money in aggregate --- tens of millions
distributed --- but the distribution is brutal: of roughly 14{,}400 models
submitting in a trailing year, only about 2{,}600 finished net positive, with a
\emph{median} positive earner around \$1 and a 90th percentile around \$235. And
you must stake your own capital in a volatile token. \textbf{Verdict: a hobby with
a lottery tail}, correctly classified alongside the trading base rates rather than
alongside income.

\textbf{Hosted quant competitions.} Prize pools typically \$100k--\$120k split
among the top handful of several thousand teams. The realistic return is a
credential and a recruiter conversation. \textbf{Given you are excluding a job, the
credential has zero cash value to you.}

\textbf{Systematic-strategy licensing --- the flagship attempt failed.} The most
serious retail alpha-licensing marketplace ever built was wound down: sales
disabled, live hosts stopped, licensees refunded. The reason is the one that
matters --- \textbf{there were not enough buyers. Institutions do not license
black-box alphas from anonymous individuals.} What remains are performance-fee and
subscription-share platforms, which are real and structurally sound; the binding
constraint is \emph{attracting third-party capital}, which is a distribution and
track-record problem, not a quant problem. \textbf{The gate is a two-to-three-year
live, audited, cost-inclusive track record before anyone subscribes --- which
consumes your entire stated horizon before the first rupee. It is not a two-year
path; it is a path that \emph{begins} at year two.}

\textbf{Paid research writing.} Real but small and brand-gated. The people earning
meaningfully almost universally arrived with an institutional brand first. Writing
does not bootstrap the brand quickly enough to matter on a two-year clock.
\srcnote{No verified income distribution data exists for this; reporting the absence.}

\textbf{Data products --- a genuine gap you are structurally blocked from filling.}
Indian minute-level index options data is expensive and poorly served, and there is
real demand for cleaned implied-vol surfaces. But redistributing exchange-derived
data requires an authorised data-vendor licence, and \textbf{the minute-wise data
you are about to receive comes through an employer relationship --- you cannot
resell what you get at work}, and attempting to would be the fastest possible way
to end both the job and the path. Closed for reasons that have nothing to do with skill.

\subsection{The arithmetic nobody states plainly}

\begin{keyfact}{Capital, not skill, is the binding constraint at your stage}
A \rs 25 lakh account returning a genuinely excellent, sustained \textbf{20\% net}
produces \rs 5 lakh gross and roughly \textbf{\rs 3.5 lakh after slab tax} --- a
fraction of a finance salary, from a strategy performing at a level the Taiwan
study places in the sub-1\% persistent tail.

To make the income question interesting rather than academic, the account has to be
an order of magnitude larger, and \textbf{the salary is the only documented
mechanism that gets it there inside two years.}

That is division, not advice.
\end{keyfact}

\srcnote{Employer personal-account-dealing policies vary substantially ---
pre-clearance requirements, minimum holding periods, restricted instrument lists,
and outright bans on instruments the firm covers. Check whatever contract you
eventually sign before any of the above is actionable.}

\newpage
\section{Current chatter --- UNVERIFIED}
\label{sec:chatter}

\begin{keyfact}{Read this section differently from the rest}
Nothing here passed the source filter. This is a \textbf{routing list} --- names,
repositories and, most valuably, \emph{complaints} --- to be sent to independent
verification. Nothing in it is evidence that anything works.
\end{keyfact}

\srcnote{Method disclosure. The dedicated social-sweep tool was invoked and
produced unusable output: incomplete first-run setup, no authenticated access to X,
rate-limiting on Reddit after six items, and a keyword collision on ``options''
that surfaced consumer-electronics threads. The fallback used the GitHub and Hacker
News APIs, a Reddit archive mirror, and web search. \textbf{X and YouTube were not
swept at all}, which is a real hole --- see the coverage failure below.}

\subsection{The complaints, which are the valuable part}

\begin{itemize}
\item \textbf{Gamma-exposure services are being publicly demolished.} A widely
upvoted thread and a companion study report that spot itself beats gamma-exposure
and max-pain as a predictor for index 0DTE. One participant claims a year of the
full greek-exposure dataset with a null result. The mechanism offered matches
Section~\ref{sec:gex} exactly: most services compute exposure from \emph{open
interest with no participant tags}, so the dealer sign is guessed.
\item \textbf{The variance risk premium in index 0DTE is real but ``razor thin''.}
One participant ran tick-level ATM implied vol against per-minute realised vol over
280 trading days and reports roughly \textbf{3--4 vol points of premium, strongly
regime-dependent, with extended runs where it is absent}. This is the most
methodologically serious retail work found in the window, and \textbf{the method
transfers directly to NIFTY minute data.} Route it to verification.
\item \textbf{Regime switching ``sounds good but disappoints''}, in a thread of
that exact title. The sharpest comment: one needs to detect a regime switch
\emph{before} it happens, and if you could do that you would be rich. This is the
community arriving independently at Section~\ref{sec:regime}'s conclusion.
\item \textbf{Latency is where backtests die, and most tools cannot see it.} The
recurring point: the gap between signal time, order time and exchange receipt is
longer than people think, and most backtesting tools have no visibility into it.
One participant models adversarial fills --- longs at the next bar's high ask,
shorts at the lowest bid, with a penalty above 1\% of that bar's volume.
\item \textbf{Look-ahead leakage is the most-confessed failure.} The cheapest
diagnostic offered: shift every feature forward one bar and rerun. \textbf{If the
edge survives, the leak is in the \emph{label}, not the features} --- usually a
rolling statistic that touches the same bars the forward return is built from.
\item \textbf{Overfitting is a trade-count question, not a time question.} Fifteen
or twenty forward trades tells you nothing; a broken system passes twenty trades
regularly. Two named curve-fit tells: high parameter sensitivity under 5--10\%
input noise, and a per-trade distribution where the mean is good but the
\emph{median} is not.
\item \textbf{Sizing, not strategy, is the reported killer.} ``Most people diagnose
a bad exit as a strategy problem when it was actually a position size that made
clear thinking impossible.''
\end{itemize}

\subsection{Repositories worth opening, and the flood to ignore}

Two are worth actually opening: a C++ options market-making engine that puts SVI
\emph{inside an execution loop} rather than in a plot, and a Rust vol-surface
library if you ever want the surface off the Python hot path. The mainstream
backtesting and crypto plumbing repositories are covered in
Section~\ref{sec:backtest}.

\begin{keyfact}{The pattern to recognise}
A GitHub search sorted by recency returns page after page of freshly pushed,
\textbf{zero-star} NIFTY and BankNifty repositories, nearly all within the last
48 hours, \emph{none} with a stated cost model. This is the most crowded and lowest
quality corner of the entire sweep. Separately, an active shill pattern is visible:
multiple new gamma-exposure dashboards launched in the same 30-day window, and one
participant's summary --- ``everyone without an edge has pivoted to using AI coding
assistants to build and sell a service.'' \textbf{Treat every new gamma dashboard in
this window as marketing until proven otherwise.}
\end{keyfact}

\subsection{Honest negatives about this sweep}

\begin{itemize}
\item \textbf{Hacker News is dead on this topic} --- nothing above single-digit
points in the window. Drop it from future sweeps.
\item \textbf{The careers-focused quant subreddit is pure noise} --- interview prep
and ``how do I break in''. Zero technical content. Drop it.
\item \textbf{Crypto quant chatter did not appear on Reddit at all.} Given your
40\% crypto weighting this is a \textbf{coverage failure, not a finding} --- crypto
vol discussion lives on X, Discord and Telegram, and X was not swept because no
credentials were available. \textbf{Re-run the crypto leg with authenticated access
before drawing any conclusion from its absence.}
\item \textbf{No cost model was stated anywhere in this chatter layer.} Under your
own filter almost nothing here would survive to a recommendation --- which is the
correct outcome for a chatter agent. This is a routing list, not a reading list.
\end{itemize}

\newpage
\section{What nobody has published: your openings}
\label{sec:openings}

Every item below is something round one or round two searched for and could not
find, on a question you are unusually well-positioned to answer. They are ordered
by expected value per week of effort.

\begin{enumerate}
\item \textbf{Where does an SVI fit residual exceed its own error bar on NIFTY
weeklies, and does it mean-revert net of the tick and STT?} Specified as a
pre-registered falsification test in Section~\ref{sec:residual}. \textbf{A rigorous
public study of this would be the first of its kind that could be found.} It is a
paper and a strategy at the same time, and the honest expectation is that it fails
--- which is itself the result.

\item \textbf{The deseasonalised realised-variance versus remaining-implied-variance
study on NIFTY minute data}, dollar-gamma weighted, with the jump decomposition
reported alongside. Specified step by step in Section~\ref{sec:rvproject}. Invisible
in end-of-day data, and it yields the active/inactive overnight variance ratio for
NIFTY --- a publishable number nobody has estimated.

\item \textbf{A public, runnable Indian SPAN implementation} driven by NSE's own
free daily files. No such thing exists; the parameter file is member-gated but
everything that matters for index derivatives is a published constant or a public
file. \textbf{About sixty lines of Python}, and a better use of a week than reading
another paper.

\item \textbf{The asymmetric hedging band for a one-sided cost.} STT is sell-side
only, which violates the symmetric-cost assumption underlying Leland,
Whalley--Wilmott, Barles--Soner and Zakamouline alike. The true no-trade region is
asymmetric and nobody has published this for the Indian market. \textbf{Days, not
weeks.}

\item \textbf{Hull--White minimum-variance delta coefficients for NIFTY}, estimated
on your own minute data and \emph{conditioned on position in the weekly Tuesday
cycle} --- a structure the US literature does not have. Nobody has published these,
and the original work explicitly excluded sub-14-day options, which is where all
your volume is.

\item \textbf{Business-time theta calibration for NIFTY weeklies.} With a Tuesday
expiry the weekend sits in the \emph{middle} of the option's life, and Indian
holidays cluster unevenly, so a weekly spans seven calendar days but three to five
trading days. Nobody has published on this, and minute data closes it.

\item \textbf{Crypto variance risk premium measured with a real cost model,
2024--2026.} The rigorous academic measurement stops earlier and includes no costs;
everything after it is vendor assertion. Free multi-year implied-vol index history
plus obtainable spot settles it. \textbf{A two-day project and genuinely original.}

\item \textbf{Dealer positioning on NIFTY, with the sign estimated rather than
assumed}, validated against the free daily participant-category file as an
aggregate constraint. No NIFTY study exists; the academic literature is entirely
S\&P 500.

\item \textbf{The weekend theta trade in crypto, costed.} The venue uses pure
calendar time by published convention, so the Friday implied-vol dip is real and
observable. Whether it survives a 25\%-of-premium round trip on short-dated options
is unpublished.

\item \textbf{Delta-neutral hedging bands with a funding carry term.} Perpetual
funding introduces a carry term in the hedge leg that no classical band literature
models, and it can dominate the transaction-cost term. Unfilled.

\item \textbf{A time series of the Deribit-versus-domestic-venue liquidity
comparison.} Section~\ref{sec:crypto} gives the method and a single snapshot.
Nobody has published the series, and you can record it.
\end{enumerate}

\begin{keyfact}{The common thread}
Nine of these eleven require exactly two things you already have or are about to
have: \textbf{minute-resolution data and a working surface fitter}. None requires
capital. Most take days to weeks. And several are simultaneously a piece of
original research and a falsification test of something you would otherwise have
believed --- which is the cheapest possible way to learn a market.
\end{keyfact}

\newpage
\section{Appendices}

\subsection{Paid, worth considering later}
\label{app:paid}

Listed separately and never mixed into the recommendations. Check institutional
access before buying any of it --- see the access note in Section~\ref{sec:people}.

\textbf{Books, in priority order.}
\begin{enumerate}
\item \textbf{Sinclair, \emph{Volatility Trading} (2nd ed.)} --- the one book with
no free substitute. Position sizing under non-normal option returns, the real
economics of the variance premium, and an unusually long treatment of why most
published option edges are cost artefacts.
\item \textbf{Bergomi, \emph{Stochastic Volatility Modeling}} --- but read the two
free chapters and the free \emph{Smile Dynamics} papers first; they carry perhaps
40\% of the practical value at zero cost.
\item \textbf{Taleb, \emph{Dynamic Hedging}} --- for the taxonomy of risk and the
only book-length treatment of vega bucketing by someone who did it on a desk.
\item \textbf{Bouchaud et al., \emph{Trades, Quotes and Prices}} --- but the free
100-page precursor by the same authors covers a large fraction of it.
\item \textbf{Sinclair, \emph{Positional Option Trading}} --- second, not first;
substantial overlap with the above.
\item \textbf{A\"it-Sahalia \& Jacod, \emph{High-Frequency Financial Econometrics}}
--- the reference behind the realised-volatility section. Library, not purchase.
\end{enumerate}

\textbf{Two paywalled papers worth pulling through your institution:}
the Dynamic Nelson--Siegel smile-harvesting paper on NSE data (\textbf{check
whether it nets transaction costs --- this is the single most important open
question in the Indian material}); and the paper establishing the Indian
single-stock implied-over-realised premium that governs the dispersion sign.

\textbf{Data, not education, and only after your pipeline is proven on free
samples:} continuous tick-level crypto options history; institutional crypto
derivatives feeds --- likely reachable through your institution before paying retail.

\textbf{Explicitly not recommended at any price:} the proprietary-firm options
programmes, the mentoring businesses, and the entire Indian option-selling course
market.

\subsection{Primary documents --- the index}
\label{app:primary}

The documents to read yourself rather than take on trust. All free.

\textbf{Indian regulation and market structure}
\begin{itemize}\footnotesize
\item SEBI, \emph{Measures to Strengthen Equity Index Derivatives Framework}, 1 Oct 2024 --- the six measures and their effective dates.
\item SEBI, \emph{Final Settlement Day (Expiry Day) for Equity Derivatives Contracts}, 26 May 2025 --- and the NSE implementation circulars, which are where the backtest trap lives.
\item NSE circular on lot-size revision, 3 Oct 2025.
\item SEBI, \emph{Review of eligibility criteria for entry/exit of stocks in derivatives segment}, 30 Aug 2024.
\item SEBI Master Circular for Stock Exchanges and Clearing Corporations, \textbf{Chapter 5} --- the complete SPAN specification, the 16-scenario table, ELM rates, and the peak-margin framework.
\item NSE Clearing risk-management FAQ --- SPAN file cadence, the ELM ladder, and the beginning-of-day parameter rule that makes margin simulation feasible.
\item NSE Clearing settlement-price and STT pages --- the method, and the per-client STT report file spec you can parse.
\item SEBI, \emph{Profitability} and \emph{Trading Behaviour of Individual Traders in the Equity Derivatives Segment (FY25--FY26)}, Aug 2026 --- the current numbers, superseding everything earlier.
\end{itemize}

\textbf{Indian daily data files, all anonymously downloadable}
\begin{itemize}\footnotesize
\item Market lot file --- current lot size per underlying per expiry month.
\item Daily volatility file --- EWMA daily and annualised volatility per symbol, \textbf{the exact input that feeds the price and volatility scan ranges}.
\item F\&O bhavcopy --- OHLC, settlement price, open interest, lot size, underlying price, per contract.
\item Settlement price file --- the number determining exercise STT intrinsic value.
\item Collateral haircut files.
\item \textbf{Participant-wise open interest, signed by category, daily} --- the file in Section~\ref{sec:gex} with no US equivalent.
\item The circulars JSON index --- every circular with a direct PDF link.
\end{itemize}
\srcnote{These require a cookie bootstrap and a browser user-agent; nothing else.}

\textbf{Tax}
\begin{itemize}\footnotesize
\item Income-tax Act 2025, official gazette text --- ss.\ 66(31), 66(33), 108--113, 121 for the F\&O treatment; s.194(1) Table Sl.\ 4 and s.393(1) Table Sl.\ 8(vi) for virtual digital assets.
\item Finance Bill/Act 2026 --- the operative STT amendment clause.
\item NSE STT rate table and tax circular --- the definitive levy table.
\item ICAI Guidance Note on tax audit --- the only authoritative statement of how derivative turnover is computed.
\item CBDT ITR-2 validation rules --- the schedule that enforces transaction-level VDA reporting.
\end{itemize}

\textbf{Foreign exchange}
\begin{itemize}\footnotesize
\item RBI Master Direction on the Liberalised Remittance Scheme.
\item \textbf{FEM (Overseas Investment) Directions 2022, definition (ix)(a)(iii)} --- the single clause that decides the US-venue question, and the one every secondary source paraphrases wrongly.
\end{itemize}

\textbf{Crypto venue specifications}
\begin{itemize}\footnotesize
\item Exchange API reference and contract introduction policy --- the listing ladder and strike intervals.
\item Index-price and settlement pages --- the capped-median index and the 30-minute settlement window, including the fact that expiring-option delta decays linearly to zero over the final thirty minutes, which will silently break a naive hedger.
\item Portfolio margin model --- a real specification with formulas. \textbf{Note the published deck is superseded}; use it for structure and pull live parameters from your own account.
\item Inverse-options specification --- the exchange's exact Black-Scholes convention, including that implied vol uses the \emph{forward}, and that time is pure calendar time.
\item Fee schedule --- the 12.5\%-of-premium cap.
\end{itemize}

\textbf{Method and measurement}
\begin{itemize}\footnotesize
\item BCBS d457, \emph{Minimum capital requirements for market risk} --- MAR32 for the P\&L attribution thresholds, MAR21 for the published vega-bucketing scheme.
\item The 0DTE replication package and its \texttt{KNOWN-ISSUES} file --- the best-documented cost-scale failure in the public options literature.
\end{itemize}

\subsection{What could not be verified}
\label{app:unverified}

Collected in one place, because a survey that hides its gaps is worse than one that
shows them.

\begin{itemize}
\item \textbf{Whether the DNS smile paper nets transaction costs} --- paywalled, no
preprint found. The single most important open question in the Indian material.
\item \textbf{SEBI turnover fee and stamp-duty rates} on derivatives; \textbf{GST at
18\%} on brokerage. Together about 0.3\,bp; they move no conclusion.
\item \textbf{Whether NSE transaction charges are still at the September 2024 level.}
\item \textbf{Whether the November 2024 expiry-day ELM add-on has since been
modified} by a later risk circular.
\item \textbf{The consolidated ``as amended'' Income-tax Act PDF} --- 403-blocked;
the gap was closed by reading the amendment clauses directly.
\item \textbf{Whether the 1\% withholding attaches to VDA \emph{derivatives}} ---
genuinely unresolved in the statutory text. A gap in the law, not a research failure.
\item \textbf{The client-level 50\% cash-margin requirement} --- verified at
clearing-member level only.
\item \textbf{A public register of registered crypto service providers} --- none
exists, so a venue's stated registration is a self-declaration.
\item \textbf{Current numeric portfolio-margin parameters} on the crypto venue ---
exposed only to logged-in accounts.
\item \textbf{The Sinclair sizing paper's body} --- the repository blocks
programmatic access, so the abstract was read but not the paper, and no formula from
it is reproduced. It is free and short; download it.
\item \textbf{Several classical hedging papers} --- Whalley \& Wilmott (1997),
Hodges \& Neuberger (1989), Davis--Panas--Zariphopoulou (1993) --- were not reached
in the original. A free working paper restates all three faithfully and was used as
a labelled substitute.
\item \textbf{Gaussian HMMs failing to reproduce volatility clustering} --- stated
in the jump-model literature, not independently confirmed. Cheap to check by simulation.
\item \textbf{Chapter-level content for several books} is from prior knowledge
rather than from opening them this round. Chapter \emph{titles} should be treated as
approximately, not exactly, correct.
\item \textbf{X and YouTube were not swept} in the social layer. A named,
acknowledged hole.
\end{itemize}

\vspace{1.5em}
\begin{keyfact}{One closing note on method}
The most valuable single output of this exercise was not a source. It was the
discovery that three claims round one presented confidently --- that the US was
reachable, that the standard dealer-positioning construction applies to NIFTY, and
that the SEBI studies show retail is short volatility --- were all \textbf{false},
and that a fourth, the SVI residual, was \textbf{unresolvable on the data round one
used}.

Every one of those was caught by making a second pass whose only job was to attack
the first. If you take one process habit from this document, take that one.
\end{keyfact}

\end{document}
